\documentclass[11pt]{ctexart}
\usepackage[a4paper,margin=1in]{geometry}
\usepackage{graphicx}
\usepackage{xcolor}
\usepackage{hyperref}
\definecolor{refgray}{gray}{0.45}
\title{\textbf{市场最新观点汇总}\\[0.4em]{\large 宏观数据噪音加剧，结构性分化成为主线：从AI算力到消费品定价，从利率偏斜到铜供给断裂，市场定价正从方向性交易转向alpha挖掘。}}
\author{KC桌面 自动生成}
\date{260603}
\begin{document}
\maketitle
\tableofcontents
\newpage
\section{一页摘要}
\begin{itemize}
  \item 瑞士手表出口4月骤降16.6\%，但剔除关税抢运基数后，美国市场仍录得正增长，大中华区企稳，中东地缘风险成为新变量。
  \item 利率期权偏斜信号在2019年后对长端利率的预测能力反而增强，当前信号指向短久期，市场对尾部风险的定价可能不足。
  \item 5月中国PMI回落至50.0，生产与订单缺口扩大，上游成本向下游传导不畅，企业利润空间受挤压，出口订单跌破荣枯线。
  \item 美国消费韧性建立在低储蓄率和实际收入下滑之上，AI资本开支是少数不受地缘风险干扰的加速领域，但劳动力市场影响有限。
  \item 大宗商品市场持仓总估值下降690亿美元，但铜逆势吸金，黄金连续三周流出，能源下跌反映中国需求结构性转变而非周期性衰退。
  \item 日本半导体设备4月出货额同比增长25\%，存储资本开支和先进封装驱动结构性增长，而非周期尾声。
  \item DRAM价格一年涨超六倍，AI需求正将存储器从周期品变为战略资产，PC和手机面临系统性供给缺口，芯片通胀正在向消费品扩散。
  \item 全球汽车4月销量微增0.7\%，但分化剧烈：欧洲和新兴市场增长，美中下滑；现代起亚HEV份额逼近20\%，改写美国市场竞争规则。
  \item 中国新能源渗透率突破63\%，但订单增长由新车型集中投放驱动，品牌间分化断层式加剧，持续性待观察。
  \item 铜价突破的关键在于供给侧三个结构性断裂：矿山产量零增长、中断率上调、废铜回收跟不上，预计一个月内触及14500美元/吨。
  \item 中国房地产市场新房成交反弹是基数效应，二手房正在成为定价锚点，去化率下降显示供给压力仍在。
  \item 新兴市场债券曲线整体持稳，但个体分化剧烈，印尼曲线大幅走平，埃及曲线最陡，定价逻辑从利率跟随转向信用分化。
\end{itemize}
\section{宏观与利率：数据噪音中的结构性信号}
\textbf{全球宏观数据被关税、基数效应和季节性因素严重扭曲，市场对真实需求的判断存在系统性偏差。利率市场方面，期权偏斜信号在长端利率上的预测能力在2019年后反而增强，当前信号指向短久期。}

\begin{itemize}
  \item 4月瑞士手表出口同比下滑16.6\%，但剔除去年关税抢运基数后，美国市场仍录得+8.9\%的增长，大中华区对香港和中国内地出口分别增长13.5\%和17.1\%，中东全线下跌显示地缘风险影响。
  \item 5月中国官方制造业PMI回落至50.0，生产指数51.2仍处扩张，但新订单指数49.9五个月来首次跌破荣枯线，新出口订单骤降1.7个百分点至48.6，生产跑在需求前面的缺口在扩大。
  \item PMI价格分项显示原材料购进价格指数60.5，出厂价格指数51.9，两者差距连续五个月未缩小，企业利润被压缩，上游库存下降而产成品库存上升。
  \item 利率期权偏斜信号在2009年至今的回测中保持预测能力，2019年后在30年期利率上表现更强，在2年期上因宏观冲击主导而变弱。
  \item 当前偏斜信号指向曲线整体偏短久期，且偏斜进一步加深，市场对尾部风险的定价可能不足。
  \item 美国Q1财报显示标普500实际营收同比增长6.3\%，但个人储蓄率偏低，实际可支配收入同比下降1.1\%，研报预计2026下半年实际消费增速仅1.3\%，低于市场共识。
\end{itemize}
\begin{figure}[htbp]
\centering
\includegraphics[width=0.88\linewidth]{figures/fig_001.jpg}
\caption{EXHIBIT 1 - EXHIBIT 1: Exports of Swiss Wristwatches watches by price categories Swiss watch exports by price - Total market (1'000 units)}
\end{figure}
\begin{figure}[htbp]
\centering
\includegraphics[width=0.88\linewidth]{figures/fig_003.jpg}
\caption{EXHIBIT 3 - EXHIBIT 3: Swiss watch exports by country in absolute numbers Swiss watch exports by country - Key markets only}
\end{figure}
\begin{figure}[htbp]
\centering
\includegraphics[width=0.88\linewidth]{figures/fig_025.jpg}
\caption{Figure 5 - Figure 5: , China's manufacturing PMI returned to exactly the 50 boom-bust line}
\end{figure}
\begin{figure}[htbp]
\centering
\includegraphics[width=0.88\linewidth]{figures/fig_030.jpg}
\caption{Figure 10 - Figure 10: Input prices continue to rise by much faster than output prices}
\end{figure}
\begin{figure}[htbp]
\centering
\includegraphics[width=0.88\linewidth]{figures/fig_004.jpg}
\caption{Exhibit 1 - Exhibit 1: 10y tail results for all combinations of lookback window, option expiry, and trading delay \# Compare LIBOR and SOFR era The LIBOR-SOFR transition provides a convenient division of the sample into two distinct macro env}
\end{figure}
\begin{figure}[htbp]
\centering
\includegraphics[width=0.88\linewidth]{figures/fig_011.jpg}
\caption{Exhibit 8 - The performance deterioration since 2019 has been most pronounced in the 2y tail, while the 30y tail is the only part of the curve that saw an improvement. One interpretation is that front-end rates have become increasingly driven by macro and policy shocks, r}
\end{figure}
\begin{figure}[htbp]
\centering
\includegraphics[width=0.88\linewidth]{figures/fig_032.jpg}
\caption{Exhibit 1 - Exhibit 1: Real Corporate Revenues Grew a Strong \$6.3\textbackslash{}\%\$ Year-over-Year in 2026Q1}
\end{figure}
\begin{figure}[htbp]
\centering
\includegraphics[width=0.88\linewidth]{figures/fig_033.jpg}
\caption{Exhibit 2 - Exhibit 2: Alternative Measures of Nominal Consumer Spending Growth Remain Solid}
\end{figure}
{\small\color{refgray}\textbf{References:} [R001] 瑞士手表出口暴跌16.6\%：关税扭曲下的数据噪音，掩盖了真正的复苏信号; [R002] 市场低估的不是波动，而是期权偏斜信号在利率市场中的结构性价值; [R003] 五月PMI回落至荣枯线，市场真正担心的不是增长放缓，而是价格传导断裂; [R004] 市场真正低估的不是需求，而是供给侧的再定价}

\section{AI/算力与半导体：从周期品到战略资产的再定价}
\textbf{AI需求正将存储器从周期品变为战略资产，DRAM成为AI建设中的首要瓶颈，供给缺口在2028年前看不到快速解决方案。日本半导体设备受益于存储资本开支和先进封装的结构性增长。}

\begin{itemize}
  \item DRAM价格过去一年累计上涨超过六倍，HBM到2028年将占据领先级存储器晶圆产能约34\%，而2023年仅为6\%，AI需求对传统市场的挤出效应显著。
  \item 到2027年，PC市场将面临约15\%的存储器短缺，智能手机市场短缺约12\%，分别对应约5800万台PC和1.34亿部手机被AI挤占产能。
  \item 日本半导体设备4月出货额同比增长25\%（日元计），组装设备同比暴增73\%，测试设备增长26\%，3个月移动平均环比仍增长6\%。
  \item 全球WFE市场2026年预计同比增长21.4\%，2027年再增长18.2\%，DRAM和NAND资本开支是主要驱动力。
  \item 存储器股票当前不到10倍的远期市盈率，研报认为市场对盈利持续性的定价仍不够充分，美光2027年EPS预测上调48\%，SanDisk上调24\%。
  \item Neocloud正在从提供GPU算力向提供生产级推理服务转型，IREN收购Mirantis、Nebius收购Eigen AI，标志AI基础设施进入2.0阶段。
\end{itemize}
\begin{figure}[htbp]
\centering
\includegraphics[width=0.88\linewidth]{figures/fig_213.jpg}
\caption{Exhibit 1 - EXHIBIT 2: April single-month billings data was -38\% MoM and +25\% YoY in JPY.}
\end{figure}
\begin{figure}[htbp]
\centering
\includegraphics[width=0.88\linewidth]{figures/fig_216.jpg}
\caption{EXHIBIT 4 - EXHIBIT 4: Monthly billings for Japanese assembly SPE was +73\% YoY / -39\% MoM.}
\end{figure}
\begin{figure}[htbp]
\centering
\includegraphics[width=0.88\linewidth]{figures/fig_224.jpg}
\caption{Exhibit 1 - Exhibit 1: DRAM prices YoY}
\end{figure}
\begin{figure}[htbp]
\centering
\includegraphics[width=0.88\linewidth]{figures/fig_227.jpg}
\caption{Exhibit 4 - Exhibit 4: We expect hyperscaler capex to surpass US\textbackslash{}\$1tr in 2027 and an aggregate \textbackslash{}\textasciitilde{}US\textbackslash{}\$2tr invested since 2024}
\end{figure}
\begin{figure}[htbp]
\centering
\includegraphics[width=0.88\linewidth]{figures/fig_229.jpg}
\caption{Exhibit 7 - Exhibit 7: AI prioritization turns 2027 supply growth into a consumer memory shortfall PC: AI priority creates residual DRAM shortfall}
\end{figure}
\begin{figure}[htbp]
\centering
\includegraphics[width=0.88\linewidth]{figures/fig_365.jpg}
\caption{Exhibit 1 - Exhibit 1: Micron Gross PPE by asset class, PPE to expand more by the end of FY27 than 1QFY20-2QFY26 We think both stocks can continue to outperform as the multiples remains reasonable on near term estimates (sub 10x PE on CY27 e}
\end{figure}
{\small\color{refgray}\textbf{References:} [R008] 市场低估了日本半导体设备的结构性机会，而非周期的尾声; [R011] 市场正在低估“芯片通胀”的扩散半径：从AI算力成本到全球消费品定价; [R020] 市场低估的不是需求，而是AI基础设施供应链的再定价; [R023] 市场真正低估的不是AI需求，而是存储器供给侧的再定价; [R028] 市场真正低估的不是算力需求，而是AI基础设施的供给结构正在被重新定义}

\section{能源与大宗商品：定价权从宏观叙事向微观供给切换}
\textbf{大宗商品市场正在从宏观叙事驱动转向微观供给结构重定价。能源下跌反映中国需求结构性转变，铜的供给断裂比需求爆发更重要，锂价上行由基本面驱动而非故事驱动。}

\begin{itemize}
  \item 过去一周全球大宗商品持仓总估值下降690亿美元至1.82万亿美元，能源市场下跌560亿美元，布伦特和WTI分别跌11\%和10\%。
  \item 中国石油需求下降9\%，但研报判断其中相当部分是永久性的消费模式转变，而非周期性需求萎缩，航空燃油和石脑油需求将大致恢复。
  \item 铜市场净流入23亿美元，逆势吸金。研报将2026年全球铜矿产量增速下调至-0.2\%，2027-2028年矿山中断率假设上调至7\%，预计一个月内铜价触及14500美元/吨。
  \item 黄金连续三周资金流出，但COMEX Managed Money净多头反而增加3900手，全球黄金ETF持仓YTD仍为正增长。
  \item 锂价年内有望测试25万元/吨，2026年前四个月中国电池产量同比增长41\%至874 GWh，储能电池产量增长99\%，需求侧高频信号强劲。
  \item 碳酸锂总库存从年初9.5万吨降至6.5万吨，冶炼厂库存去化最明显，江西某大矿因采矿证到期停产，复产时间不确定。
\end{itemize}
\begin{figure}[htbp]
\centering
\includegraphics[width=0.88\linewidth]{figures/fig_044.jpg}
\caption{Figure 1 - Figure 1: Global commodity markets saw net outflows largely driven by precious metals and energy markets USD billion, cumulative flows across tracked commodity markets}
\end{figure}
\begin{figure}[htbp]
\centering
\includegraphics[width=0.88\linewidth]{figures/fig_045.jpg}
\caption{Figure 2 - Figure 2: The estimated value of global commodity market open interest decreased by 4\% WoW USD billion, estimated open interest across tracked commodity markets}
\end{figure}
\begin{figure}[htbp]
\centering
\includegraphics[width=0.88\linewidth]{figures/fig_046.jpg}
\caption{Figure 3 - Figure 3: Net length across commodity markets decreased by \textbackslash{}\$20 billion, largely driven by crude oil and agri market real USD million, Net investor positioning aggregated across tracked commodity markets}
\end{figure}
\begin{figure}[htbp]
\centering
\includegraphics[width=0.88\linewidth]{figures/fig_051.jpg}
\caption{Figure 9 - Figure 9: The estimated value of open interest across base metals markets increased by 1\% WoW USD billion, estimated open interest across tracked base metals markets}
\end{figure}
\begin{figure}[htbp]
\centering
\includegraphics[width=0.88\linewidth]{figures/fig_279.jpg}
\caption{Figure 3 - © 2026 Citi Inc. No redistribution without Citi's written permission. Net long copper positioning is now relatively elevated but remains below the highs of late 2025 and 2024. Funds gross short positions remain subdued on COMEX, likely partly due to lingering }
\end{figure}
\begin{figure}[htbp]
\centering
\includegraphics[width=0.88\linewidth]{figures/fig_283.jpg}
\caption{Figure 8 - © 2026 Citi Inc. No redistribution without Citi's written permission. Figure 8. Scrap and mine supply growth expected to fail to match AI and energy transition demand growth this year and next (if we assume any cyclical demand recovery the gap becomes more sev}
\end{figure}
\begin{figure}[htbp]
\centering
\includegraphics[width=0.88\linewidth]{figures/fig_314.jpg}
\caption{Figure 1 - \# Battery production pipeline reaffirms the strong demand month by month Channel check with ZE Consulting suggests on-the-ground battery production pipeline extends the strong momentum YTD, easing off previous concerns that (1) NEV might be weaker than expecte}
\end{figure}
\begin{figure}[htbp]
\centering
\includegraphics[width=0.88\linewidth]{figures/fig_319.jpg}
\caption{Figure 7 - © 2026 Citi Inc. No redistribution without Citi's written permission. Figure 7. Lithium carbonate weekly inventory}
\end{figure}
{\small\color{refgray}\textbf{References:} [R006] 大宗商品市场正在经历的不是波动，而是定价权的结构性转移; [R013] 铜价突破的关键不在于需求爆发，而在于供给侧的三个结构性断裂; [R019] 市场真正低估的不是需求，而是供给侧的再定价能力}

\section{地产与消费：二手房定价权转移与消费渠道重置}
\textbf{中国房地产市场新房成交反弹是基数效应，二手房正在成为定价锚点。消费支付数据回暖但官方零售数据滞后，消费从线上回流线下正在重置渠道结构。}

\begin{itemize}
  \item 50城新房周度成交同比增长10\%，但年初至今累计同比仍为-14\%，整体去化率从66\%降至53\%，一线城市从66\%降至63\%，二线城市降至42\%。
  \item 10城二手房周度成交同比增长9\%，年初至今累计同比保持5\%正增长，中原6城二手房挂牌价指数同比-17.5\%，较前一周略有收窄。
  \item 全国二手房价格指数月环比跌幅从年初-0.5\%收窄至-0.4\%，但背后是三个价格体系在同时运行：政府回购支撑的微型单元、高租金收益率资产、仍在下跌的中高端改善型住房。
  \item 上海50平方米以下微型单元自2025年11月触底以来已涨11\%，深圳中高端住宅开始企稳，但北京、广州大部分房源仍在承压。
  \item 一季度总系统支付量同比大增29\%，银行卡消费支付同比转正至3.2\%，而官方社零增速放缓至2.4\%，消费从线上回流线下是主要解释。
  \item 家庭金融资产一季度同比增长10.8\%，主要由股票和共同基金驱动，居民部门资产负债表修复速度可能快于市场预期。
\end{itemize}
\begin{figure}[htbp]
\centering
\includegraphics[width=0.88\linewidth]{figures/fig_313.jpg}
\caption{Exhibit 1 - Exhibit 1: Primary weekly unit sales and four-week moving average – Total 50 cities}
\end{figure}
\begin{figure}[htbp]
\centering
\includegraphics[width=0.88\linewidth]{figures/fig_366.jpg}
\caption{Exhibit 1 - Exhibit 1: Total system payment volume picked up 29\% yoy in 1Q26 vs 5\% in 4Q25. Bank card consumption growth turned positive, up 3.2\% yoy, vs decline in 2025.}
\end{figure}
\begin{figure}[htbp]
\centering
\includegraphics[width=0.88\linewidth]{figures/fig_367.jpg}
\caption{Exhibit 2 - Exhibit 2: UnionPay payment growth rebounded from 4Q, up 11.2\% yoy; while NetsUnion payment growth slightly moderated to 7.2\% yoy.}
\end{figure}
{\small\color{refgray}\textbf{References:} [R017] 市场真正低估的不是需求，而是二手房定价权的结构性转移; [R024] 消费支付数据回暖，但市场真正低估的是银行非息收入的修复弹性; [R029] 市场真正低估的不是需求，而是供给侧的再定价}

\section{中国新能源车：渗透率高位下的订单分化与出口爆发}
\textbf{新能源渗透率突破63\%后，增长动力从电动化替代切换为产品周期竞争。品牌间订单分化断层式加剧，比亚迪出口爆发正在改变其估值逻辑。}

\begin{itemize}
  \item 5月新能源车批发量环比增长12\%，同比增长7\%，超市场预期。新能源渗透率在5月1-24日达到62.5\%-63.6\%。
  \item 蔚来周订单环比暴增228\%，鸿蒙智行增长92\%，比亚迪增长15\%，而理想汽车环比下滑25\%，小鹏下滑61\%，分化由新车型集中投放驱动。
  \item 比亚迪5月新能源车总销量38.35万辆，海外销量16.06万辆，环比增长19\%，前五个月累计海外销量61.69万辆，同比增长65\%，海外销量占比已从约20\%提升至约42\%。
  \item 新能源车终端折扣收窄至7.75\%，燃油车折扣同步收窄至19.47\%，上游碳酸锂价格微降至17.95万元/吨，电池价格稳定。
  \item 现代起亚在美国5月合计市占率达11.8\%，HEV市占率逼近20\%达19.9\%，同比跃升8.3个百分点，HEV销量增速是行业3-5倍。
  \item HEV在现代总销量中占比从14.7\%跃升至27.1\%，在起亚中从11.5\%飙升至31.5\%，HEV正在成为利润压舱石而非过渡产品。
\end{itemize}
\begin{figure}[htbp]
\centering
\includegraphics[width=0.88\linewidth]{figures/fig_299.jpg}
\caption{Exhibit 2 - Exhibit 2: NEV dealer discount trend}
\end{figure}
\begin{figure}[htbp]
\centering
\includegraphics[width=0.88\linewidth]{figures/fig_305.jpg}
\caption{Exhibit 1 - Exhibit 1: HMC \& Kia (together HMK) US monthly key metrics snapshot Exhibit 2: HMK US volume trend}
\end{figure}
\begin{figure}[htbp]
\centering
\includegraphics[width=0.88\linewidth]{figures/fig_307.jpg}
\caption{Exhibit 4 - Exhibit 4: HMK US M/S trend}
\end{figure}
\begin{figure}[htbp]
\centering
\includegraphics[width=0.88\linewidth]{figures/fig_309.jpg}
\caption{Exhibit 6 - Exhibit 6: HMK US HEV volume trend}
\end{figure}
\begin{figure}[htbp]
\centering
\includegraphics[width=0.88\linewidth]{figures/fig_368.jpg}
\caption{Figure 1 - Huawei Harmony – May-26 deliveries came at 46,122 units (+4\% YoY/+41\% MoM), with 5M26 deliveries +27\% YoY to 191,598 units. Orders update: all-new M9 non-refundable order over 20k units in 24 hours post launch, Luxeed V9 non-refundable orders over 10.5k units }
\end{figure}
\begin{figure}[htbp]
\centering
\includegraphics[width=0.88\linewidth]{figures/fig_369.jpg}
\caption{Figure 2 - Horizontal Axis: May-26 MoM growth © 2026 Citi Inc. No redistribution without Citi's written permission. Figure 2. Various NEV Players: May-26 Sales vs May-26 YoY}
\end{figure}
{\small\color{refgray}\textbf{References:} [R014] 新能源渗透率突破63\%的真正含义：市场正在重新定义赢家的门槛; [R015] 新能源渗透率突破63\%后，市场真正低估的是新车型的订单转化效率; [R016] 现代起亚正在悄悄改写美国汽车市场的竞争规则; [R027] 五月新能源车销量超预期，但市场真正低估的不是需求，而是供给侧的再定价}

\section{航空航天与防务：供给端的结构性重塑}
\textbf{航空航天与防务行业正在经历供给端的深层重构，从供应链结构到生产模式再到竞争格局，市场定价尚未充分反映这些不可逆的变化。}

\begin{itemize}
  \item 所有商业航空公司CEO均表示尚未看到需求疲软迹象，售后市场订单簿已覆盖未来12-18个月，但OE与后市场同时高负荷运转考验供应链极限。
  \item 波音确认737月产47架目标夏季可达，未来可能到52、57甚至63架，787年底目标月产10架，供应商库存正在下降。
  \item 防务领域导弹需求爆发式增长，计划5-7年内产能翻3-4倍，合同带补偿条款，固体火箭发动机供应商均在扩产。
  \item Anduril等新进入者从根本上挑战传统防务承包商的成本结构与交付速度，行业龙头公司的竞争优势和估值逻辑正在被重新定义。
  \item 所有防务公司预期预算会涨，但无人敢打包票能到1.5万亿美元，2027年预算仍存在不确定性。
\end{itemize}
{\small\color{refgray}\textbf{References:} [R005] 航空航天与防务：市场真正低估的不是需求，而是供给端的结构性重塑}

\section{日本工程机械与IT服务：关税窗口与AI效率竞争}
\textbf{美国对日本工程机械关税下调的窗口期设计暴露了其工业政策节奏，日本IT服务业正从系统开发转向企业AI代理设计，效率竞争成为分化关键。}

\begin{itemize}
  \item 美国将联合收割机、推土机、叉车等对美出口关税从25\%降至15\%，适用于6月8日及之后抵达货物，日落日期为2027年12月31日，为期约18个月窗口期。
  \item 研报测算关税下调对小松当期利润影响约100亿日元，约占年营业利润2\%-3\%，久保田和竹内制作所影响有限。
  \item 日本IT服务三家公司订单环境均强劲，但利润指引显著分化：Simplex预计FY3/27营业利润增长19\%，连续第七年两位数增长，而NSD和DTS仅2\%-3\%。
  \item Simplex开始专注AI驱动开发，按合同项目收费，AI提效带来的工时节省大部分能转化为公司利润，预计FY3/29开始贡献。
  \item 企业AI应用的ROI困境正在催生新的中间层市场，IT服务商是天然承接者，业务核心从系统开发转向企业AI代理的设计、管理与安全。
\end{itemize}
{\small\color{refgray}\textbf{References:} [R010] 关税下调的真正价值不在当期利润，而在日本工程机械的竞争格局重塑; [R032] 市场低估的不是AI算力需求，而是IT服务商从“系统开发”到“代理设计”的职能跃迁; [R033] 日本IT服务业的真正分化点：谁能在AI驱动的效率红利中，把成本节约转化为利润增长}

\section{医药与生物科技：ADC与双抗的交叉验证}
\textbf{肺癌一线治疗标准正在从PD-1单药或联合化疗转向TROP2 ADC联合免疫与PD-1xVEGF双抗联合化疗两条技术路线的并行竞争，全球数据验证是真正分水岭。}

\begin{itemize}
  \item sac-TMT在OptiTROP-Lung05研究中PFS的HR接近0.5，停药率仅5\%左右，毒性可控，但全球OS数据才是决定能否改变标准治疗的关键。
  \item 依沃西单抗在HARMONi-6中OS获益明显HR=0.66，但KOL提醒中国数据需在全球人群中验证，尤其是>65岁患者亚组存在争议。
  \item KRAS G12C单药在一线非小细胞肺癌中的故事线比组合疗法更清晰，组合疗法需要更精细的患者分层和更主动的毒性管理。
  \item 日本药企ONO-4578在二线胃癌Phase 2中PFS显著改善（HR 0.67），但疗效优势几乎完全集中在PD-L1阳性患者中，阴性组甚至出现恶化趋势。
  \item 武田与信达合作的IBI363在PD-L1低表达/阴性NSCLC中确认ORR达81.8\%，无论鳞癌/非鳞癌均观察到疗效。
\end{itemize}
\begin{figure}[htbp]
\centering
\includegraphics[width=0.88\linewidth]{figures/fig_373.jpg}
\caption{Exhibit 6 - Exhibit 1: Overview of select 1L metastatic NSCLC trials Histology}
\end{figure}
\begin{figure}[htbp]
\centering
\includegraphics[width=0.88\linewidth]{figures/fig_374.jpg}
\caption{Exhibit 3 - Exhibit 3: OS curves from the ITT population of HARMONi-6}
\end{figure}
{\small\color{refgray}\textbf{References:} [R009] 肺癌治疗格局的真正变量，不在PD-1，而在ADC与双抗的交叉点; [R021] ASCO 2026的真正信号：日本药企的差异化正在被重新定价; [R030] KRAS G12C的真相：不是所有抑制剂都一样，组合疗法的门槛比想象中更高; [R031] 肺癌一线治疗变局：真正值得关注的不是单一数据，而是两类新机制的“交叉验证”}

\section{中国平价零售与自动化：供给侧主动再定价}
\textbf{中国平价零售赛道正在经历供给侧的主动再定价，补贴扩大的本质是对加盟商进行供给侧筛选而非价格战。自动化需求中机床工具和通用机械走强，AI数据中心成为新引擎。}

\begin{itemize}
  \item 鸣鸣集团和万辰食品股价在5月底回调后反弹，两家公司通过官方渠道重申反对非理性竞争，补贴扩大是有纪律的优质加盟商筛选。
  \item 5月鸣鸣净增900+店，万辰净增1000+店，2Q25以来开店节奏明显提速，5M26两家合计新开超3000家。
  \item 万辰5月日均GMV转正（含新开店），鸣鸣1Q25单店GMV同比中个位数下滑，2Q预计仍有压力。
  \item 台湾气动元件龙头亚德客5月日均销售额同比增长26\%，环比仅微降1\%，订单价值连续多月超过出货价值，积压订单仍在积累。
  \item 机床工具行业销售额同比增长35\%，且较4月高位继续攀升，受益于AI数据中心投资增加，纺织机械和通用机械订单也环比改善。
  \item 电子和电池行业销售额同比分别增长12\%和55\%，但环比回落主要归因于季节性因素，光伏相关同比-5\%继续走弱。
\end{itemize}
{\small\color{refgray}\textbf{References:} [R018] 市场真正低估的不是需求，而是供给侧的再定价; [R025] 中国自动化需求的真实底色，被季节性波动掩盖了}

\section{汇率与新兴市场债券：定价机制的结构性转移}
\textbf{人民币中间价的定价机制正在经历静默转移，央行可能从被动对冲转向主动引导预期锚定。新兴市场债券曲线整体持稳，但个体分化剧烈，定价逻辑从利率跟随转向信用分化。}

\begin{itemize}
  \item 模型投影的美元兑人民币中间价系统性偏离实际公布价，模型误差从2025年1月的-1800个基点收窄至当前约600个基点，传统逆周期因子解释力下降。
  \item 最新模型预测USD/CNY中间价为6.7705，较上次预测低482个基点，计入逆周期因子后为6.7957。
  \item 新兴市场美元债券10s30s曲线整体仅收窄1bp至68bp，但印尼一个月内收窄26bp至9bp成为最平曲线，而多米尼加陡化11bp，埃及曲线斜率高达141bp。
  \item 美债10s30s曲线熊陡6bp，但EM曲线持稳甚至微幅收窄，表明EM定价锚正在从无风险利率曲线向信用基本面曲线切换。
  \item 亚洲地区是EM曲线走平的主要驱动力，从44bp大幅收窄到37bp，幅度达7bp。
\end{itemize}
\begin{figure}[htbp]
\centering
\includegraphics[width=0.88\linewidth]{figures/fig_332.jpg}
\caption{Figure 1 - Figure 1: Steepest and flattest 10s30s curves}
\end{figure}
\begin{figure}[htbp]
\centering
\includegraphics[width=0.88\linewidth]{figures/fig_333.jpg}
\caption{Figure 2 - Figure 2: Largest 1m changes in 10s30s}
\end{figure}
\begin{figure}[htbp]
\centering
\includegraphics[width=0.88\linewidth]{figures/fig_337.jpg}
\caption{Figure 7 - Figure 7: Historical EM aggregate 10s30s spread curve slope}
\end{figure}
\begin{figure}[htbp]
\centering
\includegraphics[width=0.88\linewidth]{figures/fig_338.jpg}
\caption{Figure 8 - Figure 8: Historical US treasury 10s30s yield curve slope}
\end{figure}
{\small\color{refgray}\textbf{References:} [R022] EM债券曲线正在传递一个被忽视的信号：期限溢价的结构性分化正在取代方向性交易; [R026] 人民币中间价的定价权，正在发生一次静默转移}

\section{风险偏好与资金流向：从beta交易到alpha分化}
\textbf{市场风险偏好整体降温，但资金流向呈现结构性分化。黄金流出掩盖了铜的逆势吸金，EM债券市场从beta交易转向alpha分化，投资者需从方向性交易转向挖掘个体差异。}

\begin{itemize}
  \item 大宗商品市场持仓总估值下降4\%至1.82万亿美元，黄金流出220亿美元，但COMEX Managed Money净多头反而增加。
  \item 铜市场净流入23亿美元，抵消铝和锌的流出，CTA净多头在COMEX铜上仍相对较高。
  \item EM债券市场10s30s曲线整体牛平1bp，但个体分化剧烈，印尼和埃及的曲线斜率相差15倍，均值掩盖了结构性事实。
  \item 美债曲线陡化而EM曲线持稳，背离表明EM信用风险定价没有随美债利率上升而恶化，投资者对EM信用风险评估更加独立。
  \item CTA净多头在COMEX黄金上稳定，但能源和农产品市场持仓下降，市场正从宏观叙事驱动转向微观供给结构重定价。
\end{itemize}
\begin{figure}[htbp]
\centering
\includegraphics[width=0.88\linewidth]{figures/fig_044.jpg}
\caption{Figure 1 - Figure 1: Global commodity markets saw net outflows largely driven by precious metals and energy markets USD billion, cumulative flows across tracked commodity markets}
\end{figure}
\begin{figure}[htbp]
\centering
\includegraphics[width=0.88\linewidth]{figures/fig_047.jpg}
\caption{Figure 5 - Figure 5: The CTA net long positioning in COMEX Gold stabilised over the week USD/oz (average weekly gold price); bubble size represents magnitude of weekly increase (green) or decrease (red) in CTA F\&O net length}
\end{figure}
\begin{figure}[htbp]
\centering
\includegraphics[width=0.88\linewidth]{figures/fig_051.jpg}
\caption{Figure 9 - Figure 9: The estimated value of open interest across base metals markets increased by 1\% WoW USD billion, estimated open interest across tracked base metals markets}
\end{figure}
\begin{figure}[htbp]
\centering
\includegraphics[width=0.88\linewidth]{figures/fig_337.jpg}
\caption{Figure 7 - Figure 7: Historical EM aggregate 10s30s spread curve slope}
\end{figure}
\begin{figure}[htbp]
\centering
\includegraphics[width=0.88\linewidth]{figures/fig_338.jpg}
\caption{Figure 8 - Figure 8: Historical US treasury 10s30s yield curve slope}
\end{figure}
{\small\color{refgray}\textbf{References:} [R006] 大宗商品市场正在经历的不是波动，而是定价权的结构性转移; [R022] EM债券曲线正在传递一个被忽视的信号：期限溢价的结构性分化正在取代方向性交易}

\section{结语}
当前市场最核心的特征不是方向性趋势，而是结构性分化。从宏观数据噪音到行业竞争格局，从资产定价逻辑到资金流向，分化正在各个层面加速。投资者需要从追逐beta转向挖掘alpha，从依赖宏观叙事转向关注微观供给结构的变化。
\newpage
\section{报告覆盖清单}
\begin{itemize}
  \item [R001] 宏观与利率：数据噪音中的结构性信号 - 瑞士手表出口暴跌16.6\%：关税扭曲下的数据噪音，掩盖了真正的复苏信号
  \item [R002] 宏观与利率：数据噪音中的结构性信号 - 市场低估的不是波动，而是期权偏斜信号在利率市场中的结构性价值
  \item [R003] 宏观与利率：数据噪音中的结构性信号 - 五月PMI回落至荣枯线，市场真正担心的不是增长放缓，而是价格传导断裂
  \item [R004] 宏观与利率：数据噪音中的结构性信号 - 市场真正低估的不是需求，而是供给侧的再定价
  \item [R005] 航空航天与防务：供给端的结构性重塑 - 航空航天与防务：市场真正低估的不是需求，而是供给端的结构性重塑
  \item [R006] 能源与大宗商品：定价权从宏观叙事向微观供给切换 / 风险偏好与资金流向：从beta交易到alpha分化 - 大宗商品市场正在经历的不是波动，而是定价权的结构性转移
  \item [R007] 未被主题摘要引用 - 市场真正低估的不是需求放缓，而是出口收缩对定价权的挤压
  \item [R008] AI/算力与半导体：从周期品到战略资产的再定价 - 市场低估了日本半导体设备的结构性机会，而非周期的尾声
  \item [R009] 医药与生物科技：ADC与双抗的交叉验证 - 肺癌治疗格局的真正变量，不在PD-1，而在ADC与双抗的交叉点
  \item [R010] 日本工程机械与IT服务：关税窗口与AI效率竞争 - 关税下调的真正价值不在当期利润，而在日本工程机械的竞争格局重塑
  \item [R011] AI/算力与半导体：从周期品到战略资产的再定价 - 市场正在低估“芯片通胀”的扩散半径：从AI算力成本到全球消费品定价
  \item [R012] 未被主题摘要引用 - 全球汽车市场正在经历的不是周期波动，而是竞争逻辑的根本切换
  \item [R013] 能源与大宗商品：定价权从宏观叙事向微观供给切换 - 铜价突破的关键不在于需求爆发，而在于供给侧的三个结构性断裂
  \item [R014] 中国新能源车：渗透率高位下的订单分化与出口爆发 - 新能源渗透率突破63\%的真正含义：市场正在重新定义赢家的门槛
  \item [R015] 中国新能源车：渗透率高位下的订单分化与出口爆发 - 新能源渗透率突破63\%后，市场真正低估的是新车型的订单转化效率
  \item [R016] 中国新能源车：渗透率高位下的订单分化与出口爆发 - 现代起亚正在悄悄改写美国汽车市场的竞争规则
  \item [R017] 地产与消费：二手房定价权转移与消费渠道重置 - 市场真正低估的不是需求，而是二手房定价权的结构性转移
  \item [R018] 中国平价零售与自动化：供给侧主动再定价 - 市场真正低估的不是需求，而是供给侧的再定价
  \item [R019] 能源与大宗商品：定价权从宏观叙事向微观供给切换 - 市场真正低估的不是需求，而是供给侧的再定价能力
  \item [R020] AI/算力与半导体：从周期品到战略资产的再定价 - 市场低估的不是需求，而是AI基础设施供应链的再定价
  \item [R021] 医药与生物科技：ADC与双抗的交叉验证 - ASCO 2026的真正信号：日本药企的差异化正在被重新定价
  \item [R022] 汇率与新兴市场债券：定价机制的结构性转移 / 风险偏好与资金流向：从beta交易到alpha分化 - EM债券曲线正在传递一个被忽视的信号：期限溢价的结构性分化正在取代方向性交易
  \item [R023] AI/算力与半导体：从周期品到战略资产的再定价 - 市场真正低估的不是AI需求，而是存储器供给侧的再定价
  \item [R024] 地产与消费：二手房定价权转移与消费渠道重置 - 消费支付数据回暖，但市场真正低估的是银行非息收入的修复弹性
  \item [R025] 中国平价零售与自动化：供给侧主动再定价 - 中国自动化需求的真实底色，被季节性波动掩盖了
  \item [R026] 汇率与新兴市场债券：定价机制的结构性转移 - 人民币中间价的定价权，正在发生一次静默转移
  \item [R027] 中国新能源车：渗透率高位下的订单分化与出口爆发 - 五月新能源车销量超预期，但市场真正低估的不是需求，而是供给侧的再定价
  \item [R028] AI/算力与半导体：从周期品到战略资产的再定价 - 市场真正低估的不是算力需求，而是AI基础设施的供给结构正在被重新定义
  \item [R029] 地产与消费：二手房定价权转移与消费渠道重置 - 市场真正低估的不是需求，而是供给侧的再定价
  \item [R030] 医药与生物科技：ADC与双抗的交叉验证 - KRAS G12C的真相：不是所有抑制剂都一样，组合疗法的门槛比想象中更高
  \item [R031] 医药与生物科技：ADC与双抗的交叉验证 - 肺癌一线治疗变局：真正值得关注的不是单一数据，而是两类新机制的“交叉验证”
  \item [R032] 日本工程机械与IT服务：关税窗口与AI效率竞争 - 市场低估的不是AI算力需求，而是IT服务商从“系统开发”到“代理设计”的职能跃迁
  \item [R033] 日本工程机械与IT服务：关税窗口与AI效率竞争 - 日本IT服务业的真正分化点：谁能在AI驱动的效率红利中，把成本节约转化为利润增长
\end{itemize}
\section{逐篇报告摘录}
\subsection{[R001] 瑞士手表出口暴跌16.6\%：关税扭曲下的数据噪音，掩盖了真正的复苏信号}
{\small\color{refgray}归类：宏观与利率：数据噪音中的结构性信号}

瑞士手表出口暴跌16.6\%：关税扭曲下的数据噪音，掩盖了真正的复苏信号

4月瑞士手表出口同比下滑16.6\%，这个数字足以让任何关注奢侈品行业的投资者心头一紧。但如果你只看这个标题就得出“全球奢侈品需求崩塌”的结论，很可能犯了方向性错误。

这份来自投行研报的核心判断是：当前出口数据的剧烈波动，本质上是关税政策引发的基数效应扭曲，而非终端需求的实质性恶化。真正值得关注的，不是4月单月数据本身，而是隐藏在噪音之下的三个结构性信号——美国市场在剔除关税扰动后的韧性、大中华区企稳的底部形态，以及中东地缘政治风险对奢侈品行业新增长极的威胁。

报告指出，4月出口同比跌幅较3月的-6\%（经工作日调整）显著扩大，而2月还是+9.2\%的正增长。这种过山车式的波动，恰恰印证了贸易政策对供应链节奏的深度干扰。去年4月，在“解放日”关税生效前，大量手表被提前发往美国，形成了极高的同比基数。今年4月对美出口暴跌56.4\%，正是这个基数的消化过程。如果以2024年为基准，美国市场4月仍录得+8.9\%的增长，虽然较1季度的+12.1\%略有放缓，但远非崩溃。

*市场真正低估的不是需求，而是数据噪音对判断力的系统性扭曲。**

1. 关税扰动正在制造一个“虚假的衰退”，但美国消费者已经消化了价格冲击

对美出口的剧烈波动是4月数据的核心变量。去年4月的抢运潮将出口基数推高至异常水平，今年自然面临同比回撤。但报告特别强调了一个关键细节：公司层面的反馈显示，美国消费者已经“在沉默中消化”了关税带来的涨价。

这意味着什么？意味着奢侈品公司担心的需求价格弹性断裂并未发生。高端腕表的核心客群，其消费决策对价格变动的敏感度远低于大众消费品。这一现象与过去几个季度奢侈品公司在美国市场的整体表现一致——美国仍然是全球奢侈品行业最稳定的需求支柱之一。

但这里有一个需要持续追踪的变量：关税扰动带来的基数效应将在7月和8月再次发力。去年的抢运高峰集中在这两个月份，届时同比数据可能再次出现剧烈波动。投资者需要警惕的是，不要被这些“技术性”的暴跌或反弹误导，而应该把目光聚焦在2024年为基数的连续对比上。报告没有完全展开的一个判断是：如果剔除关税噪音，美国市场的真实增速可能在个位数区间温和波动，这已经是一个相当健康的信号。

中国市场的表现可能是这份报告中最容易被误读的部分。4月对香港出口+13.5\%，对中国大陆出口+17.1\%，双双录得两位数增长。乍看之下，似乎中国奢侈品消费正在强劲反弹。

但报告明确指出，这同样是基数效应在起作用——去年4月，大量本该发往大中华区的货物被转向美国市场。如果以2024年为基准，4月出口规模仍比2024年同期低约15\%，与1季度（以2024年为基准）的走势基本一致。

这意味着什么？大中华区处于“底部震荡”而非“强劲反弹”阶段。15\%的缺口说明需求尚未恢复至疫情前水平，但也没有进一步恶化。这种“低位企稳”的状态，对于判断拐点至关重要——它意味着下行风险已经大幅释放，但上行催化剂尚未出现。

报告没有完全展开的一个问题是：中国奢侈品消费的真正复苏，需要依赖什么？是消费者信心的实质性修复，还是出境游的全面恢复，抑或是品牌定价策略的调整？这些变量目前都不明朗。大中华区最乐观的解读是“不再变得更差”，但这距离“重回增长”还有距离。

这是整份报告中我最想强调的一个结构性变化。中东地区在2025年曾是奢侈品行业罕见的增长亮点，但4月数据显示，对阿联酋出口下降9.5\%，对沙特下降17.3\%，对卡塔尔下降12.0\%，全面下滑。

报告将此归因于“持续的区域内冲突”。但更深层的含义在于：如果地缘政治风险持续发酵，中东可能从奢侈品行业的增量贡献者转变为拖累因素。这个转变的影响被严重低估，原因有二。

第一，中东高净值客群的消费力在过去两年显著提

[... middle omitted ...]

不够。

从价格带看，200-500瑞士法郎是唯一录得出口额正增长（+7.7\%）的区间，而3000瑞士法郎以上的高端区间跌幅最大（-19.0\%）。报告将此归因于基数效应，但这里可能隐藏着更深层的消费行为变化。

基数效应只能解释一部分：3000瑞士法郎以上区间去年4月增长19.8\%，基数确实最高。但200-500瑞士法郎区间的正增长，是否也反映了消费者在高通胀环境下的“降级消费”趋势？这个区间的产品，在奢侈品行业通常被定义为“入门级”或“轻奢”定位。如果高净值客群开始压缩高端消费，而中产客群仍然愿意为轻奢买单，这将是消费结构发生根本性变化的信号。

报告没有完全展开这个讨论，但它对投资决策有重要含义。如果消费分层的趋势成立，那么产品组合更侧重于高端和超高端定位的品牌（如百达翡丽、爱彼），可能面临比入门级品牌更大的需求压力。反之，产品矩阵覆盖更广、拥有更多入门级产品的集团（如斯沃琪集团），可能在当前环境下更具韧性。这一判断需要更多数据验证，但值得投资者将其纳入观察框架。

\subsection{[R002] 市场低估的不是波动，而是期权偏斜信号在利率市场中的结构性价值}
{\small\color{refgray}归类：宏观与利率：数据噪音中的结构性信号}

市场低估的不是波动，而是期权偏斜信号在利率市场中的结构性价值

利率市场参与者正在面对一个令人困惑的局面：通胀叙事尚未终结，地缘冲突此起彼伏，政策路径的可见度极低。在这样的环境中，大多数投资者本能地转向了更短的久期、更频繁的调仓、更依赖宏观叙事来形成判断。但一份来自某外资投行利率策略团队的研报提出了一个值得深思的相反观点——真正被市场低估的，不是对波动率本身的定价，而是期权偏斜（skew）作为一种独立信号，在利率方向性交易中的结构性预测能力。

这份报告的核心判断是：短期期权偏斜的变化，在长达17年的历史回测中，始终包含关于未来利率走势的预测信息，且这一关系在2019年之后不仅没有消失，反而在长端利率上变得更强。这不是一个样本内的偶然发现，而是一个经过120种参数组合压力测试后仍然成立的规律。

为什么这一点对当前市场尤其重要？因为当宏观冲击成为常态，传统的久期管理框架容易陷入“追着新闻跑”的被动状态。而偏斜信号提供了一种不同的视角——它不依赖对通胀或就业数据的预测，而是捕捉市场参与者通过期权市场表达出来的、对尾部风险的真实定价变化。这种变化往往领先于现货市场的方向性调整。

报告通过三个层次来验证这一主张：信号在不同参数设定下的稳健性、在不同市场阶段的表现分化、以及信号衰减的速度。以下是我们从中提炼出的关键洞察。

1. 偏斜信号的预测能力并非数据挖掘的产物，而是市场行为的真实映射

任何量化策略最容易被质疑的问题就是“过拟合”。报告对此做了非常扎实的处理：它没有去寻找一个最优参数组合，而是系统地测试了4种回溯窗口、3种期权期限、10种交易延迟，共计120种规格。结果发现，无论使用互换还是国债期货来表达久期暴露，无论选择1个月、3个月还是6个月的期权期限，偏斜信号的表现都高度一致。

这意味着什么？意味着偏斜所捕捉的，不是某个特定工具或期限上的定价异常，而是更底层的市场结构信息——投资者在期权市场上集中表达的尾部风险对冲需求。当这种需求发生系统性变化时，它往往预示着利率方向的转变。

从实践角度看，这一发现的意义在于：投资者不需要依赖复杂的参数优化或高频交易架构来利用这一信号。一个基于21天回溯窗口的简单策略，在2009年至今的样本期内，已经能够提供稳健的风险调整后回报。

2. 2019年之后，信号速度的重要性大幅上升，长端利率的偏斜信号反而更强

报告将样本划分为LIBOR时代（ 年中）和SOFR时代（2019年中至今），这一划分本身就揭示了市场结构性变化。在LIBOR时代，低利率、低波动的环境使得63天回溯窗口的表现与其他更短窗口相当。但在SOFR时代，63天窗口的表现显著恶化，而21天和更短窗口的策略保持了韧性。

这背后的逻辑是直观的：当市场冲击变得更加频繁和剧烈，过时的信号会迅速失去价值。2019年之后的利率市场，经历了疫情、通胀飙升、激进加息、地缘冲突等多重冲击，市场节奏明显加快。在这样的环境下，“快速反应”比“精确确认”更重要。

但报告最反直觉的发现在于期限结构的分化：2年期利率的偏斜信号预测能力在2019年后有所下降，而30年期利率的偏斜信号预测能力反而显著提升。报告给出的解释是，前端利率越来越多地被宏观和政策冲击主导，这类冲击往往在期权市场定价之前就已经通过现货市场反映出来。而长端利率仍然受到投资者持仓结构的显著影响，偏斜信号在这里捕捉到了更纯粹的“持仓拥挤度”信息。

这对当前市场的含义是：如果投资者试图利用偏斜信号来管理利率风险，长端可能比短端提供更好的信噪比。

3. 偏斜信号的信息不会立刻被市场吸收，这为战术性交易提供了窗口

报告的一个重要技术贡献是，它系统分析了信号衰减的速度。通过比较1天到10天不同交易延迟下的策略表现，报告发现：即使在延迟10天后建立仓位，策略在LIBOR

[... middle omitted ...]

场状态下都有效。当宏观冲击的规模和速度超过持仓调整所能反映的范畴时，偏斜信号可能暂时失灵。理解这一局限，恰恰是有效使用该信号的前提。

4. 报告尚未完全回答的关键问题：偏斜信号的失效边界在哪里？

报告已经做了非常详尽的稳健性检验，但它坦诚地指出了 年的回撤期，并承认这是信号最显著的失败案例。这个案例的价值在于，它揭示了偏斜信号的一个根本局限：当市场定价的驱动力从“投资者持仓调整”切换为“宏观预期跳跃”时，偏斜信号会暂时失去方向性。

从投资者的角度看，这意味着偏斜信号更适合作为“持仓拥挤度的温度计”，而不是“宏观方向的预言家”。当市场处于宏观冲击主导的阶段，偏斜信号可能提供误导性信息。而当市场进入“消化冲击、重新定价”的阶段，偏斜信号的价值会重新显现。

报告没有明确给出判断当前处于哪种状态的方法论，但这恰恰是最值得进一步探讨的问题。是否可以通过偏斜信号本身的形态（比如偏斜变化的速度、幅度、与其他市场的交叉验证）来区分这两种状态？这些是报告留下但未完全展开的伏笔。

\subsection{[R003] 五月PMI回落至荣枯线，市场真正担心的不是增长放缓，而是价格传导断裂}
{\small\color{refgray}归类：宏观与利率：数据噪音中的结构性信号}

五月PMI回落至荣枯线，市场真正担心的不是增长放缓，而是价格传导断裂

五月的中国制造业PMI数据，以精确的50.0画下一条分界线。这不仅仅是从扩张到收缩的临界点，更是一个信号：过去几个月由上游价格上涨驱动的经济叙事，正在遭遇下游需求的真实检验。

某外资投行最新发布的研报指出，五月官方制造业PMI较四月回落0.3个百分点，恰好落在荣枯线上。表面上看，这似乎只是经济复苏节奏的暂时放缓。但拆解分项数据后，一个更值得警觉的图景浮现出来：价格上涨的压力开始从上游传导至中游，但下游需求的承接能力正在减弱。这不是简单的供需错配，而是一轮潜在的价格-需求负反馈循环正在酝酿。

对于关注中国资产定价的投资者而言，五月PMI数据最核心的信息不是增长动能衰减，而是价格传导机制出现裂痕。这个判断，将直接影响我们对下半年政策节奏、行业利润分配以及资产配置逻辑的重新评估。

五月PMI最值得关注的变化，不是总量数字回落至50，而是生产指数与新订单指数之间的裂口在扩大。

生产指数录得51.2，虽然较四月回落0.3个百分点，但依然稳稳站在扩张区间。这意味着制造业企业仍在积极组织生产，供给端没有出现明显的主动收缩。然而，新订单指数从四月的50.6大幅下滑至49.9，五个月来首次跌破荣枯线。更令人警惕的是新出口订单指数，单月骤降1.7个百分点至48.6。

这两个数字放在一起，揭示了一个关键事实：生产跑在了需求前面，而且差距正在拉大。企业仍在按过去的订单节奏安排生产，但新订单的流入速度已经放缓。这种“生产惯性”与“需求冷却”之间的时间差，正是库存被动累积的典型前兆。

对于投资者而言，这个信号的含义是明确的。过去几个季度，市场对中国制造业的乐观情绪很大程度上建立在生产端持续扩张的基础上。但如果生产与订单的缺口持续扩大，那么下一阶段的焦点将从“产能利用率”转向“库存去化压力”。那些库存周期处于高位、且产品价格弹性较低的行业，可能率先面临盈利下修的风险。

研报中一个被重点讨论的指标是“输入价格-产出价格”的剪刀差。五月的输入价格指数录得60.5，产出价格指数录得51.9，两者之间的差距依然维持在8.6个百分点的较高水平。

这已经是两者连续第五个月同时处于扩张区间，但最关键的是，这个剪刀差并没有收窄。通常而言，当上游成本压力开始缓解时，输入价格会率先回落，产出价格则因滞后效应而维持相对高位，剪刀差自然收窄。然而，五月的数据显示，输入价格从四月的63.7回落至60.5，产出价格也从55.1回落至51.9，两者同步下行，但差距基本维持不变。

这意味着什么？企业不仅无法将上游涨价的压力完全传导出去，甚至在下游需求走弱的背景下，产出价格的下行速度比输入价格更快。中游制造企业正在被动承担成本压力，利润空间被双向挤压。

这种局面下，两个推演值得关注。其一，如果油价和其他工业品价格继续高位运行，而下游需求无法有效恢复，那么中游制造业的毛利率将在三季度面临系统性压力。其二，企业可能会在采购端做出更激进的调整——减少原材料采购、压缩库存、甚至主动减产。五月采购量指数从51.1降至49.8，已经暗示了这一趋势的开始。

3. 库存信号亮起黄灯：原材料去库、产成品累库的组合值得警惕

研报中一个容易被忽视但极具信号意义的组合是“原材料库存下降”与“产成品库存上升”的背离。

五月原材料库存指数从49.3降至48.6，而产成品库存指数从47.5跳升至49.3，单月上升1.8个百分点。这个组合是典型的“价格冲击下的企业防御行为”：面对居高不下的投入成本，企业减少上游采购以控制现金流；同时，下游销售放缓导致产成品被动积压。

需要强调的是，产成品库存指数目前仍在50以下，尚未进入绝对的累库区间。但方向比绝对水平更重要。过去一年，中国制造业一直处于相对健康的低库存状态

[... middle omitted ...]

经济结构转型的代价正在显现

五月PMI中为数不多的亮点是服务业PMI回升0.7个百分点至50.1，重返扩张区间。信息技术、金融服务、交通运输等细分领域均维持在55以上的高景气水平。这在一定程度上对冲了制造业的疲软。

然而，建筑业PMI虽然从四月的48.0回升至48.8，但依然连续第三个月处于收缩区间。建筑业是固定资产投资的前瞻指标，其持续低迷意味着此前市场预期的基建投资加速尚未落地。研报中提及的高频数据显示，截至五月，政府债券累计发行量仍比去年同期少5000亿元人民币。尽管五月发行节奏有所加快，但总量缺口依然存在。

服务业与建筑业的这种分化，本质上是经济结构转型的代价。服务业的高景气反映了消费升级和数字经济的韧性，但建筑业持续低迷则表明传统的投资驱动模式正在失效。对于投资者而言，这意味着两个判断需要调整：第一，不要对基建投资拉动经济有过高期待，财政发力的节奏和力度可能低于预期；第二，消费和服务业的结构性机会值得给予更高权重，但需要关注居民收入预期改善的可持续性。

\subsection{[R004] 市场真正低估的不是需求，而是供给侧的再定价}
{\small\color{refgray}归类：宏观与利率：数据噪音中的结构性信号}

这份某外资投行最新发布的Q1财报季分析报告，表面上看是在回答“企业盈利还能撑多久”，但细读下来，它真正揭示的是一组被市场定价忽略的结构性矛盾：消费者现金流的底层逻辑正在被改写，而AI投资的规模已经大到足以独立影响宏观经济账户中的资本形成。这两个方向的力量并不对冲，但它们对资产定价的含义截然不同。

报告覆盖了几乎所有标普500成分股，数据截止日接近完整。其核心判断可以浓缩为一句话：当前的经济韧性建立在两个脆弱支点之上——低储蓄率支撑的消费和超大规模企业主导的AI资本开支。前者面临油价、财政退坡和收入增速下滑的三重挤压，后者则几乎没有受到地缘风险的干扰，反而在加速。

这不是一份“乐观”或“悲观”的报告，而是一份关于“分化正在加剧”的报告。对于产业决策者和资产配置者而言，真正需要理解的问题不是经济会不会衰退，而是哪些部门正在被重新定价，以及这种重新定价的逻辑是否已经被市场消化。

财报数据显示，Q1消费者相关企业的营收增长依然强劲：标普500可选消费板块中位数公司营收同比增长9\%，必需消费板块也达到5\%。从官方统计和替代数据来看，名义消费支出增速在过去几个月保持稳定。

但这份报告真正有价值的洞察不在于这些表面数字，而在于它对消费现金流的底层拆解。报告明确指出，当前的消费韧性建立在一个不可持续的基础之上：个人储蓄率处于低位，实际可支配收入在4月个人收入报告中同比下降了1.1\%。这意味着消费者正在动用储蓄来维持支出节奏，而不是因为收入增长在加速。

报告预测2026年下半年的实际消费支出增速仅为1.3\%，显著低于市场一致预期。这个数字的背后是三个结构性因素的叠加：高油价对实际购买力的侵蚀、强劲退税效应的消退、以及通胀粘性对名义收入的持续消耗。值得注意的是，报告专门分析了不同收入阶层的分化——低收入群体的消费增速在过去一年已经从领先3.1个百分点收窄至1.1个百分点，而高油价对底层收入群体的冲击将更为显著。

这里有一个尚未被市场充分定价的含义：如果消费增速真的降至1.3\%，那么当前市场对消费类资产的盈利预期可能仍然偏高。尤其是那些依赖低收入客群的零售企业，其营收和利润的韧性可能比市场想象的要脆弱。

报告构建了一个独特的公司价格公告追踪器，用来捕捉企业层面的价格传导行为。这个指标在Q1升至2023年以来的最高水平，但仍远低于2022年的峰值。企业最频繁提及的价格上涨因素包括油价、航运成本、树脂基材料和计算机内存。

但真正值得关注的是报告揭示的一个传导机制：分析师对Q2利润率预期的下修幅度，与各行业价格公告追踪器的上升幅度高度相关。换句话说，那些面临最大输入成本压力的行业，其利润率指引的下调幅度也最大。这不是巧合，而是一个清晰的定价权信号——企业无法将全部成本上涨转嫁给消费者。

从行业分布来看，能源、必需消费、工业和非必需消费板块的价格公告追踪器上升最为明显，而这些板块的利润率预期也相应被下修。这意味着，即使整体通胀读数可能因为油价上涨而暂时走高，但企业层面的利润压力实际上在上升。

对于投资者而言，这个传导链条的含义是：油价上涨对资产价格的影响可能不是通过利率预期路径（即通胀走高导致加息预期升温），而是通过企业盈利路径。如果利润率持续承压，即使名义营收增长保持韧性，每股盈利的增长空间也将被压缩。市场目前是否充分定价了这一利润率风险，是一个需要持续跟踪的问题。

3. AI资本开支的规模已经超出“主题投资”的范畴，开始重塑宏观账户

报告在AI相关分析中给出了一个令人瞩目的数字：分析师对2026年超大规模企业资本开支的预期，自财报季开始以来被上调了12\%，达到超过7500亿美元，相当于较2025年增长83\%。这个数字已经远远超出了“科技主题”的范畴，开始对宏观经济账户中的企业固定投资产生实质性影响。

报告估算，

[... middle omitted ...]

少数几家公司的决策之上时，集中度本身就是一种脆弱性。

另一个尚未被充分讨论的问题是：AI资本开支的生产率回报何时能够体现。报告提到，目前关于AI对劳动力市场的影响仍然有限，只有极少数管理层在讨论裁员或招聘冻结时提及AI。但报告也指出，在科技公司和初创企业这个狭窄领域，将裁员与AI关联的频率在当月急剧上升，尤其是涉及运营、数据标注、软件开发和客户支持等岗位。然而，同一批公司在计算机硬件、系统工程和研究等岗位的招聘需求却在增加。

这意味着，AI对就业的影响不是简单的“替代”或“创造”，而是“再分配”。这种再分配对劳动力市场的结构性含义——包括技能溢价的变化、区域就业格局的重塑、以及工资通胀的部门分化——可能需要几个季度才能体现在宏观数据中。报告提出了这个问题，但没有给出完整的答案，这正是需要持续跟踪的关键变量。

整份报告最引人深思的部分，是它揭示了一个尚未被回答的问题：当消费增速放缓、输入成本上升、AI资本开支仍在加速这三个趋势同时发生时，企业利润率究竟会走向何方？

\subsection{[R005] 航空航天与防务：市场真正低估的不是需求，而是供给端的结构性重塑}
{\small\color{refgray}归类：航空航天与防务：供给端的结构性重塑}

航空航天与防务：市场真正低估的不是需求，而是供给端的结构性重塑

过去一周，某外资投行在纽约举办的战略决策会议上，与十二家航空航天与防务公司的CEO进行了密集交流。参会者覆盖了从波音、GE Aerospace到Anduril、L3Harris等商业航空与防务领域的核心玩家。会后发布的一份研报，表面上是对CEO观点的汇总，但深挖下去，它揭示了一个远比短期订单波动更重要的结构性判断：行业正在经历的不是一次周期性的需求起伏，而是供给端——从供应链结构、生产模式到竞争格局——的深层重构。市场当前的定价，可能更多反映了对需求端的担忧或乐观，却尚未充分消化供给端正在发生的、不可逆的变化。

这份报告的价值，不在于告诉读者“航空需求依然强劲”或“防务预算存在不确定性”这类已知信息。它的真正洞察在于，通过十二位CEO的口径，勾勒出了供给端几个同步发生、相互强化的趋势：商业航空领域，售后市场与OE生产同时高负荷运转，正在考验供应链的极限；防务领域，导弹需求的爆发式增长催生了全新的生产框架和合同模式；而Anduril等新进入者的崛起，则从根本上挑战了传统防务承包商的成本结构与交付速度。这些趋势叠加在一起，意味着行业龙头公司的竞争优势、盈利能力和估值逻辑都在被重新定义。

以下是我们从这份报告中提炼出的五个核心洞察。它们不是对报告内容的复述，而是试图回答一个更根本的问题：这些CEO的集体声音，究竟在告诉我们哪些市场尚未完全定价的变化？

1. 商业航空的真正压力点不在需求端，而在OE与售后市场同时高负荷运转下的供应链极限

报告中最容易被忽略的一个信息是，所有商业航空公司的CEO都表示，尽管燃油价格高企对航空公司盈利构成压力，但他们尚未看到任何需求疲软的迹象。售后市场的订单簿已经覆盖了未来12到18个月的产能，新机型的覆盖周期甚至更长。这听起来像是一个典型的“需求强劲”故事。但真正的故事并不在这里。

真正的压力点在于，OE（原厂设备）生产正在加速——波音确认737月产47架的里程碑已经通过FAA审核，并计划在夏季全面达成，随后向52架、57架甚至63架迈进——而售后市场的需求并未因OE增产而减弱。这两个市场同时高负荷运转，意味着供应链必须同时满足两套互不兼容的节奏：OE要求准时、大批量、标准化的交付，而售后市场则要求灵活、小批量、高时效的备件供应。报告明确指出，这种叠加正在给供应链带来“压力”。但更准确的表述是，它正在暴露供应链的瓶颈。

以发动机为例。供应商正在消化库存，使交付与生产速率更加对齐，但报告也承认，某些部件的库存消化仍需时日。这意味着，在供应链的某些节点，产能的物理瓶颈尚未解除。对于投资者而言，关键问题不是“需求会不会下降”，而是“在需求不下降的前提下，供应链能否支撑当前的增速”。如果答案是否定的，那么接下来的故事就不是需求端的放缓，而是供给端的瓶颈导致交付延迟、成本上升，进而压缩利润率。报告没有给出这个判断，但所有信号都指向这个方向。

2. 导弹需求的“3倍到4倍”增长框架，正在重塑防务供应链的价值分配

防务领域最引人注目的信号来自导弹和导弹防御。报告提到，Lockheed Martin和RTX都描述了“非常强劲”的需求前景，并建立了新的生产增长框架，计划在未来五到七年内将当前产量提升3倍到4倍。更关键的是，这些合同预计将包含“增长保障条款”——如果需求未能如期兑现，政府将对公司进行补偿。这意味着，对于参与这些框架的公司而言，增长的下行风险被显著对冲了。

但这一趋势的更深层含义在于，它正在重塑防务供应链的价值分配。报告特别提到，对于Howmet和ATI这样的材料供应商而言，导弹市场正在成为一个新的增长领域，而此前导弹并不是它们的重要市场。这听起来像是一个简单的“新增业务”故事，但它的结构含义是：随着导弹生产从

[... middle omitted ...]

和产业界都意识到了这种风险，并试图通过合同设计来分担。这本身就是一个重要的信号：导弹生产正在从“采购”模式转向“产业政策”模式。

3. 新进入者的成功不是偶然，而是对传统防务成本结构的系统性颠覆

Anduril的CEO Brian Schimpf在会议上的发言，可能是整份报告中最具颠覆性的内容。Anduril目前约有1万名员工，年收入在2025年超过20亿美元，预计2026年翻倍，五年内可能达到300亿到400亿美元。这样的增长速度在传统防务领域几乎是不可想象的。

但真正值得关注的不是增长数字本身，而是Anduril实现这种增长的逻辑。Schimpf明确指出，Anduril的优势在于它没有继承传统防务承包商“追求极致性能、忽视可制造性和可扩展性”的历史包袱。它通过内部整合关键组件——包括固体火箭发动机、作动器、航电和制导系统——来降低对低层级供应商的依赖。当大部分价值由集成商攫取时，低层级供应商往往缺乏投资动力，而Anduril的垂直整合模式恰好解决了这个问题。

\subsection{[R006] 大宗商品市场正在经历的不是波动，而是定价权的结构性转移}
{\small\color{refgray}归类：能源与大宗商品：定价权从宏观叙事向微观供给切换 / 风险偏好与资金流向：从beta交易到alpha分化}

过去一周，全球大宗商品期货市场的持仓总估值下降了690亿美元，至1.82万亿美元。表面上看，这是一次由能源价格下跌和黄金资金外流共同触发的市场回调。但这份来自某外资投行的研报揭示了一个更值得关注的信号：市场正在从“宏观叙事驱动”向“微观供给结构重定价”切换。真正重要的不是价格涨跌了多少，而是哪些资产在资金流出时仍然吸引了增量资金，哪些市场在价格下跌的同时持仓结构发生了质变。

这份报告最核心的判断可以概括为一句话：大宗商品市场的下一轮机会，不在于押注通胀或衰退的方向，而在于识别那些供给端已经发生不可逆结构性变化的品种。能源市场正在经历需求预期的修正，但LNG的供给弹性正在耗尽；黄金的资金流出掩盖了基本金属中铜的逆势流入；农产品则陷入了持仓稳定但价格信号分化的“假性平衡”。这些信号叠加在一起，指向一个结论：当前的市场定价，低估了结构性供给约束对部分品种的长期支撑，也高估了宏观情绪对整体资产类别的统摄力。

1. 能源市场的下跌不是需求崩溃，而是中国需求结构的永久性调整

过去一周，能源市场持仓总估值下降560亿美元，至8160亿美元，布伦特和WTI原油价格分别下跌11\%和10\%。市场容易将这一跌幅解读为全球需求放缓的信号，但报告提供了一个更细致的视角：中国石油需求下降了9\%，但背后的经济活动的下降幅度有限。这意味着，中国石油需求的下降并非周期性的需求萎缩，而是一次结构性的消费模式转变——汽油、柴油和燃料油的需求损失中有相当一部分将是永久性的，而航空燃油和石脑油的需求则会大致恢复。

这一判断的含义是深远的。市场目前对原油的定价，仍然隐含了一个假设：中国需求将在经济复苏后回到趋势线。但如果中国石油需求的结构性下降已成定局，那么全球原油市场的长期供需平衡点就需要重新校准。报告没有直接给出这一结论，但从数据推导，这可能是未来6-12个月原油定价中最容易被忽视的变量。

与此同时，LNG市场正在经历一个相反的供给故事。报告指出，全球LNG市场已经吸收了约60\%的霍尔木兹海峡供给冲击，依靠的是北美和非洲新项目的增产。但关键不在于已经吸收了60\%，而在于剩余的40\%将如何消化。报告明确表示，随着市场灵活性消退，且欧洲储气量远低于平均水平，价格将在注入旺季上升。这意味着，LNG市场的供给弹性正在快速收敛，而需求侧的价格敏感度有限——这是一个典型的供给约束型定价环境。

贵金属市场持仓总估值下降220亿美元，至2640亿美元，连续第三周下降。其中黄金市场净合约流出高达220亿美元，是所有商品板块中资金流出最显著的。但报告同时指出，COMEX黄金期货的Managed Money净多头头寸实际上增加了3900张合约，达到97400张净多。全球黄金ETF持仓量仍为4131吨，年初至今增长3\%。

这两组数据放在一起，揭示了一个重要的资金行为模式：流出黄金的主要是其他交易者类型（如生产商对冲、商业持仓等），而投机性资金仍在增加多头。这意味着，黄金的持仓下降并非避险需求的结构性消退，而是一次资金在不同交易者类型之间的再分配。投机者看多，而商业账户在减仓——这通常发生在价格高位附近，但并不等同于趋势逆转。

更值得关注的是，黄金ETF在4月和5月经历了新一轮卖出，但年初至今仍然是净流入。这暗示散户和机构投资者对黄金的配置意愿仍然存在，只是节奏上出现了调整。报告没有进一步分析黄金流出背后的宏观驱动力，但结合报告开头提到的美联储紧缩预期升温，一个合理的推断是：实际利率预期上行正在压制黄金的短期表现，但地缘风险和央行购金的长期趋势并未改变。

与能源和贵金属的明显流出形成对比，基本金属市场持仓总估值小幅上升18亿美元，至2390亿美元。更重要的是，基本金属是少数几个实现净合约流入的板块之一——铜市场流入23亿美元，完全抵消了铝和锌的流出。

铜的

[... middle omitted ...]

给结构性的确定性约束。

锌的情况则提供了一个反面案例。报告指出，尽管锌价预计将维持在 美元/吨的高位区间，以激励中国出口，但风险偏向下行。原因是中国出口可能超预期，而通胀压力可能抑制需求。这是一个典型的“供给支撑价格，但需求无法跟上”的格局。锌的案例提醒我们，供给约束并不自动等于投资机会——需求端的弹性同样关键。

农产品市场持仓总估值基本持平，为3930亿美元。净合约流入40亿美元，主要来自大豆市场，但被玉米、小麦、糖和棉花的价格下跌所抵消。表面上看，这是一个稳定的板块，但报告的价格动量信号揭示了一个更复杂的图景：柴油、TTF天然气、大豆和糖的短期价格动量已经转为卖出信号。

农产品市场的这种“总量稳定、内部分化”的状态，本质上是一种假性平衡。不同品种面临的供需格局完全不同：大豆有资金流入，但价格动量转负；玉米和小麦价格下跌，但净持仓变化不大；糖和棉花既无资金流入也无价格支撑。这意味着，农产品板块不再是一个可以整体配置的资产类别，投资者必须下沉到品种层面进行判断。

\subsection{[R007] 市场真正低估的不是需求放缓，而是出口收缩对定价权的挤压}
{\small\color{refgray}归类：未被主题摘要引用}

5月PMI数据已经全部落地。无论是国家统计局发布的官方PMI，还是某外资投行追踪的RatingDog PMI，都指向同一个方向：制造业扩张仍在继续，但速度明显放缓。这个结论本身并不令人意外，一季度超预期的增长之后，市场对二季度自然有回落的心理准备。

但真正值得关注的，不是扩张放缓这个事实，而是放缓的结构性特征——以及这些特征对下半年资产定价意味着什么。

这份报告的核心判断是：当前PMI读数所反映的，不是一次温和的周期性回调，而是一个正在发生的结构性切换——从内需驱动转向外需承压，从成本传导转向利润挤压。市场对后者的定价可能仍然不足。

5月RatingDog PMI中的新出口订单指数从51.1降至49.6，直接跌入收缩区间。这不是一个孤立的信号。国家统计局PMI中的新出口订单同样下行，从50.3降至48.6。两组数据指向同一个事实：外部需求正在走弱。

报告明确指出，这与航运数据以及最新的官方PMI读数一致。一季度出口对经济增长的拉动作用非常显著，但5月的数据暗示，这个引擎可能正在减速。更关键的是，出口订单的收缩不是边际性的——1.5个点的月度降幅，在PMI体系中已经属于比较显著的变化。

这意味着什么？对于依赖出口的制造业企业来说，二季度末到三季度可能面临一个双重压力：一方面是订单量的收缩，另一方面是价格传导能力的减弱。那些在一季度受益于海外补库需求的行业，可能需要重新评估下半年的营收预期。

5月的一个积极信号是通胀压力有所缓和。RatingDog的投入价格指数从57.5降至55.8，产出价格指数从53.6降至52.0。官方PMI的价格分项也呈现类似趋势。

但仔细看结构，问题在于缓解的方式。投入价格和产出价格之间的差距仍然存在，而且这个差距是以企业主动压缩利润空间的方式来弥合的。报告中的表述很克制：“producers absorbing a greater share of the upstream price shocks”——翻译过来就是，企业正在消化上游涨价，而不是完全向下游传导。

这不是一个可持续的状态。如果出口订单继续走弱，企业将面临更激烈的价格竞争，利润空间的压缩可能会加速。对于投资者而言，这意味着需要重新审视那些依赖出口且定价能力较弱的中游制造企业的盈利预期。

5月RatingDog的未来产出PMI从55.7降至55.0，官方PMI的同类指标也小幅回落至53.9。虽然两者都保持在扩张区间，但趋势是向下的。

报告提到，企业仍然对新产品、技术突破和产能改善保持乐观。但值得注意的是，调查中已经明确反映出对“外部需求放缓和成本上升”的担忧。这意味着，当前的乐观预期可能更多是基于在手订单和既有产能的惯性，而不是对新订单加速的信心。

历史数据表明，未来产出PMI领先于实际生产活动大约1-2个季度。如果这个指标在未来两个月继续下行，三季度实际生产活动放缓的速度可能会比市场当前预期的更快。

4. 就业指标持续在荣枯线附近徘徊，这是一个需要跟踪的隐性变量

5月RatingDog的就业PMI从49.9微降至49.8，官方PMI的就业分项也基本持平。两者都保持在略低于50的水平，连续多月没有明显改善。

就业指标的粘性值得关注。一方面，制造业就业通常滞后于生产活动变化，当前的数据可能已经反映了前几个月的生产扩张。另一方面，如果出口订单收缩开始传导至生产端，就业指标的下行压力可能会在三季度加速。

对于政策制定者来说，就业是比GDP增速更敏感的决策变量。如果就业PMI在未来两个月出现加速下滑，可能会触发新一轮稳增长政策的出台。对于市场而言，这意味着需要提前布局那些可能受益于政策对冲的板块。

5. 报告尚未完全回答的一个关键问题：地缘政治风险的二阶影响

报告提到了一个重要的背景变化：“Opera

[... middle omitted ...]

。这些调整的成本是结构性的，不会因为短期局势缓和而逆转。对于依赖特定原材料进口或特定区域市场的企业，这个变量的影响可能才刚刚开始被定价。

基于这份报告的逻辑，我们建议读者在接下来的两个月里，将观察重心从“制造业扩张还是收缩”转向“扩张的结构是否健康”。具体来说，可以关注三个维度：

第一，出口订单指数能否在6月企稳。如果继续下行，三季度增长预期可能需要进一步下修。第二，价格剪刀差（投入价格减去产出价格）是否继续收窄。如果收窄主要是通过产出价格更快下降实现的，说明企业的定价权在流失。第三，就业指标是否出现加速下滑。这既是经济健康度的信号，也是政策预期的先行指标。

这三个维度组合在一起，可以形成一个比单一PMI读数更丰富的判断框架。

这份报告还有一个值得深入讨论的部分，是关于“特朗普-习近平峰会”对贸易不确定性的潜在影响。报告认为这可能“注入稳定性”，但这个判断取决于很多尚未落地的细节。我们会在社群中继续拆解这个变量的几种可能情景，以及它们对不同行业资产定价的含义。

\subsection{[R008] 市场低估了日本半导体设备的结构性机会，而非周期的尾声}
{\small\color{refgray}归类：AI/算力与半导体：从周期品到战略资产的再定价}

当市场将目光聚焦于AI芯片和先进制程的激烈竞争时，一份来自某外资投行的最新研报揭示了一个被广泛忽视的信号：日本半导体设备（SPE）行业正在经历的不是周期性的尾声，而是一轮由存储资本开支和先进封装驱动的结构性增长拐点。2026年4月的数据显示，日本SPE出货额同比增长25\%（以日元计），其中组装设备同比暴增73\%，测试设备增长26\%。这些数字并非简单的月度波动——它们指向一个更深层的产业再平衡：全球晶圆厂设备（WFE）市场的增长引擎正在从逻辑芯片向存储芯片切换，而日本设备商正是这一切换的最大受益者。

该报告预测，全球WFE市场在2026年将同比增长21.4\%，2027年再增长18.2\%。这一判断的核心逻辑并非AI需求的线性外推，而是存储厂商（尤其是DRAM和NAND）在经历了两年的资本开支紧缩后，正在启动一轮大规模的产能扩张。对于投资者而言，真正重要的不是判断AI需求是否见顶，而是理解为什么日本设备商在这一轮扩张中拥有比市场预期更强的议价权和份额提升能力。

4月单月日本SPE出货额环比下降38\%，这一数字很容易被解读为需求走弱的信号。但研报通过3个月移动平均数据清晰地指出，这仅仅是季节性因素导致的基数效应——3月通常是日本财年末的出货高峰，4月自然回落。以3个月均值衡量，4月出货额环比仍增长6\%，同比增长11\%（以日元计为14\%），延续了自2023年中开始的上升趋势。

更值得关注的是设备品类的结构性分化。前道设备（主要对应东京电子TEL）同比增长12\%，增速稳健但并未超预期；组装设备（对应DISCO）同比增长73\%，虽然部分归因于低基数，但这一增速仍反映了先进封装需求的持续爆发；测试设备（对应爱德万测试Advantest）同比增长26\%，且3个月平均增速更为强劲。

这一分化揭示了一个关键事实：市场对前道设备的预期已经较为充分，但组装和测试设备的需求弹性被系统性低估。后两者直接受益于HBM（高带宽内存）和CoWoS（晶圆级封装）的产能扩张，而这恰恰是AI芯片从“设计”走向“量产”的关键瓶颈环节。

研报最引人注目的部分是运用回归模型对两家核心公司季度收入的预测。基于SEAJ（日本半导体设备协会）的出货数据与公司历史收入的强相关性（R²分别达0.83和0.95），该模型给出了与市场共识截然不同的判断。

对于东京电子，模型预测其6月季度（FQ1）收入将环比下降18\%，而市场共识预期为增长6\%。这一预期差高达24个百分点，意味着市场可能严重高估了TEL在短期内的收入弹性。研报指出，TEL的季度首月通常为慢速月，后续两个月有望改善，但即便如此，4月单月数据的疲软仍可能导致整个季度不及预期。对于投资者，这意味着TEL在接下来几周的业绩指引中可能释放低于预期的信号，从而带来短期股价压力。

然而，对于爱德万测试，模型预测其6月季度收入将环比增长13\%，远高于市场共识的3\%。尽管4月单月数据环比下降14\%，但这是建立在3月创纪录的高基数之上——4月仍然是历史上第二高的单月出货额。研报强调，爱德万测试作为HBM测试仪（市场份额约65\%）和英伟达AI GPU测试的独家供应商，其收入增长不仅受益于测试强度的提升，更受益于产品代际迁移（如HBM4）带来的ASP（平均售价）提升。这一判断意味着，市场对爱德万测试的盈利预期可能在未来几个季度持续上修。

研报对全球WFE增长的预测（2026年+21.4\%，2027年+18.2\%）建立在一个关键假设之上：存储厂商的资本开支将在2026年显著加速。这一判断并非空穴来风——过去两年，存储芯片厂商（尤其是DRAM和NAND）在库存调整和需求疲软的压力下大幅削减了资本开支，导致设备采购被压抑。随着AI对HBM和3D NAND需求的拉动，以及传统存储市场的触底反弹，一轮补偿性的资本开支周期已经

[... middle omitted ...]

4. 报告尚未完全回答的关键问题：中国市场的替代效应与估值边界

尽管研报对日本设备商给出了清晰的乐观判断，但有两个关键问题尚未得到充分论证。第一，中国本土设备商（如中微公司、北方华创、拓荆科技）的快速崛起是否会在中长期侵蚀日本设备商在中国市场的份额？研报在评级中维持对中国三家设备商（北方华创、中微公司、拓荆科技）的“跑赢大市”评级，但并未详细讨论日本设备商与中国本土厂商在技术路线、客户关系和定价策略上的直接竞争。考虑到中国目前仍是全球最大的半导体设备市场之一，这一替代效应的速度和幅度将直接影响日本设备商的收入增长假设。

第二，当前估值是否已经反映了上述乐观预期？以研报给出的目标价和2026年预测盈利计算，TEL的远期市盈率约为35.7倍，DISCO为37.0倍，爱德万测试为35.8倍。这些估值水平在历史区间中已处于较高位置。研报虽然指出了业绩上修的可能性，但并未量化“估值消化”所需的时间。对于长线投资者而言，这意味着需要仔细权衡增长确定性与当前估值之间的安全边际。

\subsection{[R009] 肺癌治疗格局的真正变量，不在PD-1，而在ADC与双抗的交叉点}
{\small\color{refgray}归类：医药与生物科技：ADC与双抗的交叉验证}

肺癌治疗格局的真正变量，不在PD-1，而在ADC与双抗的交叉点

2026年ASCO之后，市场对肺癌领域最值得关注的判断不是某款药物数据“好”或“坏”，而是一个结构性信号：ADC与PD-1/L1xVEGF双抗正在从替代选项变为潜在的基石疗法，但全球数据验证才是真正分水岭。

某外资投行在ASCO结束后迅速组织了两位肺癌领域KOL的深度讨论。这两位KOL分别来自约翰霍普金斯大学和斯坦福大学，覆盖了从临床实践到早期试验设计的完整视角。他们的判断并非简单的“看好”或“看空”，而是给出了一个更有层次的分析框架：哪些信号已经足够清晰，哪些仍需等待全球数据，以及哪些未解问题将决定未来三到五年的竞争格局。

这份讨论的核心价值不在于复述数据，而在于揭示了市场目前定价最不充分的两个维度——毒性管理能力对市场份额的直接影响，以及中国数据向全球人群的转化风险。以下是我们从KOL讨论中提炼的五个关键洞察。

1. sac-TMT的PFS数据已足够改写一线标准，但OS才是真正的定价锚

KOL对科伦博泰与默沙东合作的sac-TMT（TROP2 ADC）在OptiTROP-Lung05研究中的表现给出了高度评价。一位KOL用了一个生动的比喻：PFS的HR数据在他看过的药物中属于“少数几个真正显著的”，如果用棒球术语来量化，0.5以下的HR就是“全垒打”。而sac-TMT在PD-L1高表达人群中的表现，已经接近这个区间。

但KOL的谨慎同样值得关注。他们明确指出，Keytruda单药作为对照臂在全球范围内并非标准治疗（美国更倾向Keytruda联合化疗），因此中国数据的对照组表现可能被低估。这意味着sac-TMT的PFS优势在全球试验中可能会被压缩。

更关键的是OS。一位KOL给出了明确的临床决策门槛：对于sac-TMT下一阶段的数据，他需要看到HR在0.7到0.8之间、响应率在65\%到70\%之间，才能考虑替换标准治疗。而另一位KOL更严格，认为临床有意义的HR门槛是0.75。这些数字不是学术讨论，而是未来商业化的定价锚——如果全球OS数据达不到这个区间，即便PFS再漂亮，医生也不会轻易改变处方习惯。

Summit Therapeutics与康方生物合作的PD-1xVEGF双抗ivonescimab在HARMONi-6研究中展现了令人印象深刻的OS获益（HR=0.66），且PFS到OS的衰减很小。KOL认为这“意味着ivonescimab是一种有效的疗法”，并指出其在出血风险患者中的安全性优于传统VEGF抑制剂贝伐珠单抗。

但KOL同时点出了一个关键争议：65岁以上患者亚组的获益似乎不明显。这个细节在ASCO的全体讨论环节已被提出，而KOL的解读更为直接——中国人群的吸烟率、病理生理特征与全球人群存在差异，而全球试验（HARMONi-3）通常包含更多老年患者。如果这个亚组差异在全球数据中重现，将直接影响ivonescimab在真实世界中的使用范围。

KOL对HARMONi-3的鳞状细胞癌队列持乐观态度，认为基于ivonescimab在多个适应症中的一致性阳性数据，该队列的PFS最终分析（2026年下半年）大概率成功。但他们对非鳞状细胞癌的跨组织学解读持保守态度，认为鳞状细胞癌更均质、吸烟者比例更高，因此获益预期更强，非鳞状细胞癌的转化需要更谨慎。

*这里的核心含义是：市场可能低估了双抗在鳞状细胞癌中的竞争优势，但高估了其在非鳞状细胞癌中的可替代性。**

3. 中国数据的“可翻译性”是最大的隐性风险，也是最大的超额收益来源

两位KOL在讨论中反复提及一个主题：中国数据能否在全球人群中复现。这不仅是监管问题，更是临床实践的根本差异。

对于sac-TMT，KOL认为其安全性特征与第一三共/阿斯利康的Datroway相似，但基于

[... middle omitted ...]

行空间。这本质上是一个“可翻译性溢价”的期权。

4. 毒性管理正在成为ADC竞争的新护城河，但个体差异仍是黑箱

KOL对sac-TMT的毒性管理给予了积极评价，特别是其约5\%的治疗中断率被描述为“惊人”。但他们的讨论揭示了一个更深层的问题：ADC的毒性并非均质化风险。

一位KOL指出，患者对ADC的个体反应差异很大，同样的药物在不同患者身上可能表现出完全不同的耐受性。这意味着，即使临床试验中展现了良好的平均安全性，在真实世界中，医生仍需面对大量“黑箱”决策——谁能耐受、谁不能耐受，目前没有可靠的预测生物标志物。

这一点对商业化的含义是：ADC的竞争不仅看疗效数据，更看医生对毒性管理的“教学曲线”。KOL提到，当与患者讨论ADC时，他会预期化疗样副作用，并认为这些是可控的。但关键在于，如果某款ADC出现了独特的副作用（如间质性肺炎、眼部毒性），医生的使用意愿会大幅下降。sac-TMT目前没有发现这些“独特副作用”，这是其优势，但全球试验中是否会出现，仍是未知数。

\subsection{[R010] 关税下调的真正价值不在当期利润，而在日本工程机械的竞争格局重塑}
{\small\color{refgray}归类：日本工程机械与IT服务：关税窗口与AI效率竞争}

关税下调的真正价值不在当期利润，而在日本工程机械的竞争格局重塑

2026年6月2日，白宫一则关于修改钢铁和铝产品232条款关税的总统公告，在东京早盘交易时段引发日本机械板块的异动。表面看，这只是一次针对特定工业品类的税率调整——联合收割机、推土机、叉车等设备的对美出口关税从25\%降至15\%。但仔细拆解这份公告的生效时间、适用品类范围以及日落条款，会发现一个更值得关注的信号：美国正在有选择地降低某些工业品的进口壁垒，而其真实意图并非简单的贸易缓和，而是为重建本国工业基础争取时间。

投行研报对此事件做了量化分析，估算出对主要日本厂商当期利润的影响——小松约100亿日元增量，久保田和竹内制作所影响有限。但如果我们只停留在利润测算层面，就错过了这份报告真正有价值的部分：关税下调的结构性含义，以及它对日本工程机械行业竞争格局的长期影响。

研报明确指出，新关税税率适用于6月8日及之后抵达的货物，这意味着实际体现在终端销售价格上要等到11月至12月以后。更关键的是，这项政策的日落日期是2027年12月31日。

这不是一个永久性的税率调整，而是一个为期约18个月的窗口期。美国选择在这个时间点降低工程机械的进口关税，同时维持对钢铁和铝的基础税率不变，反映出一种精细化的政策逻辑：美国需要日本的高端工程机械来支撑其国内基建和制造业投资，但又不希望长期依赖进口。这18个月，恰好是美国本土产能爬坡的时间窗口。

对于日本制造商而言，这个窗口期意味着两件事。第一，短期利润增厚是确定的，但幅度有限——研报测算的100亿日元增量对于小松这样的体量而言，约占其年营业利润的2\%-3\%，算不上质变。第二，真正重要的是，这18个月将测试日本企业能否利用关税优势扩大市场份额、深化客户关系，从而在2027年之后即便关税恢复，也能维持更强的议价权。

2. 品类范围的扩大，揭示了一个正在被重新定义的“工业机械”边界

研报特别提到，此次公告扩大了适用15\%税率的工业机械品类范围，将推土机、叉车等移动工业设备纳入其中。这比单纯降低几个百分点的税率更有意义。

传统上，工程机械的关税分类往往按“农业机械”“建筑机械”“工业车辆”等传统标签划分。但此次调整显示出美国正在用一种更功能性、更产业化的视角来定义“工业机械”——凡是直接服务于工业基础重建的设备，都被纳入优惠范畴。这意味着，那些能够提供“设备+服务+数字化解决方案”的日本企业，将比单纯卖硬件的企业获得更大的政策红利。

对于投资者而言，需要关注的不是某一家公司有多少产品品类被纳入，而是这些品类的边界扩展是否意味着美国市场对日本工程机械的“必需品”属性在上升。如果答案是肯定的，那么即便2027年关税恢复，需求的刚性也会对冲掉大部分成本冲击。

研报估算，由于货物运输和销售周期，关税下调对当期利润的贡献要等到11月至12月才能体现。这看似是一个技术性的时间错配，但背后隐藏着一个更重要的分析维度：企业能否在这段时滞期内，将关税下降的红利转化为定价策略的优化。

理想情况下，日本企业应该维持原价销售，直接享受利润率提升。但在竞争激烈的北美市场，经销商和终端客户很可能会要求价格让利。因此，最终能有多少利润增量真正落到日本制造商的报表上，取决于三个因素：经销商库存水平、竞争对手的定价反应、以及终端需求的弹性。

研报没有给出明确的结论，但提供了一条分析线索：小松的利润增量估算为100亿日元，而久保田和竹内制作所影响有限。这暗示了不同企业在产品组合、经销商关系和市场地位上的差异。小松在北美市场更深的布局和更强的品牌溢价，可能是它能够更大程度保留关税红利的原因。这一点值得进一步验证。

4. 研报没有完全回答的关键问题：关税下调是否会触发北美市场的价格战

这是整份研报最令人意犹未尽的地方。研报详细测算了利润影响

[... middle omitted ...]

验证这两个假设，但我们可以从行业结构做一些推断。

北美工程机械市场是一个寡头格局，卡特彼勒和日本厂商之间的竞争长期保持一种微妙的平衡。关税下调可能打破这种平衡，但幅度有限。更可能的情景是，日本企业不会全面降价，而是选择在特定细分市场（如大型推土机、液压挖掘机）进行针对性促销，同时利用关税红利加大对经销商网络的投入。这样既不会引发全面价格战，又能巩固竞争地位。

5. 对投资者的观察框架：从“关税受益”转向“关税转化能力”

基于以上分析，我们认为投资者应该建立一个新的评估框架，而不是简单地给所有日本工程机械股贴上“关税受益”的标签。

这个框架的核心是“关税转化能力”——企业将关税下调转化为可持续竞争优势的能力。具体可以从三个维度衡量：

第一，北美市场收入占比和产品组合的差异化程度。占比越高、产品越差异化，企业对关税红利的掌控力越强。

第二，经销商网络的深度和忠诚度。关税下调带来的利润空间，可以用来改善经销商返点、延长保修期或提供融资支持，这些都会增强渠道粘性。

\subsection{[R011] 市场正在低估“芯片通胀”的扩散半径：从AI算力成本到全球消费品定价}
{\small\color{refgray}归类：AI/算力与半导体：从周期品到战略资产的再定价}

市场正在低估“芯片通胀”的扩散半径：从AI算力成本到全球消费品定价

过去十二个月，DRAM价格累计上涨超过六倍。如果只看这个数字，很容易将其归入“半导体周期正常波动”的叙事框架。但一份来自某外资投行研究团队的最新报告提出了一个更具结构性的判断：我们正在经历的，不是一次普通的存储器上行周期，而是一场由AI需求驱动的、跨行业、跨地域的“芯片通胀”开端。

这份报告的标题直指核心——Chipflation。它试图回答一个正在从产业界蔓延到宏观政策圈的问题：当存储器从“大宗商品”变成“战略资产”，其价格重估的冲击波，究竟会传导多远？

我们仔细研读了这份长达数十页的报告，以下是我们提炼出的核心洞察与尚未被市场充分定价的关键变量。

报告最核心的论点在于，本轮存储器价格上涨的驱动力，与传统半导体周期有本质区别。过去的涨价往往源于供给端的短期收缩或下游的补库存行为，价格会在需求正常化后迅速回落。但这一次，需求端的结构性变化是主导因素。

AI基础设施的建设，正在以前所未有的速度吞噬存储器产能。以HBM（高带宽存储器）为例，它是AI加速器的核心组件，但其制造过程极度消耗先进DRAM晶圆产能。报告提供的数据显示，到2028年，HBM将占据领先级存储器晶圆产能的约34\%，而2023年这一比例仅为6\%。这不仅是量的增长，更是质的挤出——当供应商优先将产能分配给利润更高、需求更确定的AI产品时，留给智能手机、PC、汽车、工业等传统市场的存储器供给，将出现系统性缺口。

报告通过一个“两级DRAM供给瀑布模型”量化了这一缺口：到2027年，PC市场将面临约15\%的存储器短缺，智能手机市场短缺约12\%，分别对应约5800万台PC和1.34亿部智能手机的潜在需求无法被满足。这不是一个预测性的警告，而是一个基于当前产能分配逻辑和建设周期的推演结果。

这意味着什么？存储器已经从一个“可替代的零部件”，变成了AI时代的“新石油”——它的价格、可获得性和分配优先级，将直接决定从云端算力到终端设备的成本结构。

报告揭示了一个正在发生的、但市场可能尚未充分定价的变化：存储器市场的交易机制，正在从“公开市场定价”转向“战略分配”。

报告提出的“买方层级”框架清晰地展示了这一权力格局的转移。位于顶层的超大规模云厂商和AI加速器生态，通过长期协议（LTA）、预付款和战略承诺，锁定了最优先的供给。它们可以确保HBM、服务器DRAM和企业级SSD的供应。第二梯队是服务器OEM和云基础设施公司，它们依赖战略合同获得相对稳定的供给。而第三、第四、第五梯队——从高端PC和手机厂商，到汽车、工业、医疗，再到中小OEM和初创公司——则面临越来越不确定的供给和更高的价格波动。

这种“分配经济”的出现，意味着传统的“供需决定价格”模型正在失效。对于非AI领域的买家，他们面对的不是一个可以随时按需采购的市场，而是一个需要排队、竞价、甚至接受降级规格的“配给制”环境。报告指出，这种变化使得剩余供给池变得更小、更紧张、波动性更大。

这对于下游硬件厂商而言，是一个根本性的成本函数变化。它们不再能简单地通过“涨价”来转嫁成本，因为价格的弹性在消费者端是有限的。更可能的后果是：产品规格下降、升级周期延长、或者利润率被大幅压缩。

3. 芯片通胀正在从“Big Tech”向“Everywhere”蔓延

报告最引人深思的部分，是对“芯片通胀”传导路径的分析。它指出，这种价格上涨的影响，远不止于CPI中的电子元器件分项。

在直接层面，存储器涨价已经反映在电子元器件PPI上，同比上涨约30\%。对于大型云厂商，它们可以通过长期协议锁定成本，并将增加的成本通过云服务价格转嫁给企业客户。但对于非AI领域的买家——比如一家需要采购服务器来运行ERP系统的中型企业，或者一家依赖存储器来制

[... middle omitted ...]

”的问题，而是一个“通胀将以何种形式被吸收”的问题——是消费者买单、企业消化，还是通过产品降级和需求破坏来出清。

市场或许会期待供给端的快速扩张来平抑价格。但报告清晰地展示了从“下单”到“产出”的时间线：从AI需求爆发、到签订长期协议、到订购EUV光刻机、到设备交付安装、到客户认证、到良率爬坡、再到合格产品产出——整个过程需要1到2年甚至更长的时间。

更关键的是，即使产能最终落地，其分配方向也大概率会优先满足AI需求。报告提到，即使到2027年全球DRAM晶圆总产能扩张约30\%，智能手机和PC的存储器供给仍将出现短缺。这意味着，非AI领域的“芯片通胀”，可能是一个持续数年的结构性现象。

政策层面，无论是美国的补贴和税收优惠，还是中国的产能扩张，都无法在短期内改变这一格局。报告假设美国政策保持限制性，而中国产能的增长虽然显著（到2028年占全球净新增晶圆产能约30\%），但短期内不足以缓解全球性的供需错配。政策工具更像是“慢变量”，无法对冲“快变量”驱动的价格冲击。

\subsection{[R012] 全球汽车市场正在经历的不是周期波动，而是竞争逻辑的根本切换}
{\small\color{refgray}归类：未被主题摘要引用}

这份来自某外资投行全球汽车研究团队的周度报告，覆盖了从北美到中国、从特斯拉到法拉利的广泛议题。表面上看，这是一份常规的行业跟踪更新：4月全球汽车销量同比微增0.7\%，美国和中国仍在下跌，欧洲和新兴市场成为主要支撑。但真正值得关注的不是这些数字本身，而是数字背后的结构性信号——汽车行业的竞争焦点，正在从“谁能造出更好的车”转向“谁能在执行层面把规模、技术和成本控制转化为可持续的利润”。

这不是一个周期性的底部判断。这是一个关于行业竞争逻辑已经发生不可逆变化的判断。理解这一点，比预测下一个季度的销量拐点重要得多。

1. 全球销量数据揭示的不是需求疲软，而是赢家与输家的加速分化

4月全球主要市场合计销售590万辆，同比仅增长0.7\%。但这组平淡的总量数据掩盖了两个关键分化。

第一个分化是区域性的。欧洲整体增长7\%，其中英国增长24\%、意大利增长12\%，而美国下滑5.5\%、中国下滑4.3\%。新兴市场表现更为突出：印度增长24.6\%、巴西增长20.5\%、印度尼西亚增长55\%。这不是一个“全球需求疲软”的故事，而是一个“成熟市场结构性见顶、新兴市场仍有渗透空间”的故事。对于全球布局的OEM而言，这意味着区域配置的权重正在发生实质性变化。那些过度依赖单一市场的企业，风险敞口正在放大。

第二个分化是品牌层面的。在主要OEM中，特斯拉4月销量同比增长26.8\%，市场份额提升0.5个百分点。与此同时，大众集团下滑10.1\%、福特下滑13.8\%、通用汽车下滑15.7\%。更值得关注的是，报告明确指出“高端市场在全球范围内正在失去份额，而大众市场正在增长”。高端品牌整体销量同比下降2.3\%，大众市场增长1.3\%。这一趋势如果持续，将对宝马、奔驰、保时捷等品牌的定价能力和利润率构成结构性压力。

这些数据合在一起意味着：全球汽车市场不是一个同质化的市场。不同区域、不同价格带、不同品牌正在经历完全不同的周期。用单一宏观指标来判断行业走向，已经不足以支撑投资决策。

报告中最值得深读的部分之一，是某外资投行团队对特斯拉奥斯汀Robotaxi车队的跟踪分析。他们构建了一个基于NHTSA数据的追踪框架，用以评估自动驾驶安全性的边际改善。

这里的关键洞察是：特斯拉的自动驾驶不再是一个“未来可能实现”的故事，而是一个“每季度都在发生可量化的进步”的故事。报告提到“安全性显示出改善迹象，FSD使用率正在提升，Robotaxi车队正在扩大”。这些表述背后的含义是，特斯拉正在从前期的概念验证阶段，进入一个数据驱动的迭代阶段。

但这里有一个尚未完全回答的问题：安全性的改善速度是否足以支撑监管放行和商业化的时间表？报告构建的追踪框架本身就是一个信号——市场对自动驾驶的评估正在从“相信与否”转向“测量多少”。对于特斯拉而言，这既是机会也是风险。机会在于，如果数据持续向好，市场将重新定价其自动驾驶资产的价值。风险在于，如果改善速度低于预期，目前已经反映在股价中的溢价将面临重新评估。

对于其他布局自动驾驶的企业而言，特斯拉的进展意味着竞争门槛正在提高。如果特斯拉能够通过大规模车队运营积累真实世界数据，其算法迭代速度将形成难以追赶的飞轮效应。

在整车厂面临需求不确定性和价格竞争压力的背景下，某外资投行的供应商季度记分卡给出了一个值得关注的结论：“我们对供应商的看法变得更加积极，因为尽管产量疲软、内存/大宗商品通胀和宏观不确定性，执行、成本控制和产品启动活动仍然支撑了业绩。”

这组判断的背后逻辑是：汽车供应链的价值分配正在发生微妙变化。过去几年，市场习惯于将关注焦点放在整车厂的品牌力和终端需求上。但当前环境下，当需求增长放缓、价格竞争加剧时，真正考验企业的是成本控制能力和运营执行力。而供应商，尤其是那些在特定细分领域具有技术壁垒和规模优势的供应商，反

[... middle omitted ...]

下并不算特别糟糕。但真正值得关注的是中国OEM的海外表现。报告指出，比亚迪4月海外注册量环比下降10-15\%，主要是英国销量从3月的1.5万辆回落至0.5万辆。但同比仍增长55-60\%。吉利海外注册量环比增长10-15\%，主要受巴西和澳大利亚市场驱动。

这些数据揭示了一个正在加速的趋势：中国OEM正在从“国内价格战”转向“全球渗透战”。比亚迪和吉利在海外市场的增长，不是短期冲刺，而是系统性布局的结果。对于全球传统OEM而言，这意味着竞争压力不仅来自中国市场内部，还来自中国品牌在海外市场的扩张。

但报告也揭示了一个尚未完全回答的问题：中国OEM的海外扩张能否持续？比亚迪在英国市场的月度波动表明，海外市场的需求并非线性增长。物流、渠道、品牌认知和售后服务的建设都需要时间。更重要的是，地缘政治风险——关税、补贴政策、监管壁垒——可能随时改变海外市场的竞争格局。这些不确定性意味着，中国OEM的全球化故事仍处于早期阶段，其投资逻辑需要区分“短期催化剂”和“长期结构性优势”。

\subsection{[R013] 铜价突破的关键不在于需求爆发，而在于供给侧的三个结构性断裂}
{\small\color{refgray}归类：能源与大宗商品：定价权从宏观叙事向微观供给切换}

过去六个月，铜价在每吨12,000至13,000美元之间反复拉锯。市场上最主流的两派叙事——一派押注AI与能源转型拉动结构性需求，另一派担忧全球增长放缓与地缘风险——已经在这个区间内完成了充分的定价博弈。但某外资投行最新发布的研报提出了一个更值得关注的判断：铜价将在一个月内触及每吨14,500美元，并在一年内站上15,000美元。

这个判断之所以值得认真对待，不是因为该行上调了需求预测，而是因为它系统性地重估了供给侧的三个关键变量。这三个变量过去被市场视为“已知的已知”，但该报告用最新数据表明，它们正在从“已知”滑向“严重低估”。

更重要的是，这份报告揭示了一个反直觉的定价逻辑：铜价接下来的上行动力，将主要来自“供给跟不上需求”这一事实的重新定价，而非需求本身的加速。对于产业决策者和资产配置者而言，这意味着当前的交易区间可能正在成为过去式。

报告最硬核的修正出现在矿山供给端。该行将2026年全球铜矿产量增速从正增长下调至-0.2\%——换句话说，产量几乎没有增长。这个数字本身已经低于市场共识，但真正重要的不是这个数字，而是它背后的结构性原因。

报告明确指出，其上调了2027年和2028年的全球矿山中断率假设至7\%。这意味着更多矿山的投产和爬坡将出现延误、减产或非计划停产。报告特别点名了几个关键矿区——Grasberg、El Teniente、Kamoa、Cobre Panama——这些全球级矿山在 年的实际产出都存在显著的执行风险。

更深层的问题是刚果（金）的供给风险。该地区近年已成为全球铜矿增量供给的核心来源，但报告警告，当前的硫磺市场约束以及更广泛的地缘政治不稳定，正在使这一地区的持续高产路径变得不确定。对于一个全球市场越来越依赖少数高风险司法管辖区的行业来说，这绝不是一个边际问题。

所以呢？矿山供给零增长意味着，2026年全球铜市场将完全依赖废铜回收和库存释放来满足任何新增需求。而报告的下一个判断恰恰指出，废铜供给同样不可靠。

传统铜市场分析中有一个默认假设：当铜价上涨时，废铜回收量会随之增加，从而部分对冲供给缺口。这个假设正在被现实数据挑战。

报告引用了中国2026年1-4月的废铜进口数据：尽管同期铜价同比上涨超过30\%，废铜进口量仅录得温和的同比增长。这意味着废铜的供给弹性远低于历史水平。报告给出了三个解释：能源和燃料成本上升推高了废铜的回收和运输成本；中国的“以旧换新”政策支持力度较去年有所减弱；全球废铜供应链本身已经处于相对枯竭状态——过去几年冶炼厂对原材料的需求持续超过废铜供给能力。

这个判断的产业含义非常直接：如果废铜不能像过去那样在高价环境下快速补充市场，那么任何需求的边际增长都将直接反映在精铜库存的去化和价格的上涨上。报告预计，2026年废铜和矿山供给的合计增量将无法覆盖AI和能源转型带来的需求增长，更不用说任何可能的周期性需求复苏。

所以呢？铜市场的供给安全边际正在从“有弹性”滑向“刚性”。这是一个需要被重新定价的基本面变化。

这是报告中最具洞察力的部分之一。市场普遍关注美国是否会对精炼铜加征关税，以及关税税率是多少。但报告认为，更重要的变量是美国政策制定者的行为模式。

该行的核心判断是：美国政府最终不会对精炼铜加征关税，但在2026年6月底的审查截止日期之前，也不会明确排除这一可能性。这种“战略模糊”的目的非常明确——通过悬而未决的关税威胁，激励市场参与者持续在美国境内囤积超额铜库存。只要关税风险没有被消除，套利交易就会推动铜流入美国，支撑COMEX溢价和全球铜价。

从数据上看，这一策略已经奏效。报告显示，美国精炼铜进口量在2026年初显著高于正常消费需求水平，COMEX-LME价差在审查截止日期前重新走阔。基金在COMEX上的空头持仓持续低迷，反映出市场对关税

[... middle omitted ...]

纳入其情景分析，并给出了一个明确的判断：即使海峡关闭持续到7月，全球增长和风险情绪在短期内仍可能保持韧性。但如果海峡关闭演变为持续且不稳定的中东局势，铜价以及更广泛的风险资产将面临显著的尾部风险。

这个判断的微妙之处在于：报告既不认为海峡关闭会立即引发系统性风险，也不认为其影响可以忽视。它更像是一个“已知的未知”——市场已经定价了部分风险，但尚未充分定价“重新开放”这一上行情景。

报告的基本假设是海峡将在2026年第三季度重新开放。如果这一假设兑现，铜价将获得一个额外的上行催化。而如果海峡关闭持续超预期，报告在熊市情景中已经为铜价预留了每吨11,000美元的下行空间。

所以呢？霍尔木兹海峡是铜价短期走势中最大的不确定性来源。它既可能成为推动铜价突破15,000美元的催化剂，也可能成为拖累铜价跌回12,000美元的拖累。这个二元性意味着，任何押注铜价方向的投资者都需要对这一变量保持高度敏感。

5. 2027年将出现36万吨缺口，但库存结构使这一数字比看起来更严重

\subsection{[R014] 新能源渗透率突破63\%的真正含义：市场正在重新定义赢家的门槛}
{\small\color{refgray}归类：中国新能源车：渗透率高位下的订单分化与出口爆发}

新能源渗透率突破63\%的真正含义：市场正在重新定义赢家的门槛

5月最后一周，中国新能源乘用车单周订单同比、环比双双增长超过10个百分点，渗透率在5月1日至24日期间达到62.5\%至63.6\%。这些数字本身并不令人意外——市场已经习惯了这个渗透率区间。

但这份某外资投行最新发布的周度报告揭示了一个更值得关注的结构性变化：在整体订单温和增长的表象下，品牌间的分化正在急剧加速。5月最后一周，蔚来订单环比暴增228\%，鸿蒙智行增长92\%，比亚迪增长15\%。而同一周，理想汽车订单环比下滑25\%，小鹏下滑61\%。

这不是一次普通的周度波动。这是一场关于“谁能在高渗透率时代继续增长”的压力测试。

1. 新产品周期的启动正在重塑订单结构，但持续性才是真正的考验

5月最后一周的订单数据最突出的特征是“新品驱动”效应。蔚来ES9、问界M9改款、比亚迪元PLUS改款和宋Ultra DM-i的集中上市，直接推动了这几个品牌订单的环比大幅增长。

从数据上看，蔚来5月订单达到86183辆，环比增长96\%，同比增长190\%；鸿蒙智行5月订单143479辆，环比增长17\%，同比增长154\%。这些数字在绝对量和增速上都相当可观。

但这里有一个需要冷静看待的问题：新产品上市带来的脉冲式增长，能否转化为可持续的订单曲线？从历史数据来看，大多数品牌在新品上市后的第3到4周，订单增速会出现明显回落。5月最后一周蔚来228\%的环比增长，其基数效应（前一周仅11830辆）放大了这一数字的冲击力。

对投资者和产业决策者而言，关键不是盯着单周爆发的幅度，而是观察未来4到6周订单曲线的斜率变化。真正的产品力，体现在上市一个月后能否维持环比正增长。

在众多品牌中，比亚迪的表现值得单独拆解。5月订单197343辆，虽然同比下降41\%，但环比仅下滑12\%。考虑到4月基数较高（223899辆），这个降幅在可接受范围内。

更值得关注的信号来自经销商折扣端。截至5月30日，比亚迪平均经销商折扣为5.14\%，与5月23日持平。而去年同期，这个数字是7.72\%。

这意味着什么？比亚迪正在经历一个“以价换量”到“以产品力换量”的转折。在折扣率没有扩大的情况下，订单量维持在高位。这背后可能有两个因素在起作用：一是改款车型的产品力确实支撑了定价；二是比亚迪的规模效应使其在成本端有更大的腾挪空间，不需要通过大幅降价来刺激需求。

这对整个行业的含义是：价格战的烈度可能正在见顶。当行业龙头开始收窄折扣，其他品牌跟进降价的空间就会被压缩。如果这一趋势延续，整个新能源车行业的利润拐点可能会比市场预期的更早到来。

3. 经销商折扣收窄正在改变竞争格局，但ICE品牌的困境尚未见底

报告数据显示，截至5月30日，新能源车平均经销商折扣为7.75\%，较一周前的7.79\%收窄；传统燃油车平均折扣为19.47\%，同样较一周前的19.67\%收窄。

新能源车折扣收窄，反映的是需求端改善和供给端新品周期的共同作用。当消费者对新产品的接受度提高，经销商就不需要通过大幅让利来清库存。这是一个正向循环的信号。

传统燃油车折扣收窄，则更多是供给端主动收缩的结果。在新能源渗透率突破63\%的背景下，燃油车经销商已经大幅减少了库存压力。但19.47\%的折扣率仍然处于历史高位，说明燃油车的需求端压力远未解除。更关键的是，随着新能源车继续向20万元以下市场渗透，燃油车的价格底线可能还会被进一步击穿。

对投资者而言，燃油车折扣率能否在年底前回到15\%以内，是判断传统车企转型成败的一个重要观察指标。

4. 碳酸锂价格持续下行正在改变产业链利润分配，但电池环节的议价权才是核心

报告显示，电池级碳酸锂价格降至17.95万元/吨，环比下跌0.8\%。与此同时，方形磷酸铁锂电芯和方形三元电芯价格保持稳定。

[... middle omitted ...]

术锁定，正在构建对下游整车厂的议价优势。这个优势能否在碳酸锂价格下行周期中保持，将决定产业链利润的最终分配格局。

5. 报告尚未完全回答的关键问题：高渗透率下的增长天花板在哪里

渗透率突破63\%之后，市场面临一个根本性的问题：剩余不到40\%的市场空间，谁能吃掉，以什么成本吃掉？

这个问题在报告中有所触及，但没有给出明确的答案。从数据上看，5月1日至24日，乘用车零售总量同比下降24\%，而新能源车零售量同比下降11\%。这意味着整体乘用车市场在收缩，新能源车虽然在侵蚀燃油车的份额，但自身的绝对量增长也在放缓。

如果整体乘用车市场继续收缩，那么新能源车即使渗透率达到70\%、80\%，其绝对销量增长也可能进入平台期。届时，行业竞争将从“抢份额”转向“抢利润”。

报告中提到的几个即将发布的新品——蔚来乐道L60改款、比亚迪大唐、零跑D99——将是检验这个判断的关键窗口。如果这些新品能够带动整体市场增量，而不是仅仅在存量市场中互相替代，那么行业的增长故事还可以继续讲下去。

\subsection{[R015] 新能源渗透率突破63\%后，市场真正低估的是新车型的订单转化效率}
{\small\color{refgray}归类：中国新能源车：渗透率高位下的订单分化与出口爆发}

新能源渗透率突破63\%后，市场真正低估的是新车型的订单转化效率

2026年5月最后一周，中国新能源乘用车市场交出了一份看似平淡、实则暗藏分化的成绩单。某外资投行最新发布的周度研报显示，当周主要新能源品牌合计订单环比增长11\%，同比增长20\%，新能源渗透率在5月1-24日期间达到62.5\%-63.6\%。渗透率数字本身并不令人意外——市场早已接受50\%以上的新常态。但真正值得关注的是订单增长的结构：它不是普惠式的行业回暖，而是由特定品牌的新车型集中投放所驱动。蔚来周订单环比暴增228\%，鸿蒙智行增长92\%，比亚迪增长15\%。这三家合计贡献了当周绝大部分增量。

这意味着什么？当行业渗透率进入60\%以上的高位区间，增长的动力已经从“电动化替代”切换为“产品周期竞争”。每一款新车型的订单转化效率，正在成为决定品牌季度排名的关键变量。市场对此的定价可能仍然不足。

以下，我们从订单结构、折扣信号、定价权分化三个层次，逐一拆解这份研报的核心洞察。

5月最后一周的订单数据，表面上是整体向好，但细看品牌间的差距，几乎是断层式的。蔚来周订单从1.18万辆跃升至3.88万辆，环比增长228\%；鸿蒙智行从2.11万辆增至4.07万辆，环比增长92\%。而同期，理想汽车订单环比下降25\%，小鹏下降61\%，吉利银河与极氪合计下降39\%。

这种分化并非偶然。研报明确指出，增长主要由新车型驱动：蔚来ES9、问界M9改款、比亚迪元Plus改款和宋Ultra DM-i。换句话说，当周订单增量几乎全部来自“新车效应”，而非存量需求的自然回暖。

对于产业决策者而言，这一信号的含义是明确的：在渗透率高位，任何品牌的订单增长都越来越依赖产品更新节奏。没有新车型投放的品牌，即使整体市场渗透率在提升，其订单也可能停滞甚至下滑。这意味着行业竞争已经从“谁更早电动化”演变为“谁的产品迭代更快、更准”。小鹏和理想的订单下滑，恰恰反映了其当前产品线在市场上缺乏新鲜感——尽管两者在技术和品牌层面各有积累。

从投资视角看，市场对新能源车企的估值，可能需要从“总量增长故事”转向“产品周期估值”。一个品牌未来6个月的订单可见度，很大程度上取决于其下一款新车的发布时间和市场预期，而非行业渗透率的线性外推。

2. 经销商折扣同步收窄，但NEV与ICE的定价能力差距在扩大

研报中另一个容易被忽视的细节是终端折扣的变化。截至5月30日，新能源车经销商平均折扣为MSRP的7.75\%，较前一周的7.79\%略有收窄；而燃油车经销商平均折扣为19.47\%，同样较前一周的19.67\%收窄。乍看之下，两者都在改善，但折扣率的绝对值差距——近12个百分点——才是关键。

燃油车经销商以接近20\%的折扣率卖车，意味着其定价权已基本丧失。即使折扣收窄，19.47\%的水平仍然表明燃油车终端价格处于“清仓模式”。相比之下，新能源车7.75\%的折扣虽然也不低，但至少还在产品生命周期内的正常促销区间。

更值得关注的是比亚迪的表现。其平均折扣仅为5.14\%，远低于行业平均，且连续两周持平。这表明比亚迪在拥有规模优势的同时，仍然维持了相对较强的定价能力。对于一家月订单接近20万辆的品牌而言，5\%左右的折扣率是健康的。这反过来也说明，比亚迪的利润压力可能小于市场普遍担忧的程度。

这里有一个尚未被充分讨论的推论：如果新能源车折扣继续收窄，而燃油车折扣维持高位，那么燃油车经销商的现金流压力将进一步加大。这可能导致更多燃油车经销商退出或转投新能源品牌，从而加速渠道端的结构性洗牌。研报没有展开讨论这一二阶影响，但它对传统车企的渠道价值评估可能构成重大变量。

研报中提及，电池级碳酸锂价格已降至每吨17.95万元，环比下跌0.8\%，而磷酸铁锂和三元方形电芯价格环比持平。这一组合值得拆解。

碳酸锂价格持续

[... middle omitted ...]

坚挺，那么电池环节的议价权问题就需要重新评估。

这份周度研报提供了丰富的高频数据，但有一个核心问题它没有直接回答：新车型带来的订单增长，有多少能转化为实际交付？蔚来ES9的周订单暴增228\%，但这款车的产能爬坡速度如何？交付周期是否会影响订单的留存率？鸿蒙智行M9改款的订单中有多少是存量换购，多少是新增用户？

这些问题之所以关键，是因为在新能源渗透率超过60\%的市场中，订单的“质量”比“数量”更重要。如果新车订单主要来自品牌内部老用户的换购，那么它对市场份额的净增量贡献有限。如果订单大量来自价格敏感型用户，那么一旦产能不足导致交付延迟，退单率可能上升。

研报中提到的“YTD订单增速”可以提供一些线索：蔚来YTD订单同比增长69\%，鸿蒙智行增长26\%，特斯拉增长2\%。蔚来的高增速部分受益于去年基数较低，但69\%的增长仍然超出了单纯换购所能解释的范围。然而，要确认这一趋势是否可持续，还需要观察未来4-6周的新增订单走势，以及6月11日乐道L60改款发布后的市场反应。

\subsection{[R016] 现代起亚正在悄悄改写美国汽车市场的竞争规则}
{\small\color{refgray}归类：中国新能源车：渗透率高位下的订单分化与出口爆发}

当市场把注意力集中在特斯拉的价格战和比亚迪的全球扩张时，一份来自某外资投行的最新月度跟踪数据揭示了一个容易被忽视的结构性变化：现代汽车与起亚在美国市场的合计份额在2026年5月达到11.8\%，同比提升0.3个百分点，更重要的是，两者的混合动力车（HEV）市场份额已逼近20\%大关，达到19.9\%，同比大幅跃升8.3个百分点。

这不是一次性的月度波动。从YTD数据看，现代与起亚在美国的HEV销量同比分别增长65.1\%和102.1\%，而同期美国整体HEV市场增速为17.4\%。这意味着，这两家韩国车企正在以超过行业3-5倍的速度抢占HEV赛道，而这一赛道的增速本身已是纯电动车（BEV）的3倍以上。

这份报告最值得关注的判断不是11.8\%这个数字本身，而是它背后揭示的竞争逻辑：在纯电动车需求增速放缓、消费者回归务实选择的背景下，现代起亚正在通过HEV建立起一种“成本可控的电动化过渡策略”，这种策略正在改写传统车企在北美市场的利润结构。

1. HEV正在成为现代起亚在美国的“利润压舱石”，而非过渡产品

2026年5月的数据显示，现代与起亚的HEV销量分别为25,559辆和25,392辆，同比增长高达90.1\%和179.4\%。这一增速远超行业HEV增速的32.5\%，也远超其自身整体销量的个位数增长。更关键的是结构变化：HEV在现代总销量中的占比从去年同期的14.7\%跃升至27.1\%，在起亚中更是从11.5\%飙升至31.5\%。

这意味着什么？HEV不再是“等待纯电到来之前的过渡方案”，而是正在成为这两家公司在美国市场的主力利润来源。从激励支出数据也可以佐证这一判断：现代HEV占比提升的同时，其单车激励支出仅为3,330美元，低于行业平均的3,502美元。这说明HEV的高需求允许企业以更少的折扣完成销售，从而保护了利润率。

对于关注全球汽车产业链的投资者而言，这里有一个需要持续跟踪的变量：当HEV占比超过30\%后，企业的整体利润率曲线会如何变化？报告没有直接回答，但历史经验表明，高附加值动力总成占比的提升通常会带来单车利润的结构性改善。

2. 纯电动车份额正在被HEV“蚕食”，但现代起亚的BEV布局并未停滞

美国纯电动车市场在2026年5月出现了同比19.8\%的下降，行业BEV渗透率从去年同期的7.3\%下滑至5.8\%。这是一个重要的信号：消费者对纯电动的热情正在降温，至少在中短期内如此。

然而，现代起亚的BEV表现却逆势上扬。其BEV合计销量9,305辆，同比增长22.5\%，BEV市场份额从7.0\%提升至10.8\%。这一反差说明，在整体BEV市场收缩的环境下，现代起亚的BEV产品正在赢得更多消费者的选择。

但更值得玩味的是结构：现代起亚的BEV销量增速（22.5\%）远低于HEV增速（126.1\%）。这表明，在资源有限的情况下，这两家公司正在将更多产能和营销资源倾斜到HEV上，而BEV则扮演着“技术储备和品牌形象”的角色。这种“HEV保利润、BEV保未来”的双轨策略，可能是当前最务实的电动化路径。

报告没有完全展开的一个问题是：当HEV占比持续提升后，现代起亚是否有足够的电池产能同时支撑HEV和BEV的双线扩张？这需要在后续的供应链跟踪中寻找答案。

2026年5月，现代与起亚的单车激励支出分别同比上升16.8\%和26.4\%，达到3,330美元和3,440美元，接近行业平均的3,502美元。这与YTD数据形成对比：前五个月现代激励支出同比下降10.9\%，起亚同比上升8.4\%。

这个数据点需要放在竞争格局中理解。5月是北美汽车销售的传统旺季，各家车企都会加大促销力度。更重要的是，现代起亚的市场份额从5月的11.5\%提升至11.8\%，是在行业整体销量同比仅微增0.4\%的背景下实现的。这意味着，为了从竞争对手手中

[... middle omitted ...]

一，新车型的HEV版本能否延续当前的增长势头。如果新车型主要推动的是ICE或BEV销量，那么对整体份额的提升效果可能有限。

第二，竞争对手的反应。丰田和本田在HEV领域同样强势，且正在加大北美HEV产能。如果竞争对手也推出更具竞争力的HEV产品，现代起亚的HEV市场份额能否维持20\%的水平？

这两个问题报告没有给出明确答案，但它们决定了12.0\%的预测是保守还是乐观。对于关注这一赛道的读者来说，2026年下半年将是验证这些假设的第一个时间窗口。

5. 一个尚未被充分定价的观察框架：从“市场份额”转向“HEV利润占比”

传统上，汽车股的分析框架主要关注总市场份额和整体销量。但这份报告的数据提示了一个更精细的观察维度：HEV在总销量中的占比及其对利润的贡献。

以现代起亚5月的数据为例，HEV在总销量中的占比已接近30\%，而HEV的利润率通常高于ICE车型，且低于BEV的研发摊销压力。这意味着，HEV占比每提升10个百分点，对整体利润率的拉动可能比市场预期的更大。

\subsection{[R017] 市场真正低估的不是需求，而是二手房定价权的结构性转移}
{\small\color{refgray}归类：地产与消费：二手房定价权转移与消费渠道重置}

过去几周，关于中国房地产市场的讨论重新升温。政策端的积极信号、部分城市销售数据的边际改善，让市场参与者再次开始争论“底部是否已到”。

但这份最新发布的中国房地产周度数据库追踪报告，提供了一组值得深思的读数。它揭示了一个比短期销售波动更关键的结构性事实：市场的定价逻辑正在从新房市场向二手房市场转移，而这一转移的速度，可能比大多数投资者意识到的要快。

报告显示，截至5月31日当周，50城新房周度成交同比增长10\%，相比前一周的-7\%出现明显回升。但与此同时，10城二手房周度成交同比增长9\%，虽然增速较前一周的45\%有所回落，但年初至今累计同比仍保持5\%的正增长。更关键的数据来自去化率：整体去化率从前一周的66\%降至53\%，其中一线城市从66\%降至63\%，二线城市从42\%降至一个更低的水平。

这些数字组合在一起，指向一个被市场普遍低估的结论：新房市场的短期回暖并不稳固，而二手房市场正在成为定价的锚点。真正重要的不是销售量的月度波动，而是二手房价格发现机制的形成速度，以及它对开发商定价权的侵蚀。

1. 新房成交的反弹是基数效应和集中推盘的结果，而非需求复苏的信号

50城新房周度成交同比10\%的增长，乍看之下是一个积极信号。但如果把时间拉长，年初至今累计同比仍为-14\%。这意味着，单周的增长更多是去年同期低基数（-7\%）和开发商在5月底集中推盘的结果，而非购房需求的系统性回升。

更值得关注的是去化率的下降。整体去化率从66\%降至53\%，说明即使推盘量增加，市场吸收新供应的能力并未同步提升。一线城市去化率虽仍维持在63\%的较高水平，但同样出现了边际下滑。二线城市42\%的去化率则进一步确认了库存压力的积累。

对于投资者而言，这意味着：新房市场的短期数据波动不应被过度解读为趋势反转。真正需要观察的是去化率能否在接下来的几周内重新回升至60\%以上。如果去化率持续低于50\%，那么开发商将面临更大的降价压力，而这将进一步侵蚀市场信心。

2. 二手房成交的韧性表明，市场正在经历从“新房定价”到“二手房定价”的切换

与新房市场的波动形成对比，二手房市场呈现出更强的韧性。10城二手房周度成交同比增长9\%，年初至今累计同比保持正增长。虽然增速从45\%回落，但这更可能是高基数效应，而非趋势逆转。

更关键的是，二手房成交的稳定正在改变市场的定价机制。过去几年，中国房地产市场的定价权主要掌握在开发商手中，新房价格是市场的锚点。但随着二手房成交占比持续提升，定价权正在向二手房市场转移。当二手房成交量和价格成为市场基准时，开发商的新房定价策略将受到更大的约束。

报告中的中央指数六城二手房挂牌价格追踪指数从18.2\%降至17.5\%，一线城市房产中介指数从55.5降至53.8。这些指标的下降，反映了二手房卖方的议价能力正在减弱。当卖方开始接受更低的价格成交时，新房市场的价格天花板也随之被压低。

这意味着：投资者需要重新审视对开发商的估值框架。过去基于“新房价格刚性”的假设正在被打破。在二手房定价主导的市场中，开发商的毛利率和资产周转率都将面临新的压力。

3. 三线城市新房成交的“异常”增长，可能掩盖了更深的流动性问题

报告中最引人注目的数据之一，是三线城市新房成交同比增长34\%，而前一周为-12\%。这一巨大的环比改善，很容易被解读为政策刺激在低线城市见效。

但这一判断需要谨慎。三线城市的新房市场体量较小，单周数据容易受到个别项目的集中网签影响。34\%的增长是否可持续，仍需要未来几周的数据验证。更重要的是，三线城市的二手房市场活跃度远低于一二线城市，这意味着其定价机制仍高度依赖新房市场。如果三线城市的新房成交增长只是昙花一现，那么这些城市的库存压力将进一步恶化。

对于布局三线城市的开发商而言，这一数据不应被解读为乐观信号。真

[... middle omitted ...]

售价。当去化率在50\%-70\%之间时，市场处于平衡状态，开发商需要根据竞争态势灵活定价。但当去化率低于50\%时，市场进入买方市场，开发商被迫降价以去库存。

报告显示，二线城市的去化率仅为42\%。这意味着，在二线城市，超过一半的新推房源无法在短期内被市场吸收。如果这一趋势持续，开发商的资金回笼速度将显著放缓，进而影响其拿地和投资能力。

这里有一个尚未被充分讨论的二阶效应：去化率下降不仅影响开发商的当期现金流，还会影响其再融资能力。金融机构在评估开发贷时，越来越关注项目层面的去化速度。去化率持续低于50\%的项目，可能面临银行收紧信贷的风险。这将在供给端形成负反馈循环：销售放缓导致融资困难，融资困难又迫使开发商降价促销，进一步压低市场预期。

5. 报告尚未完全回答的关键问题：二手房挂牌量激增是否已经形成“踩踏效应”

这份报告提供了周度成交和去化率的数据，但有一个关键变量没有被完全覆盖：二手房的挂牌量变化。报告显示二手房挂牌价格指数在下降，但挂牌量的变化趋势同样重要。

\subsection{[R018] 市场真正低估的不是需求，而是供给侧的再定价}
{\small\color{refgray}归类：中国平价零售与自动化：供给侧主动再定价}

最近一周，投行研报覆盖的两家中国价值零售龙头——鸣鸣集团和万辰食品，股价分别反弹9\%和6\%，部分收复了此前两周15\%和10\%的跌幅。表面看，这是市场对“补贴战扩大”担忧的短期缓解。但这份研报传递的核心信号远比“情绪修复”更值得深究：中国平价零售赛道正在经历一次供给侧的主动再定价，而市场对这一变化的系统性含义仍然定价不足。

5月27日至29日的抛售，导火索是市场担心两家公司加大了对加盟商的补贴力度，尤其是针对新店开业。补贴扩大通常被视为行业竞争加剧、单位经济模型恶化的信号。但6月2日收盘后，鸣鸣和万辰几乎同时通过官方渠道发布声明，正式重申对“健康发展”和“理性竞争”的承诺。两家公司明确反对非理性竞争，同时表达了对餐饮价值零售商在零食品类和饮料市场长期增长空间的信心，并强调有序、健康的竞争环境——这与全国性的“反内卷”政策方向高度一致。

这份研报的价值不在于复述这些公开声明，而在于它提供了三个层次的穿透性分析：第一，补贴扩大的真实边界在哪里；第二，门店扩张速度与单店GMV趋势之间的动态平衡；第三，这些变化如何重塑两家公司的估值逻辑。

以下是我们从这份研报中提炼的四个核心洞察，以及一个尚未被完全回答的关键问题。

1. 补贴扩大的本质不是价格战，而是对加盟商进行“供给侧筛选”

市场对补贴的第一反应往往是“利润受损”。但这份研报通过渠道调研指出，本轮补贴扩大的核心目的并非抢占市场份额，而是加速优质加盟商的招募与筛选。报告明确写道：“渠道调研显示，本轮补贴增加仍然是有纪律的。”这意味着补贴不是无差别的“烧钱”，而是有针对性的激励。

关键区别在于：无差别补贴会导致存量加盟商补贴依赖、新店质量参差不齐；而有纪律的补贴则是在用短期财务成本换取长期网络质量。鸣鸣和万辰的声明中反复强调“理性竞争”和“健康竞争环境”，这在政策层面呼应了反内卷方向，在商业层面则意味着两家公司正在主动定义“什么样的竞争是可接受的”。

对于产业决策者而言，这意味着：平价零售赛道的竞争正在从“谁跑得快”转向“谁选得准”。 补贴不再是规模竞赛的工具，而是筛选机制——那些能够通过补贴吸引到高质量加盟商的企业，将在下一阶段获得更强的网络效应和品牌壁垒。而市场目前对“补贴”一词的负面定价，可能忽视了这一结构性变化。

2. 5月开店加速不是偶然，而是2Q25系统性的供给扩张信号

研报中的月度门店追踪数据揭示了比市场预期更强劲的扩张节奏。鸣鸣集团在2026年5月净增门店907家，环比4月的500家增长81\%；万辰食品（好想来品牌）更是在5月净增1,051家，环比4月的600家增长75\%。将时间拉长看，两家公司在2026年前五个月累计新开门店均超过3,000家。

这不是孤立的月度波动。从数据趋势看，鸣鸣在2025年下半年至2026年初的月均净增门店维持在500-900家区间，而进入2026年4月后明显提速。万辰的扩张节奏更具爆发力：2026年5月净增1,051家，是2025年12月以来单月最高，且远超2025年下半年月均600-700家的水平。

这些数字背后是两层含义。第一层，供给侧的加速扩张说明两家公司对市场容量的判断仍然乐观。报告提到“市场仍然足够大，领先企业在良性竞争背景下可以继续扩张”。第二层，也是更重要的：门店扩张的加速正在改变整个赛道的规模经济曲线。 当门店密度达到一定阈值，供应链效率、物流成本分摊、品牌认知度都会出现非线性提升。这意味着，当前的开店提速可能正在为未来的单位经济改善埋下伏笔。

对于投资者而言，需要关注的不是开店速度本身，而是开店速度与单店GMV之间的交叉验证。如果开店加速伴随着单店GMV企稳或改善，那么规模效应正在兑现；反之，则可能是供给过剩的信号。而这份研报恰好提供了这一交叉验证的关键数据。

3. 单店GMV的拐点信

[... middle omitted ...]

支撑因素之一。对于价值零售商而言，品类扩张意味着客单价提升、复购率增加和门店坪效改善——这些都是单位经济模型的长期驱动因素。

但这里存在一个微妙但关键的问题：单店GMV的改善是否可持续？尤其是当新店加速扩张时，新店通常在开业初期会有较低的GMV贡献，这会拉低整体同店数据。万辰5月“包含新店在内的日均GMV转正”是一个积极信号，说明新店的质量正在提升。但这一趋势能否在更大样本下持续，仍需后续数据验证。

4. 估值折价正在收窄，但风险溢价尚未完全消化规模效应的二阶影响

当前鸣鸣和万辰分别交易在2026年预期市盈率的17倍和14倍。这一估值水平相对于全球价值零售同行（如美国的Dollar General、Dollar Tree，日本的良品计画、Pan Pacific International Holdings）仍有一定折价。报告给出的目标价暗示鸣鸣有约52\%的上行空间（基于23倍2027年预期市盈率），万辰有约67\%的上行空间（基于21.5倍2027年预期市盈率）。

\subsection{[R019] 市场真正低估的不是需求，而是供给侧的再定价能力}
{\small\color{refgray}归类：能源与大宗商品：定价权从宏观叙事向微观供给切换}

锂价在突破20万元/吨后短暂回落，市场随即出现分歧。一部分投资者担忧高价抑制终端需求，另一部分则怀疑涨势能否持续。但某外资投行在6月初发布的这份研报给出了一个更值得关注的判断：锂价年内有望测试25万元/吨，而权益市场尚未充分定价这一情景。

这不是一个简单的看多观点。它的底层逻辑不同于 年那轮由补贴和囤货驱动的暴涨。报告的核心主张是：当前锂价上行的驱动力，已经从“故事驱动”切换为“基本面驱动”，而基本面中最被低估的不是需求增速本身，而是供给侧的间歇性断裂与再定价能力。

这份报告基于对电池生产管道、库存动态和供给中断事件的月度追踪，给出了一个可以持续验证的框架。对于任何关注新能源产业链的决策者来说，理解这个框架，比猜测下一个月的锂价数字更有价值。

市场对2026年新能源车和储能需求的担忧，主要集中在两个点上：一是新能源车销量可能不及预期，二是原材料成本上涨可能侵蚀储能的内部收益率，从而抑制装机。但报告引用的第三方咨询数据显示，这些担忧尚未在实物层面兑现。

2026年前四个月，中国电池产量同比增长41\%至874 GWh。其中，动力电池产量增长22\%，储能电池产量增长99\%。更关键的是，这一增长趋势并非脉冲式。从月度数据看，3月和4月的产量分别达到225 GWh和247 GWh，环比持续加速。这意味着，下游电池厂并没有因为锂价上涨而放缓生产节奏。

为什么需求如此坚韧？报告给出了两个容易被忽略的结构性原因。第一，新能源车出口超预期，拉动了国内电池出货。第二，商用车电动化加速，而商用车的单车电池装机量远高于乘用车。2026年一季度，商用车电池装机量同比增长109\%至39 GWh，占全部电池装机的比例从2024年的9\%跃升至22\%。这个占比的提升，意味着每卖出一辆电动商用车，对锂的消耗量是乘用车的数倍。

对投资者而言，这意味着：即使乘用车增速放缓，商用车和储能这两个结构性增量，已经足以支撑锂需求维持较高增速。市场对“需求见顶”的叙事，可能需要重新校准。

2. 库存已降至三年低点，但市场低估了“低库存+高需求”组合的定价含义

报告跟踪的锂碳酸盐库存数据，揭示了一个比需求更紧迫的信号。截至2026年4月，中国锂碳酸盐总库存已从2025年中的约14.5万吨降至约11.5万吨。更关键的是库存结构的变化：冶炼厂库存从6.5万吨骤降至2.5万吨，而下游库存维持在4.5万吨左右。

这意味着什么？冶炼厂库存的快速消耗，表明上游供给正在被下游需求持续消化。而下游库存并未显著累积，说明电池厂并没有在囤货，而是在按需采购。这种“低库存+高需求”的组合，历史上往往预示着价格弹性会超预期——任何供给冲击都会被放大。

报告特别指出，锂价在突破20万元/吨后出现回调，但这更像是技术性修正而非趋势反转。从库存动态看，基本面并未恶化。相反，随着8-9月传统旺季的到来，电池厂为完成年度产量目标，将迎来一轮集中的原料采购需求。届时，如果冶炼厂库存仍处于低位，锂价的第二波上涨可能会比第一波更猛烈。

这个判断的关键假设是：供给端不会在旺季前出现大规模复产。而现实恰恰是，供给端的不确定性正在增加。

3. 供给侧的“间歇性断裂”正在重塑定价权，但市场仍未充分消化

如果说需求是锂价上涨的燃料，那么供给中断就是点燃燃料的火花。报告列举了多个供给侧的扰动因素，其中三个值得重点关注。

第一个是JXW矿山的停产。这座矿山自2025年8月起因采矿许可证到期和环保问题停产。虽然宁德时代在2026年初获得了新的采矿权，但环保许可证仍在审批中。报告指出，复产时间表取决于政府流程而非企业意愿，因此市场难以形成一致预期。这意味着，JXW的停产将作为一个持续的供给减量，对锂价形成支撑，直到复产被明确批准。

第二个是江西锂云母矿的潜在风险。报告提到，除了JXW，还有

[... middle omitted ...]

划，但它们的增量主要集中在2027年，2026年的实际贡献有限。

把这些碎片拼在一起，可以得出一个清晰的判断：2026年的锂供给是“净减量”的。中国国内的高成本矿山在减产，海外矿山的复产需要时间，而需求仍在增长。这种供需错配的持续时间，取决于政府审批速度和海外项目进展，这两者都存在很大的不确定性。

报告没有完全展开的一个问题是：如果锂价涨到25万元/吨，澳洲矿山是否会加速复产，从而提前结束这个上涨周期？这个问题的答案，将决定这轮锂价上涨是“脉冲式”还是“可持续的”。

市场普遍认为，锂价上涨会挤压中游材料厂商和电池厂的利润。但报告提出了一个反直觉的观点：更强的锂价反映了更强的需求，而这将正面提振整个电池链的股票情绪。

这个逻辑链条是：锂价上涨不是由供给收缩单独驱动的，而是由需求拉动和供给收缩共同作用的结果。如果需求足够强劲，电池厂可以将成本转嫁给下游车企和储能终端。事实上，报告引用的电池产量数据已经表明，电池厂并没有因为锂价上涨而减产，说明它们对成本转嫁有信心。

\subsection{[R020] 市场低估的不是需求，而是AI基础设施供应链的再定价}
{\small\color{refgray}归类：AI/算力与半导体：从周期品到战略资产的再定价}

过去几个月，关于AI算力需求见顶的讨论不绝于耳。一些投资者开始担忧，大模型训练放缓、推理成本快速下降，是否意味着整个AI基础设施的投资逻辑正在松动。

某外资投行5月发布的AI项目追踪报告提供了截然不同的信号。这份报告追踪了全球范围内neocloud、主权AI和企业级部署的最新动态，其核心判断是：AI基础设施市场正在经历从“建设能力”到“交付价值”的关键转型，而这一转型将重塑整个供应链的价值分配逻辑。

真正被市场低估的，不是AI需求本身，而是当基础设施从实验性部署走向生产化运营时，哪些环节将获得重新定价的机会。

1. Neocloud的战略转向正在改变GPU资产的定价逻辑

5月最值得关注的两笔交易，都不是简单的产能扩张。IREN收购Mirantis、Nebius收购Eigen AI，这两笔交易的共同特征是：neocloud正在从“提供GPU算力”向“提供生产级推理服务”转型。

IREN以约6.25亿美元收购Mirantis，后者是云基础设施和Kubernetes编排领域的专业服务商。这笔收购的直接效果是让IREN能够更快地部署和运维其GPU基础设施，提升运营可见性和技术服务能力。但更深层的含义在于，IREN正在为自己的GPU资产增加一层“软件溢价”——不再是简单地出租算力，而是提供一套可管理、可监控、可优化的AI云交付平台。

Nebius以约6.43亿美元收购Eigen AI的逻辑更为清晰。Eigen AI是推理和模型优化领域的领先公司，其技术栈包括Sparse Attention和Activation-aware Weight Quantization等先进优化技术。Nebius将这些技术整合进其“Token Factory”平台，目标是为企业级AI工作负载提供自动缩放端点和微调流水线。

这两笔交易的共同启示是：GPU作为硬件的价值正在被标准化，而围绕GPU构建的编排、优化、推理服务层正在成为差异化竞争的关键。当越来越多的企业开始部署AI应用，真正稀缺的不是算力卡本身，而是让算力高效运转的软件能力。

这对投资者的含义是明确的：未来neocloud的价值评估，不应仅看其拥有的GPU数量或功耗规模，而应更多地评估其软件栈的深度和客户粘性。那些能够将硬件规模转化为软件议价权的公司，将获得更高的估值倍数。

2. 企业部署AI的底层逻辑正在从“训练”转向“推理”，但市场尚未充分定价

5月的另一个关键信号来自Dell Technologies World上的一系列企业案例。三星、马自达、Eli Lilly、Hudson River Trading——这些不同行业的企业都在部署AI基础设施，但它们的部署模式有一个共同特征：强调本地化和混合部署。

Eli Lilly的LillyPod超级计算机包含1,016块Blackwell Ultra GPU，但更重要的是，Dell为其提供了覆盖制造工厂的本地计算、存储和备份方案，用于数字孪生和仿真。马自达计划用Dell AI数据平台统一其模型开发和CAD存储环境，最终构建支持AI和生成式AI工作负载的数据湖。Hudson River Trading在挪威建设了专用AI研究数据中心，采用了AMD处理器和HGX B200系统的组合。

这些案例揭示了一个被市场低估的趋势：企业AI部署正在从“去云上训练模型”转向“在本地运行推理工作负载”。训练阶段需要大规模、集中化的算力集群，而推理阶段则对延迟、数据主权和成本控制有更高的要求。这意味着，企业不会把所有AI工作负载都迁移到公有云，而是会形成“训练在云端、推理在本地”或“核心训练在云端、边缘推理在本地”的混合架构。

这一趋势的直接受益者不是云服务商，而是能够提供本地化AI基础设施的硬件厂商和系统集成商。Dell推出的Pow

[... middle omitted ...]

ra超级计算机达到全面系统验收。

主权AI的逻辑是：各国政府出于数据主权、国家安全和技术自主的考虑，正在推动在本国境内建设AI算力基础设施。这一趋势在澳大利亚、欧洲和东南亚尤为明显。

但这里有一个尚未被充分回答的问题：主权AI的需求是否可持续？目前我们看到的多是“建设型”需求——政府或国有机构出资建设算力中心。但建成之后，这些算力能否被有效利用？如果利用率不足，主权AI可能只是“一次性建设需求”，而非持续运营需求。

另一个值得关注的信号是，主权AI项目往往伴随特定的技术栈要求。Sandia选择了NextSilicon的Maverick-2加速器，而非NVIDIA GPU，这暗示在国家安全领域，对非美系技术栈的需求正在上升。这是否会催生一个独立的“主权AI供应链”？目前尚不明确，但值得持续追踪。

如果说 年是AI模型的“军备竞赛期”，那么 年正在进入“推理部署期”。Akamai与Anthropic签署的18亿美元、为期七年的云基础设施协议，是这一趋势的最新注脚。

\subsection{[R021] ASCO 2026的真正信号：日本药企的差异化正在被重新定价}
{\small\color{refgray}归类：医药与生物科技：ADC与双抗的交叉验证}

ASCO 2026的真正信号：日本药企的差异化正在被重新定价

每年ASCO都会带来一批令人振奋的临床数据，但真正值得产业决策者关注的，不是某个药物PFS延长了多少个月，而是这些数据如何改变竞争格局的底层逻辑。今年，某外资投行的日本医药研报揭示了一个被市场低估的结构性变化：日本头部药企的管线价值，正在从“跟随式创新”转向“差异化定价权”的验证期。

这份报告的判断核心并非某个具体药物能否获批，而是：当全球大药企在PD-1、ADC、KRAS等热门靶点上陷入同质化竞争时，日本公司能否通过精准的患者分层和更优的安全性/便利性组合，在细分适应症中建立不可替代的定价地位。

以下是我们从报告数据中提炼出的五个关键洞察，以及一个尚未被充分讨论的潜在风险。

1. ONO-4578的积极数据背后，隐藏着一个关于“患者筛选”的经典困境

Ono Pharmaceutical的ONO-4578在二线胃癌的Phase 2试验中，整体PFS达到了9.0个月对6.9个月的统计学显著改善（HR 0.67）。这看起来是明确的积极信号。但研报在子组分析中指出了关键细节：疗效优势几乎完全集中在PD-L1阳性（CPS≥1）患者群体中，而在CPS Personal reading notes and learning share only. Not investment advice.

最近ASCO2026大会，几家日本药企公布了重要临床数据。我快速梳理了核心看点，帮你省时间。

*1️⃣ 小野药品：ONO-4578 胃癌Ⅱ期数据不错，但别急着乐观**

联合Opdivo+化疗一线胃癌，PFS显著延长（9.0月 vs 6.9月，HR=0.67）。好消息是PD-L1阳性组效果更突出（HR=0.52）。

⚠️ 但亚组分析显示： PD-L1阴性组没有明确获益，甚至呈恶化趋势 Claudin18.2阳性患者获益偏小 研报特别标注：这是探索性分析，样本量小

*2️⃣ 武田：IBI363(TAK-928) 肺癌早期数据亮眼**

与信达合作的双功能药物，一线NSCLC（PD-L1低表达/阴性）确认ORR达81.8\%。无论鳞癌/非鳞癌、PD-L1阴性/低表达，都观察到疗效。

研报观点：与之前摘要信息一致，没有重大变化。会继续关注美国Ⅱ期数据。

*3️⃣ 第一三共：Datroway面临竞品挑战，但差异化仍在**

中国竞品sac-TMT联合K药一线NSCLC数据不错（PFS HR=0.35）。但研报认为Datroway在安全性、便利性上仍有优势。

关键看点： AVANZAR试验（TROP2阳性患者）— 预计2026下半年出数据 TROPION-Lung07/08针对不同PD-L1分层人群

竞品daraxonrasib在胰腺癌后线治疗中OS显著优于化疗（13.2月 vs 6.7月，HR=0.40）。研报表示目前对setidegrasib看法不变，但两药在一线治疗中的对决值得持续跟踪。

📌 一句话总结：日本药企管线有进展，但多数还在早期验证阶段。竞品数据虽亮眼，原研药仍有差异化空间。后续Ⅲ期数据和真实世界表现才是关键。

Japan Healthcare: Pharmaceuticals: ASCO 2026: Focus on Ono/Takeda pipeline drugs, Daiichi Sankyo/Astellas competitors

We highlight the implications for our Japan pharmaceuticals coverage from the 2026 American Society of Clinical Oncology (ASCO) Annual Meeting. Clinic

[... middle omitted ...]

ug. Regarding IBI363 (TAK-928) and sac-TMT, a potential competitor to Datroway, there is no major change to our view from the disclosure in the abstracts (released May 21 US time; see this report for our view). Favorable data were disclosed for daraxonrasib, a potential competitor to setidegrasib. While this does not change our view on setidegrasib at present, we think development trends for both drugs will need to be monitored closely.

\subsection{[R022] EM债券曲线正在传递一个被忽视的信号：期限溢价的结构性分化正在取代方向性交易}
{\small\color{refgray}归类：汇率与新兴市场债券：定价机制的结构性转移 / 风险偏好与资金流向：从beta交易到alpha分化}

EM债券曲线正在传递一个被忽视的信号：期限溢价的结构性分化正在取代方向性交易

过去一个月，新兴市场美元债券收益率曲线几乎没有任何方向性变化。EM整体10s30s利差仅收窄1个基点至68bp，EM投资级收窄1bp至61bp，CEMBI同样收窄1bp至56bp。从任何平均指标看，这都是一个沉闷的市场。

但沉闷的表象之下，个体的分化远比均值剧烈。印尼（INDON）的10s30s曲线一个月内收窄了26bp，而多米尼加（DOMREP）同期陡化了11bp。埃及（EGYPT）曲线斜率高达141bp，是印尼的15倍。这些极端差异不是噪声，它们揭示了一个被均值掩盖的结构性事实：EM债券市场正在从“beta交易”转向“alpha分化”，而大多数投资者仍然在用旧框架衡量风险。

这份来自某外资投行的EM曲线月度报告，提供的不只是一组月度变动数据。它真正值得关注的是：在美债收益率大幅波动（10s30s曲线单月熊陡6bp）的背景下，EM曲线整体保持了惊人的韧性。这背后是信用基本面的分化，而非简单的利率传导。

1. 美债曲线陡化而EM曲线持稳，意味着EM定价逻辑正在从“利率跟随”转向“信用定价”

报告中最引人注目的对比在于：过去一个月，美国国债10s30s曲线熊陡了6bp（长端利率上升更快），而美国高评级公司债曲线同步熊陡了2bp。但EM整体曲线不仅没有跟随，反而牛平了1bp。

这种背离不是偶然的。它表明EM信用资产的定价锚正在从“无风险利率曲线”向“信用基本面曲线”切换。当美债曲线陡化通常反映通胀预期或期限溢价上升时，EM曲线如果同步陡化，意味着市场认为EM信用风险在恶化——这正是2022年大部分时间发生的情况。但当前EM曲线持稳甚至微幅收窄，说明投资者对EM信用风险的评估没有随美债利率上升而恶化。

关键含义：EM曲线现在更多反映的是发行体自身的信用轨迹，而非全球利率环境的被动映射。对于资产配置者而言，这意味着EM曲线交易策略需要从宏观方向判断转向微观信用筛选。

2. 亚洲EMBIG曲线单月收窄7bp，成为全球EM中最“平坦”的板块，背后是信用质量的系统性改善

EMBIG亚洲10s30s曲线从44bp收窄至37bp，是过去一个月EM全球各板块中变动最大的。这并非偶然。报告显示，EMBIG亚洲的1年变动幅度达到-23bp，在所有EM区域中最为显著。同时，亚洲EMBIG的10年期利差仅为99bp，远低于其他区域。

两组数据叠加意味着：亚洲EM发行体的长端信用溢价正在系统性地被压缩。这不是一次性的技术性调整，而是市场对亚洲主权和准主权信用可持续性的重新定价。印尼（INDON）曲线收窄26bp是全球EM中幅度最大的，但这恰恰发生在印尼新10年期和30年期基准债券纳入EMBI指数的时点——新基准的流动性溢价和信用质量共同压低了长端利差。

关键含义：亚洲EM曲线已经进入“低斜率、低利差”的双低区间，继续压缩的空间有限。但对于尚未充分定价的非洲或拉丁美洲发行体，曲线陡峭化可能对应着未被市场充分认知的信用改善机会。

3. 埃及曲线斜率141bp居全球EM之首，但真正需要关注的不是斜率本身，而是其背后的流动性结构

埃及（EGYPT）以141bp的10s30s曲线斜率位居全球EM最陡之首，远高于第二名PEMEX的121bp。直观上，这通常被解读为市场对埃及长端信用的极度悲观——投资者要求极高的期限溢价来补偿长期不确定性。

但报告中的一个细节值得深究：埃及在5月底指数再平衡时，新10年期和30年期基准债券被纳入EMBI指数。这意味着141bp的斜率部分反映的是新旧基准之间的流动性差异和指数再平衡效应，而非纯粹的信用风险定价。同样的情况也出现在印尼，但方向相反——印尼新基准的纳入反而压平了曲线。

关键含义：在指数再平衡窗口期，曲线斜

[... middle omitted ...]

指数中退出），拉美曲线的实际收窄幅度可能要大得多。

这引出一个更普遍的问题：EM曲线平均指标在指数成分变动时可能产生误导。当最平坦（或最陡峭）的发行体因到期、评级变动或指数规则调整而退出时，整体曲线斜率会发生机械式变化，与信用基本面无关。

关键含义：对于依赖EM曲线平均指标进行资产配置的投资者，必须定期核查指数成分变动的影响。报告提供了一组“1年高低”和“1月前”的辅助参考数据，但并未系统性地剔除成分变动效应。这是一个尚未被充分解答的问题：如果剔除指数再平衡影响，EM曲线的真实变动幅度是多少？

5. EM高收益债曲线斜率反而低于投资级，这是一个反直觉但合理的定价现象

报告数据显示，EM高收益（HY）10s30s曲线为82bp，而EM投资级（IG）为61bp。乍看之下，高收益债曲线斜率更高，符合风险溢价逻辑。但进一步拆解后会发现：CEMBI高收益曲线斜率仅为48bp，远低于CEMBI投资级的58bp；EMBIG高收益曲线为88bp，而EMBIG投资级为62bp。

\subsection{[R023] 市场真正低估的不是AI需求，而是存储器供给侧的再定价}
{\small\color{refgray}归类：AI/算力与半导体：从周期品到战略资产的再定价}

过去两年，市场对AI基础设施的投资热情从未冷却，但一个关键问题始终悬而未决：当GPU供给逐步跟上，算力瓶颈是否会从“芯片”转移到“存储”？某外资投行在最新发布的研报中给出了一个清晰的判断——DRAM已经成为AI建设中的首要瓶颈，而这一结构性短缺在2028年前看不到快速解决方案。这意味着，存储器股票当前不到10倍的远期市盈率，不是估值泡沫，而是市场对盈利持续性的定价仍不够充分。

这份报告上调了美光科技的目标价至1050美元，SanDisk至1750美元，并大幅调高了2027年的盈利预测。上调幅度之大——美光2027年EPS预测上调48\%，SanDisk上调24\%——本身就是一个信号：分析师认为市场对存储器公司未来两年盈利能力的理解存在系统性偏差。

本文将从五个维度拆解这份报告的核心逻辑，并指出报告尚未完全回答的关键问题。

1. 这轮存储器短缺不是周期性的，而是由AI建筑的结构性需求驱动

传统上，存储器行业遵循“繁荣-萧条”周期：供不应求时价格上涨，厂商扩产，随后供过于求，价格暴跌。但这一次，驱动因素发生了根本性变化。

报告明确指出，DRAM短缺的核心原因不是消费电子需求回暖，而是AI基础设施建设的持续扩张。超大规模云服务商对HBM（高带宽存储器）的需求几乎无止境，而DRAM的供给增长却受到两个硬约束：一是洁净室建设周期，二是EUV光刻设备的可用性。这两个因素都不是短期能解决的。

报告引用了在台湾的调研反馈：供应链预期DRAM价格在5月季度上涨40\%，8月季度再涨15\%，甚至高于分析师的模型假设。这意味着，即便分析师已经上调了价格预期，真实市场可能比模型更紧。

这里的关键洞察是：存储器短缺已经从“周期性事件”转变为“结构性特征”。对于投资者而言，这意味着不能再用过去10年的估值框架来定价这些公司。过去，存储器公司的高盈利往往被视为不可持续，市场给予极低的远期市盈率。但如果短缺持续2-3年甚至更久，这些盈利的“持续性”就需要被重新评估。

美光的市盈率从3月底的4倍扩张到目前的11倍，但报告认为这远未结束。核心逻辑是：市场正在逐步接受“高盈利将持续更长时间”这一叙事，但定价尚未完全反映。

报告将美光的目标价上调至1050美元，估值基础是基于周期平均EPS的29.5倍。这个倍数本身并不激进——与半导体板块整体估值相当。但关键在于，分析师将周期平均EPS从21美元大幅上调至35美元。这意味着，即便采用与之前相同的估值倍数，目标价也几乎翻倍。

更值得关注的是，报告提到美光预计将在2027财年启动大规模股票回购，模型假设 年回购规模达500亿美元。这一信息有两个含义：第一，公司自身对现金流前景有足够信心；第二，回购将进一步提高每股收益，形成正反馈。

但这里有一个尚未被充分讨论的问题：回购启动的前提是美光不再受芯片法案相关条款的限制。如果政策环境发生变化，这一催化剂可能推迟。报告没有对此展开，但这是一个需要持续跟踪的政策变量。

许多投资者习惯将目光集中在AI需求的增长曲线，但报告提醒我们：供给侧的约束可能才是决定价格走势的更关键变量。

美光的资本支出计划非常激进。报告模型显示，2027财年资本支出将达440亿美元，2028财年400亿美元，远超市场预期的376亿和369亿美元。如此大规模的投资，是否意味着供给即将过剩？

报告的回答是：不会。原因在于，这些资本支出中有相当大一部分被HBM的“产率折损”所消耗。所谓的“产率折损”，是指生产HBM所需的DRAM晶圆数量远高于传统DRAM——因为HBM需要将多个DRAM die堆叠在一起，每一层都需要额外的测试和封装步骤。这意味着，即便资本支出大幅增加，实际可用的位元供给增长仍然有限。

报告预计，位元供给增长将在 年加速，但节奏仍然受制于工厂建设和设备安装的

[... middle omitted ...]


报告明确表达了对美光的相对偏好，核心原因在于两家公司的终端市场结构不同。

美光的DRAM业务中，数据中心和AI相关收入占比更高。而SanDisk的NAND业务虽然同样受益于超大规模云服务商的采购，但其需求结构中PC和智能手机的权重更大。报告提到，Kioxia（SanDisk的合资伙伴）在投资者日上仅将长期位元增长预期从20\%小幅上调至22\%，主要增量来自AI推理，而PC和移动市场增长乏力。

这一差异意味着：如果消费电子需求持续疲软，SanDisk的业绩弹性可能不如美光。但报告同时指出，消费电子疲软并未影响NAND的定价环境——这恰恰说明，超大规模云服务商的需求已经足够强大，足以支撑整个市场。

这里有一个值得深思的观察：报告将长期供应协议（LTA）视为“症状而非原因”。也就是说，超大规模云服务商之所以愿意签署多年供应协议，是因为他们预期供应将持续紧张，而非这些协议本身导致了价格上涨。这一视角的转换很重要：它意味着需求端的承诺是真实且持续的，而非投机性的库存堆积。

\subsection{[R024] 消费支付数据回暖，但市场真正低估的是银行非息收入的修复弹性}
{\small\color{refgray}归类：地产与消费：二手房定价权转移与消费渠道重置}

这份来自某外资投行的研报，在2026年一季度的支付数据中捕捉到了一个被官方零售数据掩盖的转折信号。总系统支付量同比大幅反弹29\%，银行卡消费支付在经历了2025年全年负增长后，首次转正至同比上升3.2\%。这些数字本身并不令人意外，真正值得关注的是，它们正在讲述一个与官方消费数据截然不同的故事。

官方社会消费品零售总额增速在2026年一季度进一步放缓至2.4\%，但支付数据的表现却明显优于这一读数。某外资投行的分析师认为，这种背离可能源于消费从线上向线下渠道的结构性回流。如果这一判断成立，那么市场对银行盈利能力的评估将需要重新校准——尤其是那些长期被低估的非息业务收入。

报告中最具冲击力的数字，是分析师估算的家庭金融资产在2026年一季度实现了10.8\%的同比增长，主要由股票和共同基金驱动。这意味着，尽管宏观叙事仍然充满不确定性，但居民部门的资产负债表正在以比大多数人预期的更快的速度修复。而这一点，恰恰是市场目前定价最不充分的变量。

1. 支付数据的“背离”揭示的不是消费疲弱，而是渠道结构的深刻重置

2026年一季度，官方零售销售增速从2025年全年的水平进一步下滑至2.4\%。如果只看这个数字，很容易得出“消费依然疲弱”的结论。但支付数据给出了完全不同的信号。

银联支付增速反弹至11.2\%，网联支付增速也维持在7.2\%的合理水平。两者均显著高于官方零售增速。更关键的是，总系统支付量同比飙升29\%，对比2025年四季度的仅5\%增速，这是一个数量级的跃升。

这种背离并非统计误差。某外资投行在报告中提出了一个尚未被市场充分讨论的解释：消费正在从线上渠道回流至线下。如果这一判断成立，那么官方零售数据可能系统性低估了真实的消费活跃度，因为线下消费的统计覆盖和时效性往往弱于线上。对于银行而言，线下消费的回流意味着银行卡收单、商户手续费等传统业务的交易量基础正在改善，这直接关系到中间业务收入的修复。

2. 银行卡消费从负增长转正，是银行非息收入最直接的先行指标

报告中最具操作意义的信号，是银行卡消费支付在2026年一季度实现了3.2\%的同比增长。这个数字本身并不惊人，但它终结了2025年全年负增长的局面。某外资投行明确表示，虽然银行在信用卡贷款投放上仍然保守，但这一趋势标志着银行卡消费增长正在回归3\%-4\%的正常节奏。

对于银行股投资者而言，这个信号的意义怎么强调都不过分。过去两年，银行非息收入的最大拖累来自银行卡业务——信用卡分期手续费、商户回佣、年费等收入项全面承压。如果银行卡消费能够稳定在3\%-4\%的同比增长通道，那么银行非息收入的最低点大概率已经过去。

但这里有一个关键假设需要验证：银行是否愿意在信用卡贷款投放上从“保守”转向“适度积极”？报告指出银行目前仍然谨慎，这意味着消费支付量的增长更多来自存量用户的活跃度提升，而非信用扩张。如果后续银行信贷政策出现边际松动，银行卡业务的弹性可能会超出当前市场预期。

3. 家庭金融资产增速10.8\%：一个被忽视的银行收入支撑因子

某外资投行估算，2026年一季度家庭金融资产同比增长10.8\%，主要由权益和共同基金驱动。这个数字的隐含含义是：居民财富效应正在逐步改善，而这将直接转化为银行的财富管理业务收入。

报告没有深入展开的是，家庭金融资产增长对银行收入的影响路径至少有两条。第一，资产管理规模扩张带来管理费收入增长；第二，居民风险偏好回升推动高费率产品（如权益类基金、结构性存款）的销售占比提升。如果一季度家庭金融资产增速能够维持，2026年银行代销基金、理财等业务的收入贡献将显著高于2025年。

4. 报告尚未完全回答的关键问题：支付数据改善的可持续性有多强？

尽管一季度数据令人鼓舞，但报告并未对可持续性给出明确判断。这里至少有两个需要进一步验证的变量

[... middle omitted ...]

映的是消费结构变化，而非总量扩张。这意味着银行非息收入的改善可能更多是结构性的，而非周期性的全面反转。

5. 投资者需要建立的新观察框架：从“宏观消费”转向“支付结构”和“家庭资产负债表”

基于这份报告的启示，评估银行股投资价值时，传统的宏观消费数据可能不再是核心先行指标。以下三个维度值得纳入日常跟踪框架。

第一，支付结构指标。比较银联与网联支付增速的差值，可以判断线下消费回流的速度。如果银联增速持续高于网联，说明消费渠道重置正在深化，对银行卡业务占比较高的银行构成利好。

第二，家庭金融资产增速。这一指标直接关联财富管理业务的收入基础。如果后续季度增速能够维持在10\%左右，银行代销业务的压力将显著缓解。

第三，银行信用卡贷款增速。这是判断银行是否从“保守”转向“积极”的关键观察点。如果信用卡贷款增速开始回升，意味着银行对消费复苏的信心正在增强，非息收入的弹性将得到进一步释放。

6. 一个值得追问的伏笔：如果支付数据持续优于零售数据，市场何时会重新定价银行？

\subsection{[R025] 中国自动化需求的真实底色，被季节性波动掩盖了}
{\small\color{refgray}归类：中国平价零售与自动化：供给侧主动再定价}

理解中国制造业投资，不能只看月度环比。2026年5月的数据提供了一个关键信号：当市场将注意力集中在电子和电池行业的季节性回调时，机床工具、纺织机械和通用机械领域的订单正在悄然走强。这份来自某外资投行的研报，通过对台湾气动元件龙头亚德客（Airtac）最新销售数据的拆解，揭示了一个正在发生的结构性变化——AI数据中心投资正在成为自动化需求的新引擎，而这一轮增长的质量和持续性，可能与过去依赖新能源和消费电子的周期有本质不同。

报告的核心判断是：5月亚德客日均销售额同比增长26\%，环比仅微降1\%，这远好于市场对“五月淡季”的悲观预期。更重要的是，订单价值连续第几个月超过出货价值，意味着积压订单仍在积累。这不是一个需求疲软的信号，而是供给端正在被结构性力量重塑的迹象。

5月亚德客销售额环比下降10\%，表面看符合季节性规律——劳动节长假减少了工作日。但日均销售额仅环比下降1\%，说明需求的内在强度远超表面数据。更值得关注的是分行业表现。

电子行业销售额同比增长12\%，电池行业增长55\%，但两者均较4月有所回落。报告明确指出，这很大程度上归因于季节性因素。真正超出预期的是机床工具行业，销售额同比增长35\%，且较4月的高位继续攀升。纺织机械、包装机械和通用机械的订单也呈现环比改善。

这意味着什么？过去两年，市场习惯将中国自动化需求与消费电子周期和新能源扩产周期挂钩。当这两个领域进入季节性调整时，市场容易得出“需求走弱”的结论。但本轮的不同之处在于，机床工具行业的持续走强，暗示着更底层的资本开支逻辑正在切换——从产能扩张驱动的自动化采购，转向效率提升和技术升级驱动的设备更新。

报告在讨论机床工具订单走强时，给出了一个关键归因：“我们认为这在很大程度上归因于AI数据中心投资的增加。”这一判断虽然简短，但值得深入推敲。

机床工具是制造业的母机，其需求变化往往领先于整体自动化周期。如果AI数据中心投资确实在拉动机床工具订单，那么传导链条可能是：数据中心建设本身需要大量的服务器机架、冷却设备、电力基础设施，这些都需要高精度加工设备；更深一层，AI芯片的封装测试、服务器主板的生产，也在拉动电子制造设备的升级需求。

报告没有展开这一逻辑，但数据提供了佐证。2026年第一季度，中国工业机器人产量同比增长明显，而半导体集成电路产量同样处于高位。这两条线的交汇点，正是AI算力基础设施的国产化进程。如果这一判断成立，那么自动化需求的驱动力正在从“消费电子换机”和“新能源扩产”的二元结构，转向“AI基础设施投资”这一新增量。

5月数据中最被低估的信号，是“订单价值继续超过出货价值，且订单和出货均超过公司指引”。这意味着亚德客的产能利用率已经接近极限，而需求仍在涌入。

对于一家在中国市场占比约90\%的气动元件企业来说，订单积压通常有两种解释：一是下游客户在抢产能，二是企业的交货周期在拉长。无论哪种情况，都指向一个事实——自动化供应链的紧张程度可能比市场预期的更持久。

报告提到的“十五五规划强调智能制造和产业升级”，为这一判断提供了政策背景。中国制造业正在经历从“有没有”到“好不好”的转型，自动化设备的采购决策正在从成本导向转向效率导向。这种转变不会因为月度波动而逆转。

在所有下游行业中，能源和照明（太阳能相关）是唯一销售额同比下降5\%的领域。这一趋势与2024年以来的行业数据一致——太阳能产业链的产能过剩问题仍在消化中。

这提醒我们，自动化需求的复苏并非均匀分布。新能源领域的资本开支高峰已经过去，而AI和智能制造领域的资本开支才刚开始。对于投资者而言，区分“周期性底部反弹”和“结构性新增长”至关重要。太阳能相关领域的疲软不是自动化需求的拖累，而是行业结构变化的正常表现。

第一，机床工具订单的走强有多少是由AI数据中心直接拉动，

[... middle omitted ...]

户的扩产节奏会产生什么影响，报告没有展开。

这些问题恰恰是理解这份研报价值的关键。报告提供了高质量的数据和初步归因，但二阶效应和情景分析需要更深入的讨论。

对于关注中国自动化市场的读者，这份研报提供了一个更精细的观察框架。

传统的做法是盯着PMI和工业增加值。但这份报告展示了一个更有效的组合：将亚德客的日均销售额（代表气动元件需求）、Komtrax中国运行小时数（代表工程机械开工率）和日本机床订单（代表高端设备需求）三者放在一起交叉验证。当这三个指标同时走强时，说明自动化需求具有广泛的行业基础；如果只有其中一个走强，则需要警惕结构性分化。

目前来看，三个指标虽然月度波动各异，但方向一致——均处于历史较高水平。这支持了“需求韧性”的判断。

此外，报告提到的“十五五规划”和“智能制造”政策方向，为这一轮自动化周期提供了不同于以往的政策锚点。过去十年，自动化需求更多依赖出口拉动和房地产投资，而未来可能更多依赖产业升级和AI基础设施投资。这个框架转换，值得持续跟踪。

\subsection{[R026] 人民币中间价的定价权，正在发生一次静默转移}
{\small\color{refgray}归类：汇率与新兴市场债券：定价机制的结构性转移}

市场对人民币汇率的讨论，长期集中在两个焦点：一是中美利差，二是关税政策的短期冲击。但某头部外资投行最新发布的中间价模型报告，揭示了一个更值得关注的信号——决定人民币每日交易基准的定价机制，正在经历一次结构性的、但尚未被充分定价的调整。

这份报告的核心判断是：模型投影的美元兑人民币中间价，正在系统性偏离实际公布的中间价，且偏离的方向和幅度，已经不能用传统的“逆周期因子”来解释。这意味着，央行对汇率的管理框架，可能已经从“被动对冲外部波动”转向了“主动引导预期锚定”。

本文将基于该报告的核心数据与模型逻辑，拆解这一判断背后的三个层次：模型本身在说什么、模型错误意味着什么、以及这些变化对资产定价和交易策略的启示。

报告给出了两组关键数字。第一组是纯粹的模型投影：基于一篮子货币的隔夜变动，模型计算出2026年6月3日的中间价应为6.7705，较前一日官方收盘价低482个基点。第二组是加入逆周期因子后的投影：6.7957，较前一日实际中间价低230个基点。

这两组数字的差值，就是市场通常理解的“逆周期因子”的贡献。但真正值得关注的不是差值本身，而是差值的动态变化。报告没有直接给出历史序列，但从模型误差的历史图表（Fig. 2）中可以看到，模型误差在2025年1月曾达到-1800个基点，随后在2025年中收窄至0附近，2026年初又反弹至+600个基点。

这意味着什么？模型误差的波动，本质上反映了央行对模型结果的“干预强度”。当误差为正时，说明实际中间价高于模型投影（即央行在引导人民币贬值）；当误差为负时，说明实际中间价低于模型投影（即央行在引导人民币升值）。2026年以来的正误差，表明央行正在主动容忍甚至引导人民币的适度贬值，这与2025年上半年的“维稳”姿态形成了鲜明对比。

这一变化对市场参与者的含义是：过去依赖模型预测中间价并以此构建交易策略的投资者，需要重新校准框架。模型本身仍然是有效的基准，但“逆周期因子”的调整方向和幅度，已经从一个技术性参数，变成了一个政策意图的先行指标。

2. 隔夜贡献的国别结构，暴露了人民币定价权的外部锚正在重组

模型投影的隔夜变动，并非均匀地来自所有货币对。报告给出了前四大贡献来源：俄罗斯卢布贡献33.5个基点，欧元21个基点，泰铢9个基点，日元8个基点。

这个结构本身就是一个值得拆解的信号。人民币中间价的定价，传统上以美元指数和欧元/美元为核心。但在这份投影中，卢布的贡献已经超过了欧元。这并非偶然。随着中俄贸易结算中人民币占比的持续提升，以及俄罗斯央行将人民币纳入外汇储备，卢布/人民币之间的联动性正在增强。这种联动不是通过传统的外汇市场套利实现的，而是通过贸易结算和资本账户的“真实流动”传导的。

欧元仍然重要，但其贡献从历史高点回落，反映的是欧洲经济周期与中国经济周期的脱钩。泰铢和日元的贡献则更多体现区域产业链的联动——泰国和日本是中国在东盟和东亚的主要贸易伙伴，其货币变动会通过“竞争性贬值”或“成本传导”的渠道影响人民币的均衡估值。

这个结构的重组，对投资者意味着两件事。第一，预测人民币中间价，不能只看美元指数和欧元/美元，还需要跟踪卢布、泰铢等“非传统锚”的变动。第二，人民币汇率的波动来源正在多元化，这既增加了预测难度，也创造了新的套利机会——那些能够更早识别卢布或泰铢异动的交易者，可能获得相对于市场平均水平的优势。

3. 模型误差的“非线性”特征，暗示政策框架可能已发生根本性变化

报告没有直接讨论模型误差的统计特征，但从Fig. 2的图表中可以观察到，模型误差并非随机波动，而是呈现出明显的“聚类”特征。2025年1月出现-1800个基点的极端负误差后，误差在2025年中迅速收敛至0附近，随后在2026年初又反弹至+600个基点。

这种非线性的误差模式，

[... middle omitted ...]

出口部门提供了竞争性汇率环境。

对资产定价而言，这一框架变化的影响是深远的。如果中间价的波动率被系统性压低，那么基于人民币波动率进行交易的外汇期权策略，其定价模型可能需要重新校准。同时，人民币资产（如A股、在岸债券）的汇率风险溢价，也可能因波动率下降而出现结构性压缩。

4. 报告尚未完全回答的关键问题：逆周期因子的“可预测性”正在消失

这份报告最坦诚的地方，在于它没有试图给出一个“万能公式”来预测逆周期因子的调整。模型本身是一个基于一篮子货币的线性回归，而逆周期因子是一个“黑箱”——报告没有披露其具体的函数形式或参数。

更关键的是，从历史数据看，逆周期因子的调整似乎并不完全跟随可观测的经济变量（如中美利差、出口增速、外汇储备变动）。2025年1月的极端负误差（-1800个基点）发生在中美利差仍然较大的背景下，而2026年初的正误差则发生在中美关系出现缓和信号的背景下。这暗示，逆周期因子的调整可能还包含了“非经济”的考量，例如外交谈判的节奏、国内政治周期的需要。

\subsection{[R027] 五月新能源车销量超预期，但市场真正低估的不是需求，而是供给侧的再定价}
{\small\color{refgray}归类：中国新能源车：渗透率高位下的订单分化与出口爆发}

五月新能源车销量超预期，但市场真正低估的不是需求，而是供给侧的再定价

五月新能源车销量数据出炉，整体好于市场预期。某外资投行最新研报显示，纳入已披露数据的品牌后，五月新能源车行业批发量环比增长12\%，同比增长7\%，超出此前市场一致预期。比亚迪、蔚来、零跑、华为鸿蒙智行等品牌均录得显著环比增长。

这份数据本身并不令人意外。真正值得关注的，是数据背后正在发生的结构性变化：行业竞争格局正在从“谁能卖得多”转向“谁能在规模增长中守住定价权”。5月销量超预期的背后，隐藏着一个尚未被市场充分定价的供给侧再平衡过程。

报告的核心判断并非“需求回暖”，而是“供给侧的效率分化正在加速”——头部企业的规模优势开始转化为成本与定价的双重能力，而这一转化过程，将重新定义未来12个月中国新能源车行业的投资逻辑。

1. 比亚迪的出口爆发正在改变其估值逻辑，国内市场的同比持平反而成为安全垫

比亚迪5月新能源车总销量38.35万辆，同比持平，环比增长19\%。乍看之下，同比零增长似乎缺乏亮点。但深入拆解结构，真正的增量来自海外：5月海外销量达到16.06万辆，环比增长19\%，前五个月累计海外销量61.69万辆，同比增长65\%。

这一数据意味着什么？比亚迪的海外销量占比已从去年的约20\%提升至当前的约42\%。当一家公司的海外业务增速达到65\%，且月销量已经突破16万辆时，其估值框架就不能再仅用国内市场份额来定价。

更值得关注的是，国内市场的同比持平，恰恰说明比亚迪在国内的定价策略正在从“以量换价”转向“以价换量”。5月纯电动车型销量19.87万辆，同比下降3\%，而插电混动车型销量17.83万辆，同比增长3\%。这一结构变化表明，比亚迪正在国内主动控制纯电产品的促销力度，转而通过插混产品维持量价平衡。

对于投资者而言，需要重新思考的问题是：当一家公司海外业务增速达到65\%，且国内业务不再需要大幅降价来维持份额时，其盈利能力的改善空间可能被市场系统性低估。

2. 蔚来的环比反弹并非偶然，子品牌放量验证了多品牌战略的可行性

蔚来5月交付3.77万辆，同比增长62\%，环比增长28\%，是本次数据中环比增速最高的品牌之一。更值得拆解的是其内部结构：主品牌蔚来交付约2万辆，子品牌乐道交付1.2万辆，环比增长125\%；萤火虫交付5663辆，环比增长14\%。

乐道的环比爆发是本次数据中最值得关注的信号之一。一个子品牌在上市后迅速实现月销过万，且环比翻倍增长，说明蔚来的多品牌战略正在从概念走向落地。此前市场对蔚来的核心担忧之一，就是其高端定位能否向下兼容中低端市场。乐道的表现初步回答了这个问题。

但这里有一个尚未完全回答的关键问题：乐道的快速放量是否会对主品牌形成内部竞争？从5月数据看，蔚来主品牌并未因乐道放量而出现明显下滑，这至少说明当前阶段内部替代效应有限。然而，随着乐道继续上量，这一风险仍是需要持续跟踪的变量。

对行业而言，蔚来的案例表明，在新能源车市场，品牌矩阵的构建可能比单一爆款更具长期价值。但这也意味着，企业的管理复杂度将呈指数级上升。

3. 理想汽车的低迷暴露了增程式路线的阶段性瓶颈，而非终局问题

理想汽车5月交付3.34万辆，同比下降18\%，环比下降2\%，是本次主要品牌中唯一同比和环比均下滑的企业。前五个月累计交付16.26万辆，同比下降3\%。

理想的低迷并非需求问题，而是产品周期问题。报告指出，理想2026年的重心集中在L系列改款上。这意味着当前正处于老款清库与新款上市之间的过渡期，销量承压是预期内的。

但更深层的问题是：增程式路线的竞争壁垒正在被削弱。当比亚迪、零跑、华为等品牌纷纷推出增程产品时，理想在该技术路线上的先发优势正在被稀释。理想真正的护城河不在于增程技术本身，而在于其家庭用户定位和智能化体验的整合能力。

[... middle omitted ...]

0-20万元价格带获得市场份额。极氪则依托吉利集团的资源，在20-30万元价格带构建品牌溢价。

两条路线同时实现高增长，说明当前中国新能源车市场并非零和博弈。不同价格带、不同用户群体的需求正在同时扩张。这背后的宏观驱动力是新能源车渗透率的持续提升——当市场从早期用户转向主流用户时，需求结构会变得更加多元化，而非单一化。

这对投资决策的含义是：不应简单地将行业竞争理解为“谁赢谁输”的二元格局。在渗透率从40\%向60\%迈进的过程中，多个品牌可以在不同细分市场同时获得增长。

5. 报告尚未完全回答的关键问题：海外市场的利润转化率与贸易政策风险

本次报告提供了详尽的销量数据，但有一个关键假设尚未得到充分验证：海外销量的利润转化率。

以比亚迪为例，其海外销量增速达到65\%，但海外业务的毛利率、净利率水平如何？海外市场的定价策略是否可持续？这些问题在报告中并未完全展开。此外，欧盟对中国电动车的反补贴调查、美国的关税政策等贸易风险，也可能对海外业务的盈利能力产生重大影响。

\subsection{[R028] 市场真正低估的不是算力需求，而是AI基础设施的供给结构正在被重新定义}
{\small\color{refgray}归类：AI/算力与半导体：从周期品到战略资产的再定价}

市场真正低估的不是算力需求，而是AI基础设施的供给结构正在被重新定义

当NVIDIA CEO黄仁勋在GTC上宣布Vera Rubin已进入全面量产，以及Qualcomm在Computex上推出数据中心品牌Dragonfly时，市场大多将注意力集中在算力升级的节奏上。但某外资投行最新研报揭示了一个更重要的信号：这两场演讲共同宣告，AI行业已从“基础大语言模型”的静态问答范式，正式转入“代理式AI”与“物理AI”时代。这个转变的真正含义，不是GPU卖得更多，而是整个AI基础设施的供应链结构、芯片架构和商业模式都在经历一次根本性的重构。

报告的核心判断值得认真对待：市场低估的不是需求端，而是供给侧的再定价。当NVIDIA从一家GPU系统公司进化为全栈AI基础设施公司，当Qualcomm将边缘芯片的触角延伸至数据中心，当Vera CPU开始侵蚀x86在服务器处理器中的份额——这些变化叠加在一起，意味着中国AI产业链上许多公司的竞争位置和估值逻辑需要被重新审视。

以下是我们从这份研报中提炼出的五个关键洞察，以及一个报告尚未完全回答的关键问题。

1. Vera Rubin量产不是一次常规升级，而是NVIDIA从“卖芯片”到“卖工厂”的转折点

Vera Rubin进入全面量产，对于供应链而言，确认了2026年下半年的拉货周期。但研报真正值得注意的，是NVIDIA在GTC上对自己定位的重新定义：它不再是一家GPU公司，而是一家全栈AI基础设施公司。

这个表述不是公关话术。报告指出，NVIDIA暗示其供应链控制将从L10（板卡级）延伸到L11（机柜级）和L12（集群级）。Vera Rubin DSX AI Factory参考设计已经展示了一个完整的AI工厂基础设施堆栈，包括NVL机柜、Vera GPU机柜、Spectrum-6 SPX交换机、Groq 3 LPX机柜和Vera BlueField-4 STX存储机柜。

这意味着什么？NVIDIA正在将“AI工厂”作为一种标准化产品来销售。客户购买的将不再是一块块GPU，而是一个从计算、网络到存储都已预集成、预优化的完整系统。这将对供应链产生结构性影响：那些能够提供更高集成度服务（如机柜组装、液冷系统、整机交付）的中国厂商，其价值量将显著提升。研报中提到的FII、VGT、WUS等公司，以及联想作为生态系统合作伙伴，都可能受益于这一趋势。

但反过来说，那些仅提供单一零部件的供应商，如果无法进入NVIDIA的新供应链层级，其议价能力和增长空间可能被压缩。这是供给侧正在发生的结构性变化。

2. Vera CPU的真正威胁不是性能，而是它正在改写服务器处理器的竞争规则

Vera CPU的发布是这场GTC中最容易被忽视但可能影响最深远的信号。报告给出了几个关键数据：Vera CPU采用88核单片网状拓扑，避免模块化芯粒架构带来的跨芯片通信延迟；在复杂SQL数据库任务中比x86快3倍，在实时数据流处理中快6倍。

但真正重要的不是性能数字，而是Vera CPU的架构选择。它是首个采用LPDDR5内存的企业级CPU，提供1.2TB/s的内存带宽。这直接决定了它使用的内存接口芯片类型是SOCAMM/LPDDR5X，而非传统的RDIMM。

研报由此提出了一个值得警惕的推论：如果Vera CPU在服务器处理器中的份额持续扩大，那么对传统DIMM芯片的需求可能被侵蚀。报告中明确点出了澜起科技（Montage）面临的风险——其内存接口芯片业务可能受到长期需求冲击。

这个判断的逻辑链条是清晰的：NVIDIA正在用Vera CPU重新定义服务器处理器的内存架构，而这一变化将沿着供应链向上游传导。对于中国半导体公司而言，这既是挑战也是机会——那些能够提供SOCAMM/LPDDR5X相关芯

[... middle omitted ...]

的Blackwell架构RTX GPU和20核Grace CPU，可提供1 petaflop的本地化AI算力和128GB统一内存。

这个产品的战略意义在于：它让复杂的多模态自主AI代理能够在终端设备上原生运行，无需持续依赖云连接。研报指出，这将对中国的软件和ICT公司产生深远影响——它们正在向“代理提供商”和“基于Token的商业模式”迁移。

“软件应该被呈现为代理；Token是收入和利润的生成器。”黄仁勋在GTC上的这句话，实际上在重新定义软件的价值衡量方式。过去，软件按许可证或订阅收费；未来，按Token消耗量收费将成为主流。这对中国SaaS公司和AI应用开发者的商业模式设计、定价策略和成本结构都将产生根本性影响。

但这里有一个尚未被充分讨论的问题：Token化定价是否适合所有类型的AI应用？对于高频、低价值的推理任务，Token定价可能导致用户成本失控；对于低频、高价值的任务，又可能定价过低。中国公司需要找到适合自己细分市场的定价模型，而这需要时间和试错。

\subsection{[R029] 市场真正低估的不是需求，而是供给侧的再定价}
{\small\color{refgray}归类：地产与消费：二手房定价权转移与消费渠道重置}

2026年6月，某外资投行举办了其中国房地产专家电话系列的首场会议，邀请趋势动物创始人聂聪分享对市场的判断。这份研报的核心信号并非“市场何时见底”这类老问题，而是一个更结构性的判断：中国房地产市场的复苏路径正在从“是否复苏”演变为“不同资产类别之间的持续分化”。这种分化不是暂时的波动，而是供给结构、政策工具和需求偏好共同作用下的新均衡。

多数投资者仍在等待一个统一的拐点信号——全国房价普涨、成交量全面回暖。但这份研报呈现的证据指向另一个方向：市场正在被切割成三个截然不同的价格体系——政府回购支撑的微型单元、高租金收益率的非核心城市资产、以及仍在加速下跌的中高端改善型住房。这三条线的逻辑不同，驱动因素不同，对投资组合的含义也不同。

更值得关注的是，报告提出了两种截然不同的复苏情景。在基准情景下，复苏将是一个“从两端向中间渗透”的缓慢过程——优质新房和高收益率资产先行，中端住宅随后跟进，且节奏取决于各城市的库存水位。而在大规模政策刺激情景下，市场可能复制香港模式——风险偏好快速扩张，中价位住宅价格迅速反弹。这两种情景对资产定价的含义完全不同，而当前市场定价隐含的是哪一种，恰恰是最大的未解问题。

以下是我们从这份研报中提炼出的几个关键判断，以及它们对投资框架的含义。

1. 全国房价跌幅收窄的背后，是三个完全不同的价格体系在同时运行

如果只看全国层面的月度数据，结论可能是温和的：全国二手房价格指数月环比跌幅已经从年初的-0.5\%收窄到近期的-0.4\%。但这份收窄并不是因为“所有房子都跌得少了”，而是因为三个价格体系的力量在同时作用。

第一个体系是政府回购驱动的政策性市场。上海50平方米以下的微型单元自2025年11月触底以来已上涨4\%，50-70平方米的小户型自2026年1月以来上涨11\%。这不是市场自发的需求回暖，而是上海政府加速存量住房回购用于保障性租赁住房的直接结果。这部分价格回升是政策性的、有上限的——一旦回购规模放缓，价格动能可能随之消退。

第二个体系是高租金收益率的非核心城市资产。乌鲁木齐部分核心区域在连续六个月下跌后开始回升，背后是3.6\%的平均租金收益率；大理的旅游和养老地产在过去一年上涨了8\%，租金收益率约3\%。这些资产的定价逻辑更接近固定收益——在利率下行周期中，高收益率资产自然获得重定价。这部分回升是自发的，但规模有限。

第三个体系是仍在承压的主流住宅。北京和广州的房价全线承压，杭州和成都核心区的高端资产自2026年3月以来加速下跌。这些市场既没有政策性回购支撑，也没有高收益率保护，仍在经历供需失衡的调整。

这三个体系并存意味着，全国平均数据掩盖了比表面数字更大的结构性差异。投资者不能再用“全国房价”这个笼统概念来指导决策，而需要分别理解每个价格体系的驱动因素和可持续性。

2. 深圳和上海的分化路径正好相反，但都指向同一个逻辑：需求结构已经变了

上海和深圳是这轮分化中最具代表性的两个城市，但它们的反弹路径截然不同。

上海的反弹集中在微型和小户型，由政府回购主导。深圳的稳定则来自中高端市场——总价1000万至2000万元的住房已经触底并小幅回升，南山和福田等核心子市场录得明显上涨动力，背后是改善型需求的复苏。换句话说，上海是“政策托底低端”，深圳是“需求拉动中高端”。

这两个案例的共同点在于：需求已经不再是统一的“买房需求”，而是高度分层的。政府的回购需求、首次置业者的上车需求、改善型家庭的升级需求、以及投资客的收益率需求，正在被不同的资产类别分别满足。这意味着，未来的市场复苏不会是“雨露均沾”的普涨，而是每个细分市场各自找到自己的均衡价格。

这对开发商的启示是：产品定位比城市选择更重要。在上海，开发微型单元的企业可能比开发豪宅的企业更早受益；在深圳，专注于改善型产品的项

[... middle omitted ...]

中，租金收益率上升意味着房产的投资价值在改善。乌鲁木齐3.6\%的租金收益率、大理约3\%的收益率、成都部分单元超过3.5\%的收益率——这些数字在利率下行环境中变得有吸引力。

报告进一步指出，如果抵押贷款利率再下降50-100个基点，这些高租金收益率资产可能获得20-30\%的价值重估。这个判断的逻辑是清晰的：在无风险利率下降的背景下，任何提供稳定现金流且收益率高于无风险利率的资产，都会吸引资本流入。

但这个逻辑有一个尚未验证的关键假设：这些高租金收益率资产的租金能否维持稳定？如果经济下行导致租金下跌，收益率优势可能迅速消失。报告没有完全展开这个风险，但对于任何考虑配置这类资产的投资者来说，这是必须独立验证的问题。

4. 复苏的传导机制是“从两端向中间”，但中端市场的库存压力是最大变量

报告提出的基准复苏路径是：优质新房和高收益率资产先行，然后通过置换需求逐渐传导到中端住宅。这个传导机制在理论上是合理的——高端市场回暖会释放卖房者的购买力，从而带动中端市场交易活跃。

\subsection{[R030] KRAS G12C的真相：不是所有抑制剂都一样，组合疗法的门槛比想象中更高}
{\small\color{refgray}归类：医药与生物科技：ADC与双抗的交叉验证}

KRAS G12C的真相：不是所有抑制剂都一样，组合疗法的门槛比想象中更高

ASCO 2026进入第二天，讨论的焦点从前沿探索转向了一线治疗的格局重塑。这份某外资投行发布的第二天要点解析，透露出一个清晰但容易被忽略的信号：市场正在为KRAS G12C抑制剂寻找合理的临床定位，但真正的变量不是机制本身，而是分子层面的差异和组合策略的选择门槛。

报告的核心判断可以浓缩为一句话：KRAS G12C单药在一线非小细胞肺癌中的故事线比组合疗法更清晰，而组合疗法的有效性虽然真实存在，但需要更精细的患者分层、更主动的不良反应管理和更严格的剂量连续性。这不是一个“谁更好”的问题，而是一个“谁更适合哪种场景”的格局判断。

1. 单药在一线的故事线比组合更干净，但“干净”不等于“容易”

讨论者对于Divarasib和Elisrasib的数据给出了一个有趣的信号：对单药在一线的前景更有信心。这不是否定组合疗法的价值，而是承认了一个现实——当疗效与毒性需要权衡时，单药方案的确定性更高。

报告特别强调了一个容易被忽视的细节：讨论者认为，化疗联合方案相比免疫联合化疗的长期生存优势有限。这意味着，如果组合疗法的设计只是简单地“把化疗加进去”，而不是考虑与免疫检查点抑制剂的协同，那么这条路径的长期价值可能被高估。

对于专注于KRAS G12C的中国生物科技公司——如Abbisko、Genfleet、Jacobio——这个判断意味着，单药适应症的推进速度可能比市场预期的更具竞争力，前提是分子本身具备足够好的药代动力学特性和安全性窗口。但报告也隐含了一个未完全回答的问题：什么样的单药活性才算“足够好”？这个阈值尚未被明确界定。

2. “不是所有G12C抑制剂都生而平等”——分子层面的差异正在成为竞争的分水岭

这是报告中最重要的一个观点。讨论者明确指出了这一点，并且给出了具体的维度：ON/OFF读通（read-through）的差异、分子水平的药代动力学、靶点结合能力、中枢神经系统穿透性以及毒性谱。

市场上长期以来有一种倾向，把KRAS G12C抑制剂视为一个同质化的类别，认为只要机制相同，临床结果就会趋同。但这份报告用数据说明，事实恰恰相反。一些OFF抑制剂（即非共价结合的抑制剂）正在展示持久的单药活性，而ON抑制剂（共价结合）的读通情况仍然模棱两可。

这意味着，投资者在评估不同公司的KRAS G12C资产时，不能再仅仅看“有没有进入临床”或者“初步有效率是多少”，而需要深入到分子设计层面去理解差异。肝毒性和剂量调整的频率将成为区分优劣的关键指标——报告明确提到，肝毒性和剂量修改仍然是核心关注点。

对于中国公司而言，这是一个双刃剑：分子设计能力强的公司可能获得溢价，而那些仅仅依靠跟随策略的资产可能面临更大的竞争压力。

3. 非G12C和pan-RAS的探索让RAS图谱保持流动性，但化疗依赖和质量控制是硬约束

报告提到，非G12C和pan-RAS抑制剂的探索正在从多个方向推进，尤其是在胰腺导管腺癌（PDAC）中，组合策略仍然以化疗为基础。这是一个重要的现实检验：尽管RAS抑制剂的进展令人振奋，但在一线治疗中，化疗仍然是大多数组合方案的支柱。

报告还提出了两个需要关注的约束：MAPK通路的合作伙伴选择、以及谱系可塑性（lineage plasticity）带来的耐药风险。这些都不是短期可以解决的问题，而是需要更长时间的临床验证。

对于投资者来说，这意味着RAS领域的竞争格局仍然高度不确定。那些声称可以覆盖所有RAS突变的平台，在临床数据出来之前，仍然是一个需要验证的假设。报告没有完全展开的是：在非G12C领域，哪些合作伙伴（如SHP2、MEK）的组合策略最有可能首先突破？这个问题的答案可能决定下一轮RAS资产的价值重估。

[... middle omitted ...]

ormat depth）和载荷/连接子的逻辑正在成为核心竞争力。那些拥有多样化载荷平台和差异化连接子技术的公司，在一线竞争的格局中可能占据更有利的位置。但报告没有完全展开的是：交叉耐药性的具体机制是什么？如果患者在接受ADC治疗后进展，下一步的序贯治疗应该如何选择？这些问题将是未来临床实践和投资判断的核心。

5. 报告尚未完全回答的关键问题：剂量强度与生活质量之间的权衡如何量化？

报告在多个地方提到了剂量强度、生活质量（QoL）和疗效/毒性权衡。尤其是在KRAS G12C的讨论中，肝毒性和剂量调整被反复提及；在ADC的讨论中，皮肤和胃肠道毒性对生活质量的影响也被强调。

但报告没有给出一个清晰的框架来量化这种权衡。对于临床医生和患者来说，什么样的疗效改善可以接受一定程度的毒性增加？这个阈值在不同的癌种、不同的治疗线数、不同的患者群体中可能完全不同。对于投资者来说，这意味着那些能够提供更好的安全性窗口（从而维持更高的剂量强度）的分子，可能获得超出市场预期的商业价值。

\subsection{[R031] 肺癌一线治疗变局：真正值得关注的不是单一数据，而是两类新机制的“交叉验证”}
{\small\color{refgray}归类：医药与生物科技：ADC与双抗的交叉验证}

肺癌一线治疗变局：真正值得关注的不是单一数据，而是两类新机制的“交叉验证”

ASCO 2026 的 NSCLC 专场，数据密度极高。但如果我们只提炼一个判断，它应该是：一线肺癌的标准治疗正在从“PD-1 单药或联合化疗”的单一范式，转向“TROP2 ADC 联合免疫”与“PD-1xVEGF 双抗联合化疗”两条技术路线的并行竞争。 市场对 Keytruda 2028 年专利到期的担忧早已存在，但 ASCO 这次提供的关键信号在于：挑战者不再只是“另一个 PD-1”，而是两种作用机制截然不同的新药。它们各自的数据，正在交叉验证一个更大的命题——Keytruda 的统治地位并非不可撼动，但取代它的方式可能不止一种。

某外资投行在 ASCO 后的研报中，系统梳理了这两条路径的最新进展。本文基于该报告的核心逻辑，提炼出对产业决策者和投资者真正重要的几个层次。

1. TROP2 ADC 的数据增量有限，但“联合用药”的策略信号远比单药 PFS 更重要

MRK/Kelun 的 Sac-TMT（TROP2 ADC）在 ASCO 上公布了 OptiTROP-Lung05 的中国 III 期数据：Sac-TMT 联合 Keytruda 对比 Keytruda 单药治疗 PD-L1 阳性 NSCLC。中位 PFS 未达到 vs 5.7 个月，HR 显著。从数字上看，这符合市场预期——此前模型预测的联合组中位 PFS 约 15-17 个月，与本次公布的参数模型预测值（ 个月）高度一致。研报明确指出，这是“增量正面更新”，而非颠覆性突破。

真正值得关注的是两个隐含信息。第一，Sac-TMT 在 PD-L1 TPS 1-49\% 亚组中的 PFS HR 为 0.28，显著优于 TPS >= 50\% 亚组的 0.47。这意味着，TROP2 ADC 联合免疫可能在“中等表达”人群中展现出不成比例的优势——这恰恰是 Keytruda 单药效果不够理想的区间。如果后续全球 III 期数据能验证这一趋势，Sac-TMT 将有机会重新定义 PD-L1 分层下的治疗选择。

第二，研报特别提到 MRK 正在推进 Sac-TMT 针对“PD-L1 低表达人群”的策略，但细节尚未公布。这是一个关键的未解问题：如果 TROP2 ADC 能在 PD-L1 Personal reading notes and learning share only. Not investment advice.

ASCO 刚结束，肺癌一线治疗（1L NSCLC）的格局正在被重塑。目前这个市场被某款 PD-1 单抗+化疗统治，但它的专利 2028 年就到期了。这次大会，最值得关注的两条技术路线：TROP2 ADC 和 PD-(L)1xVEGF 双抗。

某款 TROP2 ADC 在中国 3 期临床（PD-L1 阳性人群）中，联合 PD-1 单抗 vs PD-1 单抗，中位 PFS 还没达到，但模型预测在 15-17 个月左右，和预期一致（约 15-20 个月）。客观缓解率（ORR）是 70.2\% vs 42.0\%。

⏳ 关键点：OS 数据还不成熟，但趋势有利（HR 0.55）。副作用主要是 ADC 带来的血液毒性和口腔炎，但整体可控，停药率低。

⚠️ 注意：这只是中国数据，更关键的全球 3 期数据要等今年晚些时候。另外，它在 PD-L1 低表达人群中的策略也还不清楚。

2️⃣ PD-(L)1xVEGF 双抗：首次证明 OS 获益，但挑战仍在

这是本次 ASCO 的亮点。某款 PD-1xVEGF 双抗（Ivo）联合化疗，在鳞状 NSCLC 中，首次证明 PFS 获益能转化为 OS 获益（中位 OS 27.9 vs 23.7 个月）。

💡 但这只是在中国人群中，且

[... middle omitted ...]

1 单抗标准疗法。 TROP2 ADC 数据扎实，但需要全球数据证明自己。 PD-(L)1xVEGF 双抗首次看到 OS 获益，但全球验证才是关键。 2026 年下半年到 2027 年，会是这些新药的关键数据读出期。

Front-line lung cancer takeaways from ASCO

ASCO conference featured a number of presentations of Ph2/3 data emerging in front-line NSCLC, which is currently dominated by MRK's Keytruda+/-chemo. We were most focused on MRK/Kelun's Sac-TMT (TROP2 ADC) and Akeso/SMMT's Ivo (PD1xVEGF bispec). See below and within for our takeaways.

\subsection{[R032] 市场低估的不是AI算力需求，而是IT服务商从“系统开发”到“代理设计”的职能跃迁}
{\small\color{refgray}归类：日本工程机械与IT服务：关税窗口与AI效率竞争}

市场低估的不是AI算力需求，而是IT服务商从“系统开发”到“代理设计”的职能跃迁

当全球投资者还在争论AI的资本开支何时见顶、企业AI应用的ROI是否成立时，一份来自某外资投行的日本IT服务行业研报，提供了一个更具结构性的观察视角。这份报告的核心判断并非关于算力采购的规模，而是关于AI价值链正在发生的、不易被察觉的职能转移：企业IT服务商（Sler）的业务核心，正在从传统的系统开发，转向企业级AI代理（Agent）的设计、管理与安全。这个转变，将重新定义该行业的定价权与增长天花板。

当前市场对AI的讨论，大多集中在模型能力、算力需求和基础设施投资上。但真正决定未来几年产业格局的，可能是“谁来把通用AI能力转化为企业可用的、成本可控的专用代理”。这份报告从NVIDIA Computex主题演讲中捕捉到的信号表明，这个问题的答案正在向IT服务商倾斜。而这恰恰是当前估值模型中尚未完全定价的部分。

1. AI ROI的困境正在催生一个新的中间层市场，而IT服务商是天然承接者

企业AI应用的ROI问题，已成为制约下一阶段部署的关键瓶颈。报告引用了Uber COO的公开表态——生产力提升是否足以证明AI的token成本是合理的，目前仍不清晰。与此同时，Amazon的策略正在从“最大化token使用量”转向“优化token使用量”。这不是孤立现象。随着模型能力提升，OpenAI API等服务的成本也在持续上升，部分企业已经开始质疑这些支出的合理性。

这个困境的本质是什么？不是AI没有价值，而是通用模型的“一刀切”模式在企业场景中产生了大量无效成本。企业需要的是针对特定任务、经过优化、成本可控的专用代理，而不是每次调用都让一个庞大的通用模型从头“思考”。

NVIDIA在Computex上发布的Nemotron 3 Ultra模型，正是针对这一痛点。它被定位为一个开源、低成本、高效率的企业代理开发基础模型。NVIDIA明确表示，企业需要自己设计专用的、低成本的代理。而谁最擅长做这件事？正是那些长期嵌入企业业务流程、理解行业逻辑、拥有客户数据和领域知识的IT服务商。

这意味着，一个全新的中间层市场正在形成：介于基础模型提供商和企业最终用户之间，由IT服务商承担的“企业代理设计、管理与安全”服务。这个市场的规模，可能远超当前市场对IT服务行业的线性增长预测。

2. 从“系统开发”到“代理设计”是一次商业模式的升维，而非简单的服务内容扩展

理解这个转变的深度，需要将其放在IT服务行业的历史演进中来看。过去二十年，日本IT服务商的核心业务是系统集成与软件开发——根据客户需求，编写代码，构建信息系统。这是一个按人天计费、项目制交付、边际成本递减空间有限的模式。

而企业AI代理的设计与管理，本质上是完全不同的业务形态。它不是一次性交付的项目，而是持续优化、迭代、运维的服务。企业需要决定“哪个任务用哪个模型”，需要根据任务难度和重要性进行模型路由，需要管理代理的安全性与合规性。这些工作，报告明确指出，将成为Sler未来的业务。

这意味着三件事：第一，收入模式从项目制向订阅制或持续服务制转变，提升收入的可预测性与客户生命周期价值。第二，定价权从“人天费率”向“价值定价”迁移——一个能够帮助企业将token成本降低30\%的代理设计方案，其价值远高于同等工时投入的系统开发。第三，竞争壁垒从“项目经验”转向“领域数据+模型调优能力”的复合壁垒，后者的可复制性更低。

当然，这里仍需验证。目前尚不清楚企业愿意为这类代理管理服务支付多高的溢价，也不清楚IT服务商能否在短期内建立起足够的人才储备和方法论。但趋势的方向是明确的。

3. 物理AI是下一个前沿，而拥有行业软件和数据的服务商将占据独特生态位

报告在讨论NVIDIA主题演讲

[... middle omitted ...]

的推进，这些资产的价值将被重新评估。报告对富士通表达了明确的关注，认为其将在物理AI时代获得增长预期。

这一判断对投资者的启示是：在评估IT服务商时，不能只看其AI相关的营收占比，更要看其行业软件的深度和数据资产的独特性。那些拥有不可替代的行业知识库和规则系统的公司，在物理AI时代可能获得超出市场预期的生态位。

4. 报告尚未完全回答的关键问题：模型层的竞争格局会如何影响服务商的议价能力？

NVIDIA发布Nemotron系列模型，意味着它不再仅仅是算力提供商，也开始进入模型层。报告提到NVIDIA计划后续发布Nemotron 4和5，并希望跟踪NVIDIA作为企业AI模型提供商的定位变化。

这引出一个关键但尚未被充分讨论的问题：如果模型层出现“赢家通吃”或“寡头垄断”的格局，服务商的议价空间有多大？如果基础模型本身足够便宜、足够好用，企业是否会绕过服务商直接使用？相反，如果模型层呈现碎片化、多模型并存的局面，服务商的“模型路由”和“代理管理”价值反而会上升。

\subsection{[R033] 日本IT服务业的真正分化点：谁能在AI驱动的效率红利中，把成本节约转化为利润增长}
{\small\color{refgray}归类：日本工程机械与IT服务：关税窗口与AI效率竞争}

日本IT服务业的真正分化点：谁能在AI驱动的效率红利中，把成本节约转化为利润增长

过去几个季度，市场对日本IT服务业的关注焦点一直集中在“需求是否可持续”。金融系统升级、企业数字化转型、政府IT现代化——这些驱动因素看起来都足够坚实。但最近一份覆盖三家非覆盖公司的研报，揭示了一个更微妙的信号：需求侧的故事已经不够用了。真正区分赢家与输家的，不再是订单增长的快慢，而是企业能否在AI驱动的开发效率提升中，把节省下来的人月成本重新转化为利润，而不是被客户的价格压力或内部的费用膨胀所吞噬。

这份研报通过对Simplex Holdings、NSD和DTS三家公司的实地调研，提供了一个极具价值的“对照实验”视角。三家公司同属金融IT服务领域，面对相似的市场环境，但它们的盈利前景却出现了显著分化。Simplex预计FY3/27营业利润增长19\%，连续第七年实现两位数增长；而NSD和DTS却只能给出2\%-3\%的利润增幅指引。这种分化不是需求端的问题——三家公司的订单环境都相当强劲——而是供给端的结构性差异在起作用。

对于关注日本IT服务板块的投资者而言，这意味着传统的“订单增速-收入增速-利润增速”线性分析框架已经不够用了。需要建立一个更精细的观察框架，来评估哪些公司真正具备将AI驱动的效率提升转化为可持续利润的能力。

1. 订单增长强劲但利润指引分化，说明行业已进入“效率竞争”阶段

三家公司最新的订单数据都相当健康。Simplex 4Q订单同比增长19\%，连续第三个季度实现两位数增长；NSD系统开发订单4Q同比增长17\%，较3Q的3\%明显加速；即使是表现最弱的DTS，虽然订单同比下滑4\%，但公司表示大型项目的回撤影响已基本消化，公共和电信领域订单正在复苏。

然而，三家的FY3/27利润指引却呈现出截然不同的面貌。Simplex预计营业利润增长19\%，利润率达24.6\%；NSD仅预计增长2\%，利润率维持在15.5\%左右；DTS预计增长3\%，利润率约11.2\%。

这种订单与利润之间的脱节，核心原因在于“费用投入”的差异。NSD预计FY3/27的SG\&A费用将大幅增长22\%（+23.5亿日元），主要用于AI解决方案开发、广告费用和办公空间扩张。DTS同样面临SG\&A费用增长13\%（+17亿日元），主要投入在AI和人力资源领域。而Simplex虽然也在增加研发投入（+19亿日元，主要用于生成式AI和Web3.0），但其收入增长足够强劲，足以吸收这些增量费用后仍保持利润率稳定。

这意味着什么？当行业需求从“爆发增长”转向“稳定增长”阶段时，企业面临的真正考验不再是“能不能拿到订单”，而是“能不能在投入新技术的同时保持盈利纪律”。那些能够将增量投资精准聚焦于高回报领域、避免费用无差别膨胀的公司，将在这一轮竞争中脱颖而出。

研报中最具洞察力的部分，是Simplex关于AI驱动开发的战略阐述。公司表示，AI驱动开发预计将从FY3/29开始对盈利产生贡献，如果成功，将帮助公司达到FY3/30营业利润目标区间的上限（250-300亿日元）。更重要的是，Simplex明确指出，由于其业务主要基于合同开发模式，AI带来的工时削减将直接转化为公司自身的利润增长。

这一逻辑的关键在于“剩余产能的再分配”。Simplex认为，即使客户在中期内施加价格压力，公司也可以通过将工时削减产生的剩余产能重新分配到其他项目，来提升整体盈利。研报作者对此表示认同，并认为这适用于整个行业。

这实际上揭示了一个重要的二阶效应：AI对IT服务业的冲击，不是简单的“替代人月”，而是“释放产能”。那些拥有强大项目管理和客户关系能力的公司，可以将释放出来的工程师资源快速重新部署到新的高价值项目上，从而在不变动人员规模的情况下实现收入增长。而那些缺乏这种再分配

[... middle omitted ...]

quasi-mandate contracts），即基于月度人月定价的合同模式。在这种模式下，公司计划通过增加上游业务来避免客户的价格压力。

而Simplex则主要采用项目制合同。这种模式的优势在于，当AI提升开发效率后，公司可以直接享受工时缩短带来的利润提升，而不需要与客户重新谈判定价。Simplex管理层明确表示，即使中期内客户价格压力增加，公司也可以通过剩余产能再分配来提升整体盈利。

这两种模式的差异，本质上决定了AI效率提升的“利润归属”问题。在准委托合同下，效率提升带来的好处更容易被客户通过定价谈判所侵蚀；而在项目制合同下，公司可以更直接地保留效率提升的收益。

当然，项目制合同也有其风险——如果项目估算不准确或技术能力不足，容易产生亏损项目。DTS就是一个典型案例。公司在4Q出现了约7亿日元的亏损项目，主要集中在新客户、新领域（如云服务）的项目上，原因正是缺乏相关经验。这提醒我们，项目制合同的红利并非自动获得，而是需要强大的项目管理和技术能力作为支撑。

\section{Disclaimer}
{\scriptsize\color{refgray}
This community is a paid learning and information-sharing community created for general educational purposes only. The membership fee is charged solely for access to community discussions, educational content, organization, and operational costs. It is not an advisory fee, management fee, performance fee, brokerage fee, or compensation for personalized financial, investment, legal, tax, or professional advice.\par
I participate and share information solely as an individual and for educational discussion among community members. I am not acting as a financial adviser, investment adviser, broker, dealer, portfolio manager, legal adviser, tax adviser, accountant, fiduciary, or any other licensed professional adviser. Nothing shared in this community constitutes, and should not be interpreted as, financial advice, investment advice, legal advice, tax advice, a recommendation, solicitation, offer, endorsement, instruction, or guarantee to buy, sell, hold, short, trade, or invest in any security, cryptocurrency, fund, derivative, strategy, or other financial product.\par
All content shared in this community, including messages, discussions, links, files, charts, examples, market commentary, watchlists, AI-generated or AI-polished summaries, and third-party materials, is general, impersonal, and educational in nature. It does not take into account any individual's financial situation, investment objectives, risk tolerance, investment horizon, liquidity needs, tax status, legal restrictions, or suitability requirements. No content should be relied upon as suitable for any specific person, account, portfolio, or financial decision.\par
Information shared in this community may be incomplete, inaccurate, outdated, speculative, AI-generated, AI-edited, or based on third-party internet sources that have not been independently verified. I make no representation or warranty as to the accuracy, completeness, timeliness, reliability, or suitability of any information shared.\par
Financial markets involve substantial risk, including the possible loss of principal. Past performance, historical data, hypothetical examples, simulated results, backtests, personal opinions, or educational examples do not guarantee future results. Each member is solely responsible for conducting their own research, exercising independent judgment, and making their own decisions. Before making any financial, investment, legal, or tax decision, members should consult qualified and licensed professionals.\par
Participation in this community, payment of any membership fee, or communication with me or other members does not create any adviser-client, fiduciary, professional, contractual, agency, partnership, employment, or other advisory relationship. By joining or participating in this community, you acknowledge and agree that you are solely responsible for your own decisions, actions, transactions, profits, losses, risks, and consequences, and that I shall not be liable for any direct, indirect, incidental, consequential, financial, legal, tax, or other loss or damage arising from or related to any content, discussion, information, or material shared in this community.\par
}
\end{document}
